% =====================================================================
%  ncwu_linear_algebra_book
%  来源：线性代数讲义（水墨青主题，15 文件 / 9 章 + 封面 + 附录）（YangTex 模板库 · 用户原稿全量保留版）
%  处理：仅删除编译产物 PDF / digest.lock，把 \input 递归展开为单文件；
%        正文、章节、例题、每张幻灯片、封面与参考文献一字未改。
%  字体：原稿的 macOS 字体声明全部保留，外加 \IfFontExistsTF 存在性判断；
%        在 Mac 上视觉与原 PDF 完全一致，在 Windows / Linux 上字体缺失时
%        自动跳过该声明改用默认字体，仍能正常编译。
%  编译：xelatex（目录 / 交叉引用需连续编译 2~3 遍，产品已自动处理）
% =====================================================================

% ===================================================================
%  线性代数讲义（2026—2027 学年新版）
%  编译：xelatex main.tex   （需编译两遍以生成目录；亦可 latexmk -xelatex）
%  主题：水墨青 —— 低调沉稳、单双栏混排、图文并茂
% ===================================================================
\documentclass[11pt,a4paper,openany]{ctexbook}

% ---------- 版式 ----------
\usepackage{geometry}
\geometry{
  a4paper,
  top=2.35cm, bottom=2.25cm,
  inner=2.75cm, outer=3.05cm,
  marginparsep=0.32cm, marginparwidth=1.55cm,
  headheight=22.5pt, headsep=0.58cm, footskip=1.05cm
}
\usepackage{setspace}
\setstretch{1.16}
\usepackage{multicol}
\usepackage{marginnote}
\usepackage{emptypage}

% ---------- 字体 ----------
\usepackage{fontspec}
\usepackage{unicode-math}
\IfFontExistsTF{Times New Roman}{\setmainfont{Times New Roman}}{}
\IfFontExistsTF{Hiragino Sans GB}{\setsansfont{Hiragino Sans GB}}{}
\IfFontExistsTF{Menlo}{\setmonofont{Menlo}}{}
\setmathfont{texgyretermes-math.otf}[Scale=MatchLowercase]
\IfFontExistsTF{Songti SC}{\setCJKmainfont{Songti SC}[BoldFont=STHeiti, AutoFakeSlant=0.16]}{}
\IfFontExistsTF{Hiragino Sans GB}{\setCJKsansfont{Hiragino Sans GB}}{}
\IfFontExistsTF{Menlo}{\setCJKmonofont{Menlo}}{}
\IfFontExistsTF{FandolKai}{\setCJKfamilyfont{kai}{FandolKai}}{\setCJKfamilyfont{kai}{FandolKai}}
\IfFontExistsTF{STHeiti}{\setCJKfamilyfont{hei}{STHeiti}}{\setCJKfamilyfont{hei}{FandolHei}}
\newcommand{\cjkkai}{\CJKfamily{kai}}
\newcommand{\cjkhei}{\CJKfamily{hei}}

% ---------- 颜色（水墨青 · 低调沉稳）----------
\usepackage{xcolor}
\definecolor{ink}{RGB}{40,52,58}        % 墨（正文主色）
\definecolor{inkl}{RGB}{104,116,122}    % 淡墨（次级说明）
\definecolor{qing}{RGB}{42,88,102}      % 天青（主强调色）
\definecolor{qingd}{RGB}{30,66,78}      % 深天青
\definecolor{dai}{RGB}{84,110,98}       % 黛绿
\definecolor{zhe}{RGB}{152,100,64}      % 赭石
\definecolor{lan}{RGB}{72,98,126}       % 靛蓝
\definecolor{line}{RGB}{212,218,218}    % 细线
\definecolor{soft}{RGB}{246,247,245}    % 极浅底
\definecolor{paper}{RGB}{253,252,249}   % 宣纸白

% ---------- 数学 ----------
\usepackage{amsmath}
\usepackage{mathtools}
\allowdisplaybreaks[3]
\DeclareMathOperator{\rank}{rank}
\DeclareMathOperator{\tr}{tr}
\DeclareMathOperator{\Span}{span}
\DeclareMathOperator{\Col}{Col}
\DeclareMathOperator{\Nul}{Nul}
\DeclareMathOperator{\Row}{Row}
\DeclareMathOperator{\diag}{diag}
\DeclareMathOperator{\proj}{proj}
\DeclareMathOperator{\sign}{sign}
\DeclareMathOperator{\cond}{cond}
\newcommand{\R}{\mathbb{R}}
\newcommand{\C}{\mathbb{C}}
\newcommand{\norm}[1]{\left\lVert #1 \right\rVert}
\newcommand{\abs}[1]{\left\lvert #1 \right\rvert}
\newcommand{\ip}[2]{\left\langle #1,\,#2 \right\rangle}
\newcommand{\Rnum}[1]{\mathbb{R}^{#1}}

% ---------- 图表 ----------
\usepackage{graphicx}
\usepackage{tikz}
\usetikzlibrary{arrows.meta,calc,backgrounds,fit,positioning,
                shapes.geometric,patterns,matrix,decorations.pathreplacing,
                decorations.pathmorphing,3d,perspective}
\usepackage{pgfplots}
\pgfplotsset{compat=1.18}
\usepgfplotslibrary{fillbetween}
% 知识图谱节点样式（front.tex 的“全书知识地图”使用）
\tikzset{
  nmap/.style={
    draw=#1!55, fill=white, rounded corners=1.6pt, line width=0.6pt,
    minimum width=3.9cm, minimum height=0.88cm, inner sep=3pt,
    text=ink, font=\small\sffamily,
  },
  nmap/.default=qing,
}
\usepackage{booktabs,array,tabularx,multirow}
\usepackage{adjustbox}
\usepackage{enumitem}
\usepackage{listings}

% ---------- 标题 / 页眉 / 目录 ----------
\usepackage{titlesec}
\usepackage{titletoc}
\usepackage{fancyhdr}

% ---------- 定理盒子 ----------
\usepackage[most]{tcolorbox}
\tcbuselibrary{skins,breakable,theorems}

% ===================================================================
%  页眉页脚
% ===================================================================
% openany 版式下取消 cleardoublepage 产生的空白页
\makeatletter
\renewcommand{\cleardoublepage}{\clearpage}
\makeatother
\pagestyle{fancy}
\fancyhf{}
\fancyhead[LE]{\small\sffamily\color{qing}\nouppercase{\leftmark}}
\fancyhead[RO]{\small\sffamily\color{qing}\nouppercase{\rightmark}}
\fancyfoot[C]{\small\sffamily\color{inkl}\thepage}
\renewcommand{\headrulewidth}{0.45pt}
\renewcommand{\headrule}{%
  \hbox to\headwidth{\color{qing!28}\leaders\hrule height \headrulewidth\hfill}}
\renewcommand{\footrulewidth}{0pt}
\fancypagestyle{plain}{%
  \fancyhf{}
  \fancyfoot[C]{\small\sffamily\color{inkl}\thepage}
  \renewcommand{\headrulewidth}{0pt}}

% ===================================================================
%  标题样式
% ===================================================================
\titleformat{\chapter}[display]
  {\sffamily\bfseries\color{ink}}
  {\filleft\fontsize{58}{58}\selectfont\color{qing!17}\thechapter}
  {4pt}
  {\fontsize{23}{28}\selectfont\raggedright}
  [\vspace{3pt}{\color{qing}\rule{\textwidth}{0.9pt}}]

\titleformat{\section}[hang]
  {\sffamily\bfseries\Large\color{qing}}{\thesection}{0.55em}{}
\titleformat{\subsection}[hang]
  {\sffamily\bfseries\large\color{ink}}{\thesubsection}{0.5em}{}
\titleformat{\subsubsection}[runin]
  {\sffamily\bfseries\normalsize\color{dai}}{\thesubsubsection}{0.5em}{}[~~]

\titlespacing*{\chapter}{0pt}{-6pt}{16pt}
\titlespacing*{\section}{0pt}{15pt plus 3pt}{6pt}
\titlespacing*{\subsection}{0pt}{11pt plus 2pt}{4pt}

\titlecontents{chapter}[0em]
  {\addvspace{7pt}\sffamily\bfseries\color{qing}}
  {\thecontentslabel\qquad}{}
  {\titlerule*[5pt]{$\cdot$}\contentspage}
\titlecontents{section}[1.9em]
  {\sffamily\color{ink}}
  {\thecontentslabel\qquad}{}
  {\titlerule*[5pt]{$\cdot$}\contentspage}

% ===================================================================
%  插图 / 表格 环境（避免浮动带来的排版不可控）
% ===================================================================
\newcounter{fignum}[chapter]
\newcounter{tabnum}[chapter]
\newenvironment{lafig}[2][]{%
  \par\addvspace{9pt}%
  \stepcounter{fignum}%
  \gdef\lafigcap{#2}%
  \begin{center}%
  \begin{adjustbox}{max width=\textwidth,center}%
  \begin{tikzpicture}[#1]%
}{%
  \end{tikzpicture}%
  \end{adjustbox}%
  \par\vspace{4pt}%
  {\small\sffamily\color{qing}\textbf{图~\thechapter.\arabic{fignum}}\quad \lafigcap\par}%
  \vspace{3pt}%
  \end{center}%
  \addvspace{8pt}%
}
\newenvironment{latab}[1]{%
  \par\addvspace{9pt}%
  \stepcounter{tabnum}%
  \gdef\latabcap{#1}%
  \begin{center}%
  \begin{minipage}{0.94\textwidth}%
}{%
  \end{minipage}%
  \par\vspace{4pt}%
  {\small\sffamily\color{qing}\textbf{表~\thechapter.\arabic{tabnum}}\quad \latabcap\par}%
  \vspace{3pt}%
  \end{center}%
  \addvspace{8pt}%
}
\newcommand{\marnote}[1]{%
  \marginnote{\footnotesize\sffamily\color{qing!80}\raggedright\setstretch{1.05}#1}}

% ===================================================================
%  定理 / 定义 / 例题 盒子
% ===================================================================
\tcbset{
  labox/.style={
    enhanced, breakable, boxrule=0.35pt, arc=1.5mm,
    left=8pt, right=8pt, top=6pt, bottom=6.5pt,
    fonttitle=\sffamily\bfseries\small,
    fontupper=\rmfamily,
    separator sign={~},
    description delimiters={~~（}{）},
    description font=\normalfont\rmfamily,
    before skip=9pt, after skip=9pt,
    colbacktitle=white, coltitle=qing,
  }
}

\newtcbtheorem[number within=chapter]{defn}{定义}{
  labox, colframe=qing!45, colback=qing!3,
  borderline west={2.4pt}{0pt}{qing!65},
}{dfn}
\newtcbtheorem[use counter from=defn]{thm}{定理}{
  labox, colframe=zhe!55, colback=zhe!3, coltitle=zhe,
  borderline west={2.4pt}{0pt}{zhe!75},
}{thm}
\newtcbtheorem[use counter from=defn]{prop}{命题}{
  labox, colframe=lan!50, colback=lan!3, coltitle=lan,
  borderline west={2.4pt}{0pt}{lan!70},
}{prp}
\newtcbtheorem[use counter from=defn]{lem}{引理}{
  labox, colframe=lan!45, colback=lan!2.5, coltitle=lan,
  borderline west={2.4pt}{0pt}{lan!60},
}{lem}
\newtcbtheorem[use counter from=defn]{coro}{推论}{
  labox, colframe=zhe!40, colback=zhe!2.5, coltitle=zhe,
  borderline west={2.4pt}{0pt}{zhe!55},
}{cor}
\newtcbtheorem[use counter from=defn]{exmp}{例}{
  labox, colframe=line, colback=soft,
  borderline west={2.4pt}{0pt}{dai!70}, coltitle=dai,
}{exm}
\newtcbtheorem[use counter from=defn]{remk}{注记}{
  labox, colframe=line, colback=soft, boxrule=0.3pt,
  borderline west={2.4pt}{0pt}{inkl!55}, coltitle=inkl,
}{rmk}
\newtcbtheorem[use counter from=defn]{alg}{算法}{
  labox, colframe=dai!45, colback=dai!3, coltitle=dai,
  borderline west={2.4pt}{0pt}{dai!70},
}{alg}

% 证明结束符（unicode-math 下 \square 不可用，改用 TikZ 空心方块）
\newcommand{\pfbox}{\tikz[baseline=-0.55ex]{\draw[qing, line width=0.65pt] (0,0) rectangle (0.60em,0.60em);}}

% 证明环境
\newenvironment{pf}{%
  \par\vspace{2pt}\noindent
  {\sffamily\bfseries\color{zhe}证明\quad}%
  \ignorespaces
}{%
  \hfill\pfbox\par\vspace{6pt}%
}
% 证明梗概
\newenvironment{pfsketch}{%
  \par\vspace{2pt}\noindent
  {\sffamily\bfseries\color{zhe}证明梗概\quad}%
  \ignorespaces
}{%
  \hfill\pfbox\par\vspace{6pt}%
}

% 章首导览（双栏）
\newtcolorbox{chapoutline}{
  enhanced, breakable, colback=qing!2.6, colframe=qing!22, boxrule=0.35pt,
  arc=1.8mm, left=9pt, right=9pt, top=7pt, bottom=7pt,
  coltitle=qing, fonttitle=\sffamily\bfseries\small,
  title=\sffamily 本章导览,
  borderline west={2.6pt}{0pt}{qing!55},
  before skip=4pt, after skip=10pt,
}
% 章末要点（双栏）
\newtcolorbox{chapkeys}{
  enhanced, breakable, colback=dai!2.6, colframe=dai!22, boxrule=0.35pt,
  arc=1.8mm, left=9pt, right=9pt, top=7pt, bottom=7pt,
  coltitle=dai, fonttitle=\sffamily\bfseries\small,
  title=\sffamily 本章要点,
  borderline west={2.6pt}{0pt}{dai!65},
  before skip=10pt, after skip=8pt,
}
% 历史 / 拓展边栏
\newtcolorbox{sidebar}[1]{
  enhanced, breakable, colback=zhe!2.4, colframe=zhe!20, boxrule=0.35pt,
  arc=1.6mm, left=9pt, right=9pt, top=6pt, bottom=6pt,
  coltitle=zhe, fonttitle=\sffamily\bfseries\small,
  title=\sffamily #1,
  borderline west={2.6pt}{0pt}{zhe!55},
  before skip=10pt, after skip=10pt,
}
% 小结提示条
\newtcolorbox{keypoint}{
  enhanced, breakable, colback=qing!4, colframe=qing!28, boxrule=0.35pt,
  arc=1.2mm, left=8pt, right=8pt, top=5pt, bottom=5pt,
  borderline west={2.4pt}{0pt}{qing},
  before skip=8pt, after skip=8pt,
}

% 习题小节标题（双栏内使用）
% 满宽细线（自动适配单/双栏宽度，且不受段落缩进影响）
\newcommand{\fullrule}[2][0.8pt]{%
  \par\vspace{2pt}\noindent{\color{#2}\rule{\linewidth}{#1}}\par\vspace{4pt}}
\newcommand{\exhead}{%
  \par\addvspace{12pt}%
  {\noindent\sffamily\bfseries\Large\color{qing} 习题\par}%
  \vspace{-4pt}\noindent{\color{qing!32}\rule{\linewidth}{0.7pt}}\par\vspace{6pt}}
\newcommand{\minihead}[1]{%
  \par\addvspace{8pt}%
  {\noindent\sffamily\bfseries\large\color{qing} #1\par}%
  \vspace{-4pt}\noindent{\color{qing!28}\rule{\linewidth}{0.6pt}}\par\vspace{5pt}}

% ---------- 代码环境 ----------
\lstset{
  basicstyle=\ttfamily\small,
  keywordstyle=\color{qing}\bfseries,
  commentstyle=\color{inkl}\itshape,
  stringstyle=\color{zhe},
  backgroundcolor=\color{soft},
  frame=single, framerule=0.3pt, rulecolor=\color{line},
  framesep=5pt, xleftmargin=6pt, xrightmargin=6pt,
  tabsize=2, showstringspaces=false, breaklines=true,
  captionpos=b, abovecaptionskip=6pt, belowcaptionskip=2pt,
}

% ---------- 列表 ----------
\setlist{itemsep=3pt, parsep=2pt, topsep=4pt, leftmargin=2.2em}
\setlist[enumerate]{labelsep=0.5em}

% ---------- 杂项 ----------
\setcounter{tocdepth}{1}
\renewcommand{\textfraction}{0.10}
\hyphenpenalty=100
\tolerance=900
\emergencystretch=2em
\raggedbottom

\usepackage{hyperref}
\hypersetup{
  colorlinks=true, linkcolor=qing, citecolor=qing, urlcolor=qingd,
  bookmarksnumbered=true, bookmarksopen=true, pdftitle={线性代数讲义},
  pdfauthor={杨建伟}, pdfcreator={XeLaTeX},
}

% ===================================================================
\begin{document}
\frontmatter
\pagenumbering{Roman}
% ===================== 封面 =====================
\begin{titlepage}
\thispagestyle{empty}
\enlargethispage{1.4cm}
\centering
\begin{tikzpicture}[remember picture,overlay]
  % 背景：居中的淡网格（线性世界的底纹）
  \begin{scope}[shift={([xshift=10.5cm,yshift=-6.0cm]current page.north west)}]
    \foreach \i in {0,...,8}{
      \draw[qing!6, line width=0.8pt] (\i*0.82-3.28,0) -- (\i*0.82-3.28,6.56);
      \draw[qing!6, line width=0.8pt] (-3.28,\i*0.82) -- (3.28,\i*0.82);
    }
  \end{scope}
  \draw[qing!16, line width=1.5pt]
    ([xshift=3.4cm,yshift=-26.1cm]current page.north west) --
    ([xshift=17.6cm,yshift=-26.1cm]current page.north west);
\end{tikzpicture}

\vspace*{1.25cm}

{\sffamily\bfseries\color{qing}\fontsize{12.5}{17}\selectfont
 华北水利水电大学\quad·\quad 数学与统计学院\quad·\quad 课程讲义系列\par}

\vspace{0.52cm}
{\color{qing}\rule{0.80\textwidth}{1.1pt}\par}
\vspace{0.80cm}

{\sffamily\bfseries\fontsize{50}{58}\selectfont\color{ink} 线 性 代 数\par}
\vspace{0.34cm}
{\sffamily\fontsize{19}{27}\selectfont\color{qing} 讲\qquad 义\par}
\vspace{0.16cm}
{\sffamily\fontsize{14.5}{21}\selectfont\color{qing!85} 2026—2027 学年新版\par}

\vspace{0.34cm}
{\rmfamily\fontsize{11.5}{16}\selectfont\color{inkl}
 Linear Algebra: A Course Companion\par}
\vspace{0.06cm}
{\rmfamily\fontsize{11.5}{16}\selectfont\color{inkl}
 for the 2026--2027 Academic Year\par}

\vspace{0.78cm}

\begin{tikzpicture}[>=Stealth, line cap=round, line join=round, scale=0.92]
  % ---------- 左：单位圆与正交基 ----------
  \begin{scope}
    \draw[line!80, very thin, step=0.5] (-1.75,-1.75) grid (1.75,1.75);
    \draw[qing!45, line width=0.7pt, ->] (-2.0,0) -- (2.05,0);
    \draw[qing!45, line width=0.7pt, ->] (0,-2.0) -- (0,2.05);
    \draw[qing, line width=1.35pt] (0,0) circle (1.35);
    \draw[qing!80, line width=1.5pt, ->] (0,0) -- (0.955,0.955)
      node[midway, above left, xshift=-1pt]{\scriptsize $\boldsymbol{v}_1$};
    \draw[dai!80, line width=1.5pt, ->] (0,0) -- (-0.955,0.955)
      node[midway, above right, xshift=2pt]{\scriptsize $\boldsymbol{v}_2$};
    \node[below=3pt, font=\small, color=inkl] at (0,-2.25) {单位圆：$\norm{x}=1$};
  \end{scope}
  % ---------- 变换箭头 ----------
  \draw[zhe, line width=1.6pt, ->] (2.45,0) -- (3.75,0)
    node[midway, above=4pt, font=\small]{$\boldsymbol{A}$};
  % ---------- 右：椭圆与奇异方向 ----------
  \begin{scope}[xshift=6.7cm]
    \draw[line!80, very thin, step=0.5] (-2.6,-1.75) grid (2.6,1.75);
    \draw[qing!45, line width=0.7pt, ->] (-2.85,0) -- (2.9,0);
    \draw[qing!45, line width=0.7pt, ->] (0,-2.0) -- (0,2.05);
    \draw[zhe, line width=1.35pt] (0,0) ellipse (2.0 and 0.85);
    \draw[zhe!85, line width=1.5pt, ->] (0,0) -- (2.0,0)
      node[midway, below, yshift=-2pt]{\scriptsize $\sigma_1 u_1$};
    \draw[zhe!70, line width=1.5pt, ->] (0,0) -- (0,0.85)
      node[midway, right, xshift=1pt]{\scriptsize $\sigma_2 u_2$};
    \node[below=3pt, font=\small, color=inkl] at (0,-2.25) {像：椭圆 $\norm{Ax}=1$ 的边界};
  \end{scope}
\end{tikzpicture}

\vspace{0.20cm}
{\centering\sffamily\footnotesize\color{inkl}
 封面主图：单位圆在线性变换 $A$ 下的像是椭圆，其半轴即奇异值 $\sigma_1,\sigma_2$\par}

\vspace{0.52cm}
{\color{qing}\rule{0.80\textwidth}{1.1pt}\par}
\vspace{0.62cm}

\sffamily\fontsize{10.5}{17}\selectfont\color{ink}
\begin{tabular}{r@{\hspace{0.85em}}l}
\textcolor{qing}{主\quad 编}   & 杨建伟（华北水利水电大学 数学与统计学院）\\[3pt]
\textcolor{qing}{适用对象}     & 数学类、统计类、信息与计算科学、工科各专业\\[3pt]
\textcolor{qing}{学时建议}     & 48—64 学时（含上机实践 8 学时）\\[3pt]
\textcolor{qing}{版\quad 次}   & 2025年 9 月第 1 版\quad·\quad 2026 年9 月第 2 次修订\\[3pt]
\textcolor{qing}{配套资源}     & 习题参考解答、Python / MATLAB 演示程序、课件源文件
\end{tabular}

\vfill
{\sffamily\fontsize{10.5}{15}\selectfont\color{inkl}
 二〇二六年九月\quad·\quad 谨以此讲义献给初入线性世界的同学们\par}
\vspace{0.15cm}
\end{titlepage}

% ===================== 版本说明页 =====================
\cleardoublepage
\thispagestyle{empty}
\vspace*{1.6cm}
{\sffamily\bfseries\Large\color{qing}版本说明\par}
\vspace{4pt}\noindent{\color{qing}\rule{\textwidth}{0.8pt}}\par
\vspace{0.6cm}

本讲义依据 2026—2027 学年《线性代数》课程教学大纲编写，参考同济大学《线性代数》、
David C.\ Lay《线性代数及其应用》、Gilbert Strang《线性代数导论》等国内外经典教材，
并结合编者十余年一线教学实践重新组织与改写。全书在以下方面作了特别处理：

\begin{multicols}{2}
\sffamily\small\setstretch{1.25}
\begin{itemize}[leftmargin=1.5em]
\item[$\diamond$] \textbf{几何先行。}每一概念先给出平面或空间的直观图像，再进入代数推演。
\item[$\diamond$] \textbf{矩阵即作用。}始终强调矩阵是“线性变换的一种描述”，而非单纯的数表。
\item[$\diamond$] \textbf{单双栏混排。}正文推导单栏排印，便于长公式与大幅插图；
      导览、要点、习题与附录采用双栏，便于检索与对照。
\item[$\diamond$] \textbf{计算与理论并重。}每章配置算法框、上机提示与工程背景材料。
\item[$\diamond$] \textbf{层次分明。}定义、定理、例题、注记分别用不同色系的边框标识，
      一眼即可分辨内容性质。
\item[$\diamond$] \textbf{面向应用。}第 6、7、9 章分别给出动力系统、最小二乘拟合与主成分分析的完整案例。
\end{itemize}
\end{multicols}

\vspace{0.3cm}
\begin{center}
\begin{tikzpicture}[>=Stealth]
  \foreach \x/\lab in {-6.0/\text{向量},-3.6/\text{矩阵},-1.2/\text{方程组},
                       1.2/\text{空间},3.6/\text{变换},6.0/\text{特征值}} {
    \node[draw=qing!45, fill=qing!3, rounded corners=1.6pt, minimum width=2.05cm,
          minimum height=0.72cm, font=\small\sffamily, text=ink] at (\x,0) {\lab};
  }
  \foreach \a/\b in {-6.0/-3.6,-3.6/-1.2,-1.2/1.2,1.2/3.6,3.6/6.0}{
    \draw[qing!55, line width=0.8pt, ->] (\a+1.05,0) -- (\b-1.05,0);
  }
\end{tikzpicture}
\end{center}
\vspace{0.2cm}
{\centering\sffamily\small\color{inkl}图 0.1\quad 全书知识主线：向量 $\to$ 矩阵 $\to$ 方程组 $\to$ 空间 $\to$ 变换 $\to$ 谱\par}

\vfill
\noindent\sffamily\footnotesize\color{inkl}
本书为校内讲义，仅供教学使用。排版引擎 XeLaTeX，插图由 TikZ / pgfplots 现场绘制。

% ===================== 全书知识地图 =====================
\cleardoublepage
\thispagestyle{empty}
\vspace*{0.9cm}
{\sffamily\bfseries\Large\color{qing}全书知识地图\par}
\vspace{4pt}\noindent{\color{qing}\rule{\textwidth}{0.8pt}}\par
\vspace{0.15cm}
{\sffamily\small\color{inkl}
 下图给出九章内容的依赖关系：自下而上为“基础 $\to$ 核心 $\to$ 桥梁 $\to$ 结构 $\to$ 应用”，
 实线箭头表示直接依赖，虚线表示交叉应用。\par}
\vspace{0.25cm}

\noindent\makebox[\textwidth][c]{%
\begin{tikzpicture}[>=Stealth, line cap=round, font=\small\sffamily]
  % ---------- 分层底纹 ----------
  \foreach \y/\lab/\cf in {0.55/基\quad 础/qing, 2.85/核\quad 心/lan,
                           5.15/桥\quad 梁/dai, 7.45/结\quad 构/zhe,
                           9.75/应\quad 用/qingd}{
    \fill[\cf!5] (0,\y-0.62) rectangle (14.5,\y+0.62);
    \draw[\cf!22, line width=0.4pt] (0,\y-0.62) -- (14.5,\y-0.62);
    \node[\cf, font=\footnotesize\sffamily] at (0.95,\y) {\lab};
  }
  \draw[line!70, line width=0.4pt] (1.85,0.05) -- (1.85,10.3);

  % ---------- 节点 ----------
  \node[nmap=qing] (n1) at (7.6,0.55) {第1章\quad 向量与矩阵运算};
  \node[nmap=lan]  (n2) at (3.3,2.85)  {第2章\quad 线性方程组};
  \node[nmap=lan]  (n3) at (7.6,2.85)  {第3章\quad 行列式};
  \node[nmap=lan]  (n4) at (11.9,2.85) {第4章\quad 向量空间};
  \node[nmap=dai]  (n5) at (7.6,5.15)  {第5章\quad 线性变换};
  \node[nmap=zhe]  (n6) at (4.4,7.45)  {第6章\quad 特征值与对角化};
  \node[nmap=zhe]  (n7) at (10.8,7.45) {第7章\quad 正交性与最小二乘};
  \node[nmap=qingd](n8) at (4.4,9.75)  {第8章\quad 二次型与极值};
  \node[nmap=qingd](n9) at (10.8,9.75) {第9章\quad 奇异值分解与应用};

  % ---------- 依赖箭头 ----------
  \foreach \a/\b in {n1/n2,n1/n3,n1/n4,n2/n5,n3/n5,n4/n5,n5/n6,n5/n7,n6/n8,n7/n9}{
    \draw[line!75, line width=0.85pt, ->] (\a) -- (\b);
  }
  \draw[line!60, line width=0.7pt, ->, dashed] (n6) to[bend left=22] (n9);
  \draw[line!60, line width=0.7pt, ->, dashed] (n7) to[bend right=22] (n8);
  \draw[zhe!60, line width=0.7pt, <->, dashed] (n6) -- (n7)
    node[midway, above=3.5pt, font=\scriptsize\sffamily, color=inkl]
      {对称矩阵 $\Rightarrow$ 正交对角化};
\end{tikzpicture}}
\vspace{0.12cm}
{\centering\sffamily\small\color{inkl}
 图 0.2\quad 全书九章的知识依赖图谱\par}

\vspace{0.35cm}

\begin{multicols}{2}
\sffamily\small\setstretch{1.26}
\minihead{贯穿全书的三条主线}

\textbf{代数主线}（可计算）：\\
矩阵运算 $\to$ 高斯消元 $\to$ 行列式展开 $\to$ 特征多项式 $\to$ QR / SVD 迭代。
这条线回答“\emph{怎么算}”。

\textbf{几何主线}（可想象）：\\
向量叠加 $\to$ 平面相交 $\to$ 有向体积 $\to$ 空间的拉伸旋转 $\to$ 单位圆变椭圆。
这条线回答“\emph{是什么}”。

\textbf{结构主线}（可推广）：\\
线性组合 $\to$ 子空间与基 $\to$ 四个基本子空间 $\to$ 谱分解 $\to$ 低秩逼近。
这条线回答“\emph{为什么}”。

\minihead{三者的配合方式}

每章都按“\textbf{几何直观 $\to$ 代数定义 $\to$ 算法步骤 $\to$ 应用案例}”
四步推进。读者若时间有限，可只沿几何主线浏览插图，
亦可只沿代数主线演练算法，二者在结构主线处汇合。

\minihead{几条常用捷径}

\begin{itemize}[leftmargin=1.4em]
\item 只求解方程组 $\Rightarrow$ 第 2 章 + 第 4 章前三节；
\item 只备考特征值与二次型 $\Rightarrow$ 第 3、6、8 章；
\item 只做数据分析与降维 $\Rightarrow$ 第 7、9 章；
\item 补习几何直觉 $\Rightarrow$ 第 1、5 章的全部插图。
\end{itemize}
\end{multicols}

\vfill
\noindent\sffamily\footnotesize\color{inkl}
图中虚线箭头提示：特征值理论（第 6 章）是奇异值分解（第 9 章）的对称化特例，
而正交化方法（第 7 章）则是二次型标准化（第 8 章）的主要工具。

% ===================== 前言 =====================
\cleardoublepage
\thispagestyle{empty}
\vspace*{1.2cm}
{\sffamily\bfseries\Large\color{qing}前\quad 言\par}
\vspace{4pt}\noindent{\color{qing}\rule{\textwidth}{0.8pt}}\par
\vspace{0.5cm}

\setstretch{1.22}
线性代数是一门“越往后越觉得当初没学透”的课程。初学时，它似乎只是行列式的展开、
矩阵的行变换与一堆莫名其妙的运算规则；等到学习微分方程、最优化、统计学、
机器学习乃至量子力学时，我们才恍然发觉：那些看似繁琐的记号，其实是在描述
\emph{空间本身的结构}——向量如何叠加，空间如何被拉伸、旋转与压缩，
一个变换在哪些方向上最“省力”。

本讲义试图换一种讲法：\textbf{把几何放在前面，把代数作为几何的语言}。
第一章从向量的加法与数乘出发，让读者先在平面上“看见”线性组合；
第二章用三张平面的相交情形解释方程组解的结构；第三章把行列式解释成有向体积；
第五章用一张网格的变形说明矩阵乘法；第六章在箭头的迭代中看到特征方向；
第九章则以“单位圆变成椭圆”这一幅图收束全书，把奇异值分解的几何含义说清楚。
\marnote{建议读者在每一章结束时，合上书默画一遍该章的主图。}

与此同时，讲义并未放松对计算能力的要求。高斯消元、格拉姆—施密特正交化、
特征多项式、二次型的配方法、QR 迭代等基本技能，都以算法框的形式给出完整步骤，
并配相应的 Python / MATLAB 命令（见附录 B）。每章末的习题分三级：
A 组为基础计算，B 组为推理证明，C 组为建模与应用，供不同需求的读者取舍。

讲义的排版也刻意向正式出版物看齐：正文单栏便于长推导，导览与要点双栏便于速查；
定义、定理、例题、注记分别用天青、赭石、黛绿与灰色的边框加以区分；
所有插图均以矢量方式现场绘制，缩放不失真。希望能为读者提供一份
既能\textbf{读得下去}、也能\textbf{查得起来}的教材。

\vspace{0.8cm}
\begin{flushright}
\cjkkai
教学相长，疏漏之处在所难免。\\
欢迎读者将意见与勘误反馈至编者，以便下一学年修订。\\[6pt]
\sffamily\color{qing}杨建伟\\[2pt]
\sffamily\footnotesize\color{inkl}2026 年 9 月\quad 于华北水利水电大学
\end{flushright}

% ===================== 使用说明（双栏） =====================
\cleardoublepage
\thispagestyle{empty}
\vspace*{1.0cm}
{\sffamily\bfseries\Large\color{qing}使用说明\par}
\vspace{4pt}\noindent{\color{qing}\rule{\textwidth}{0.8pt}}\par
\vspace{0.45cm}

\begin{multicols}{2}
\sffamily\small\setstretch{1.28}

\minihead{版面与栏式}

本书采用\textbf{单双栏混排}：正文推导、长公式与大幅插图使用单栏，
以保证可读性与公式完整；章首“本章导览”、章末“本章要点”与“习题”、
以及全部附录使用双栏，便于快速检索与左右对照。

\minihead{彩色边框的含义}

每类内容都有固定的色系，读者可凭颜色判断性质：

\begin{center}
\begin{tikzpicture}[font=\small\sffamily]
  \node[draw=qing!50, fill=qing!3, text width=4.4cm, text=ink,
        rounded corners=1.5pt, inner sep=4pt,
        label={[color=qing]left:\bfseries 定}] at (0,0) {\ 定义：概念的精确表述};
  \node[draw=zhe!55, fill=zhe!3, text width=4.4cm, text=ink,
        rounded corners=1.5pt, inner sep=4pt,
        label={[color=zhe]left:\bfseries 理}] at (0,-0.95) {\ 定理 / 命题 / 推论};
  \node[draw=line, fill=soft, text width=4.4cm, text=ink,
        rounded corners=1.5pt, inner sep=4pt,
        label={[color=dai]left:\bfseries 例}] at (0,-1.90) {\ 例题：完整演算过程};
  \node[draw=line, fill=soft, text width=4.4cm, text=ink,
        rounded corners=1.5pt, inner sep=4pt,
        label={[color=inkl]left:\bfseries 注}] at (0,-2.85) {\ 注记：补充与提醒};
\end{tikzpicture}
\end{center}

\minihead{习题分级}

\begin{itemize}[leftmargin=1.4em]
\item[\textbf{A}] 基础计算，覆盖全部必会技能；
\item[\textbf{B}] 推理与证明，训练逻辑表达；
\item[\textbf{C}] 建模与应用，鼓励上机实现。
\end{itemize}

\minihead{符号约定}

向量默认写作列向量 $\boldsymbol{x}$；矩阵用大写拉丁字母 $A,B,\dots$；
标量用小写字希腊或拉丁字母。若未特别说明，数域默认为实数域 $\R$。

\minihead{上机提示}

标有 \textcolor{dai}{\textbf{$\triangleright$ 上机}} 的段落给出可在
Python（NumPy）或 MATLAB 中直接运行的核心命令，建议边学边练。
浮点运算与理论结果存在舍入误差，判定时应使用容差。

\minihead{阅读建议}

第一遍可按教师进度逐章阅读；考前复习时，只看各章“本章要点”与
“本章导览”两处双栏内容，再回头做 A 组习题，效率最高。

\end{multicols}

\vfill
\noindent\sffamily\footnotesize\color{inkl}
题图说明：封面主图为“单位圆在线性变换下的像”，它同时是全书第 9 章
奇异值分解的几何出发点。


\cleardoublepage
{\sffamily\hypersetup{linkcolor=ink}\tableofcontents}

\mainmatter
\pagenumbering{arabic}
% ===================== 第 1 章 向量与矩阵 =====================
\chapter{向量与矩阵}
\label{ch:vecmat}

\begin{chapoutline}
\small
\begin{multicols}{2}
\textbf{学习目标}
\begin{itemize}[leftmargin=1.5em, itemsep=1.5pt]
\item 用几何直观理解向量的加法、数乘与线性组合；
\item 掌握矩阵的三重身份：数表、算子与映射；
\item 熟练完成矩阵乘法，并理解“行乘列”与“列的线性组合”两种视角；
\item 认识常见特殊矩阵及其运算性质。
\end{itemize}
\columnbreak
\textbf{关键词}
\par\smallskip
线性组合\quad 张成\quad 列向量\quad 行向量\\
矩阵乘法\quad 单位矩阵\quad 转置\quad 分块
\par\smallskip
\textbf{建议学时}
\par\smallskip
8 学时（含上机 1 学时）
\end{multicols}
\end{chapoutline}

\section{向量：既有大小，又有方向}

把 $n$ 个实数排成一列，就得到一个\textbf{列向量}：
\[
\boldsymbol{x}=\begin{bmatrix}x_1\\ x_2\\ \vdots\\ x_n\end{bmatrix}\in\Rnum{n},
\qquad
\boldsymbol{x}^{\mathsf T}=[\,x_1\ \ x_2\ \ \cdots\ \ x_n\,]\in\R{1\times n}
\]
称为行向量。在 $\Rnum{2}$ 与 $\Rnum{3}$ 中，向量可以被“看见”：
它是有长度、有方向的箭头，起点约定在原点。

\begin{defn}{向量的加法与数乘}{vec-op}
设 $\boldsymbol{u},\boldsymbol{v}\in\Rnum{n}$，$c\in\R$。定义
\[
\boldsymbol{u}+\boldsymbol{v}=(u_1+v_1,\dots,u_n+v_n)^{\mathsf T},
\qquad
c\,\boldsymbol{u}=(cu_1,\dots,cu_n)^{\mathsf T}.
\]
\end{defn}

\begin{exmp}{两种几何解释}{geo-add}
下图给出 $\Rnum{2}$ 中的向量加法：$\boldsymbol{u}+\boldsymbol{v}$ 是以
$\boldsymbol{u},\boldsymbol{v}$ 为邻边的平行四边形的对角线，
也是“先走 $\boldsymbol{u}$ 再走 $\boldsymbol{v}$”的三角形折线的终点。
数乘 $c\boldsymbol{v}$ 则是把 $\boldsymbol{v}$ 拉长（$|c|>1$）、
压短（$|c|<1$），并在 $c<0$ 时掉头反向。
\begin{lafig}[>=Stealth, line cap=round]{向量加法的平行四边形法则（阴影）与数乘的伸缩（虚线）}
  \draw[line!80, very thin, step=0.5] (-1.2,-1.0) grid (3.9,3.9);
  \draw[qing!55, line width=0.7pt, ->] (-1.4,0) -- (4.05,0);
  \draw[qing!55, line width=0.7pt, ->] (0,-1.2) -- (0,4.05);
  \fill[qing!9] (0,0) -- (2.0,1.0) -- (3.0,3.0) -- (1.0,2.0) -- cycle;
  \draw[dashed, ink!40] (2.0,1.0) -- (3.0,3.0) -- (1.0,2.0);
  \draw[dai, line width=1.6pt, ->] (0,0) -- (2.0,1.0);
  \node[dai] at (1.25,0.28) {\small $\boldsymbol{u}$};
  \draw[zhe, line width=1.6pt, ->] (0,0) -- (1.0,2.0);
  \node[zhe] at (0.52,1.35) {\small $\boldsymbol{v}$};
  \draw[qing, line width=1.9pt, ->] (0,0) -- (3.0,3.0);
  \node[qing] at (2.4,3.3) {\small $\boldsymbol{u}+\boldsymbol{v}$};
  \draw[zhe!60, dashed, line width=0.9pt, ->] (2.0,1.0) -- (2.5,2.0)
     node[midway, right, yshift=-1pt]{\scriptsize $\tfrac12\boldsymbol{v}$};
  \draw[dai!60, dashed, line width=0.9pt, ->] (0,0) -- (3.0,1.5)
     node[midway, below, xshift=6pt]{\scriptsize $1.5\,\boldsymbol{u}$};
\end{lafig}
\end{exmp}

\begin{remk}{加法是唯一的主角}{only-add}
线性代数的全部内容，几乎都建立在两种“保结构”的运算之上：
加法与数乘。凡是把加法与数乘“送进去、送出来”的映射，
我们都称之为线性的。这句话现在听起来像口号，
第五章会把它变成精确定义。
\end{remk}

\section{线性组合与张成空间}

\begin{defn}{线性组合与张成}{span-def}
对向量 $\boldsymbol{v}_1,\dots,\boldsymbol{v}_k\in\Rnum{n}$ 与标量
$c_1,\dots,c_k\in\R$，称
\[
c_1\boldsymbol{v}_1+c_2\boldsymbol{v}_2+\cdots+c_k\boldsymbol{v}_k
\]
为一个\textbf{线性组合}。全体线性组合之集记为
\[
\Span\{\boldsymbol{v}_1,\dots,\boldsymbol{v}_k\}
=\{\,c_1\boldsymbol{v}_1+\cdots+c_k\boldsymbol{v}_k : c_i\in\R\,\},
\]
称为由这组向量\textbf{张成}（span）的集合。
\end{defn}

\begin{keypoint}
张成的几何图像非常清晰：一个非零向量张成一条过原点的直线；
两个不成比例的向量张成一个过原点的平面；
三个“不太巧”的向量在 $\Rnum{3}$ 中张成整个空间。
反之，若 $\boldsymbol{v}_2=c\,\boldsymbol{v}_1$，
则 $\Span\{\boldsymbol{v}_1,\boldsymbol{v}_2\}$ 只是一条直线——
第二个向量没有带来任何新东西。
\end{keypoint}

\begin{exmp}{三个向量的三种命运}{span-ex}
\[
\Span\Big\{\begin{bmatrix}1\\0\end{bmatrix}\Big\}
=\text{$x$ 轴（一条直线）};\qquad
\Span\Big\{\begin{bmatrix}1\\0\end{bmatrix},
              \begin{bmatrix}0\\1\end{bmatrix}\Big\}=\Rnum{2};\qquad
\Span\Big\{\begin{bmatrix}1\\0\end{bmatrix},
              \begin{bmatrix}2\\0\end{bmatrix}\Big\}
=\text{$x$ 轴（仍是直线）}.
\]
第三组说明：张成集合的大小取决于向量的“方向是否重复”，
而不取决于向量的个数。这为第四章“线性相关性”埋下伏笔。
\end{exmp}

\begin{lafig}[>=Stealth, line cap=round]{线性组合的几何：左图两个不共线向量张成整个平面，
右图两个共线向量只能张成一条直线}
  % ---------- 左：张成平面 ----------
  \begin{scope}
    \fill[qing!6] (-2.4,-1.7) rectangle (2.4,1.7);
    \draw[line!80, very thin, step=0.5] (-2.4,-1.7) grid (2.4,1.7);
    \begin{scope}
      \clip (-2.4,-1.7) rectangle (2.4,1.7);
      \foreach \t in {-1,0,1}{
        \draw[qing!30, line width=0.6pt] (1.5*\t,-1.7) -- (1.5*\t,1.7);
      }
      \foreach \s in {-1,0,1}{
        \draw[dai!35, line width=0.6pt] (-2.4,1.0*\s) -- (2.4,1.0*\s);
      }
    \end{scope}
    \draw[qing!45, line width=0.7pt, ->] (-2.65,0) -- (2.65,0);
    \draw[qing!45, line width=0.7pt, ->] (0,-1.95) -- (0,1.95);
    \draw[qing, line width=1.5pt, ->] (0,0) -- (1.5,0)
      node[midway, below, yshift=-3pt, fill=paper, inner sep=1pt]{\small $\boldsymbol{v}_1$};
    \draw[dai, line width=1.5pt, ->] (0,0) -- (0,1.0)
      node[midway, left, xshift=-2pt]{\small $\boldsymbol{v}_2$};
    \draw[ink!70, line width=0.9pt, ->, dashed] (0,0) -- (1.5,1.0)
      node[right, xshift=1pt]{\small $c_1\boldsymbol{v}_1+c_2\boldsymbol{v}_2$};
    \node[font=\footnotesize\sffamily, color=inkl] at (0,-2.35)
      {两族平行网格线：$\Span\{\boldsymbol{v}_1,\boldsymbol{v}_2\}=\Rnum{2}$};
  \end{scope}
  % ---------- 右：共线退化 ----------
  \begin{scope}[xshift=9.6cm]
    \draw[line!80, very thin, step=0.5] (-2.4,-1.7) grid (2.4,1.7);
    \draw[qing!45, line width=0.7pt, ->] (-2.65,0) -- (2.65,0);
    \draw[qing!45, line width=0.7pt, ->] (0,-1.95) -- (0,1.95);
    \draw[qing!16, line width=4.2pt] (-2.35,0) -- (2.35,0);
    \draw[qing, line width=1.5pt, ->] (0,0) -- (1.5,0)
      node[midway, below, yshift=-2pt]{\small $\boldsymbol{v}_1$};
    \draw[dai, line width=1.5pt, ->] (0,0) -- (0.9,0)
      node[midway, above, yshift=2pt]{\small $\boldsymbol{v}_2=\tfrac35\boldsymbol{v}_1$};
    \node[font=\footnotesize\sffamily, color=inkl] at (0,-2.35)
      {$\boldsymbol{v}_2$ 与 $\boldsymbol{v}_1$ 共线：只能张成一条直线};
  \end{scope}
\end{lafig}

\section{矩阵：一张数表，一种作用}

\begin{defn}{矩阵}{mat-def}
由 $m\times n$ 个数排成的矩形数表称为 $m\times n$ 矩阵，记作
\[
A=\begin{bmatrix}
a_{11} & a_{12} & \cdots & a_{1n}\\
a_{21} & a_{22} & \cdots & a_{2n}\\
\vdots & \vdots & \ddots & \vdots\\
a_{m1} & a_{m2} & \cdots & a_{mn}
\end{bmatrix}\in\R{m\times n},
\]
其中 $a_{ij}$ 位于第 $i$ 行第 $j$ 列。当 $m=n$ 时称 $A$ 为\textbf{$n$ 阶方阵}。
\end{defn}

\begin{remk}{矩阵的三重身份}{three-roles}
同一张数表，在不同场合扮演三种角色：
\begin{enumerate}[label=(\arabic*)]
\item \textbf{数表}：储存数据，如 $m$ 名学生 $n$ 门课的成绩矩阵；
\item \textbf{算子}：对向量做线性变换，$A\boldsymbol{x}$ 把 $\boldsymbol{x}$ “搬”到别处；
\item \textbf{方程组的浓缩}：$A\boldsymbol{x}=\boldsymbol{b}$ 用一行符号写下了 $m$ 个方程。
\end{enumerate}
本讲义最强调第二种身份。
\end{remk}

\subsection{矩阵加法、数乘与转置}

设 $A=(a_{ij}),B=(b_{ij})\in\R{m\times n}$，$c\in\R$，则
\[
(A+B)_{ij}=a_{ij}+b_{ij},\qquad (cA)_{ij}=c\,a_{ij},\qquad
(A^{\mathsf T})_{ij}=a_{ji}.
\]
三者都按“逐元素”进行，只有\textbf{乘法}例外。这些运算满足
\[
A+B=B+A,\quad (A+B)+C=A+(B+C),\quad c(A+B)=cA+cB,
\]
\[
(A^{\mathsf T})^{\mathsf T}=A,\quad (A+B)^{\mathsf T}=A^{\mathsf T}+B^{\mathsf T},\quad
(cA)^{\mathsf T}=cA^{\mathsf T},\quad (AB)^{\mathsf T}=B^{\mathsf T}A^{\mathsf T}.
\]
最后一式尤其值得记住：\emph{转置把乘积的顺序颠倒过来}。

\section{矩阵乘法：两种视角，一个事实}

\begin{defn}{矩阵乘法}{matmul-def}
设 $A\in\R{m\times n}$，$B\in\R{n\times p}$，则 $C=AB\in\R{m\times p}$ 定义为
\[
c_{ij}=\sum_{k=1}^{n}a_{ik}b_{kj}
=\text{$A$ 的第 $i$ 行}\cdot\text{$B$ 的第 $j$ 列}.
\]
\end{defn}

\begin{keypoint}
乘法能进行的条件是\textbf{内维相等}：$(m\times n)(n\times p)$。
结果的形状由外侧两数决定。这个“内维相等”的口诀，比任何公式都好用。
\end{keypoint}

按定义算乘法是“行乘列”。但理解乘法，更靠下面两种视角。

\begin{prop}{列组合视角}{colview}
$AB$ 的第 $j$ 列，等于 $A$ 乘以 $B$ 的第 $j$ 列：
\[
AB=A\,[\,\boldsymbol{b}_1\ \boldsymbol{b}_2\ \cdots\ \boldsymbol{b}_p\,]
=[\,A\boldsymbol{b}_1\ \ A\boldsymbol{b}_2\ \ \cdots\ \ A\boldsymbol{b}_p\,].
\]
即\textbf{乘积的每一列都是 $A$ 各列的线性组合}，组合系数由 $\boldsymbol{b}_j$ 给出。
因此 $\Col(AB)\subseteq \Col(A)$。
\end{prop}

\begin{prop}{行组合视角}{rowview}
对称地，$AB$ 的第 $i$ 行是 $B$ 各行的线性组合：
\[
(AB)_{i\,*}=a_{i\,1}\,\boldsymbol{b}_{1\,*}+a_{i\,2}\,\boldsymbol{b}_{2\,*}+\cdots
+a_{i\,n}\,\boldsymbol{b}_{n\,*},
\]
因此 $\Row(AB)\subseteq \Row(B)$。
\end{prop}

\begin{lafig}[>=Stealth, line cap=round]{矩阵乘法的两种视角：左为“行乘列”的逐元素定义，右为“乘一个向量即取列的线性组合”}
  % 左：行乘列
  \begin{scope}
    \node[font=\small\sffamily, color=qing] at (1.9,2.5) {视角一：行乘列};
    \matrix (A1) [matrix of math nodes, left delimiter={[}, right delimiter={]},
                  nodes={minimum width=0.72cm, minimum height=0.46cm, font=\small},
                  row sep=-0.1pt, column sep=-0.1pt] at (0,0.62)
      { a_{11} & a_{12} & a_{13}\\ a_{21} & a_{22} & a_{23}\\ a_{31} & a_{32} & a_{33}\\ };
    \node[qing, font=\footnotesize] at (1.85,0.62) {$\times$};
    \matrix (B1) [matrix of math nodes, left delimiter={[}, right delimiter={]},
                  nodes={minimum width=0.72cm, minimum height=0.46cm, font=\small},
                  row sep=-0.1pt, column sep=-0.1pt] at (2.9,0.62)
      { b_{11} & b_{12}\\ b_{21} & b_{22}\\ b_{31} & b_{32}\\ };
    \node[qing, font=\footnotesize] at (4.35,0.62) {$=$};
    \matrix (C1) [matrix of math nodes, left delimiter={[}, right delimiter={)},
                  nodes={minimum width=0.72cm, minimum height=0.46cm, font=\small},
                  row sep=-0.1pt, column sep=-0.1pt] at (5.7,0.62)
      { c_{11} & c_{12}\\ c_{21} & c_{22}\\ c_{31} & c_{32}\\ };
    \draw[zhe, line width=1.1pt] ($(A1-2-1.north west)+(-0.06,0.02)$) rectangle
         ($(A1-2-3.south east)+(0.06,-0.02)$);
    \draw[zhe, line width=1.1pt] ($(B1-1-2.north west)+(-0.06,0.02)$) rectangle
         ($(B1-3-2.south east)+(0.06,-0.02)$);
    \draw[zhe, line width=1.1pt] ($(C1-2-2.north west)+(-0.06,0.02)$) rectangle
         ($(C1-2-2.south east)+(0.06,-0.02)$);
    \node[font=\footnotesize, color=inkl] at (2.9,-0.75)
      {第 $2$ 行 $\cdot$ 第 $2$ 列 $=c_{22}$};
  \end{scope}
  % 右：列组合
  \begin{scope}[xshift=9.7cm]
    \node[font=\small\sffamily, color=qing] at (2.3,2.5) {视角二：列的线性组合};
    \matrix (A2) [matrix of math nodes, left delimiter={[}, right delimiter={]},
                  nodes={minimum width=0.72cm, minimum height=0.46cm, font=\small},
                  row sep=-0.1pt, column sep=-0.1pt] at (0,0.62)
      { a_{11} & a_{12}\\ a_{21} & a_{22}\\ a_{31} & a_{32}\\ };
    \node[qing, font=\footnotesize] at (1.55,0.62) {$\times$};
    \node[draw=zhe, line width=1.1pt, inner sep=3pt, font=\small] (B2) at (2.8,0.62)
      {$\begin{matrix}x\\ y\end{matrix}$};
    \node[qing, font=\footnotesize] at (3.85,0.62) {$=$};
    \node[font=\small] at (5.05,0.62) {$x\,\boldsymbol{a}_1+y\,\boldsymbol{a}_2$};
    \node[font=\footnotesize, color=inkl] at (2.5,-0.75)
      {其中 $\boldsymbol{a}_i$ 为 $A$ 的第 $i$ 列};
  \end{scope}
\end{lafig}

\begin{prop}{矩阵乘法的基本法则}{matmul-law}
在维数允许的前提下：
\[
A(BC)=(AB)C,\qquad A(B+C)=AB+AC,\qquad (B+C)D=BD+CD,
\]
\[
AI=IA=A,\qquad A\,\boldsymbol{0}=\boldsymbol{0}.
\]
\end{prop}

\begin{remk}{交换律一般不成立}{no-commute}
矩阵乘法\textbf{不满足交换律}。取
$A=\begin{bmatrix}0&1\\1&0\end{bmatrix}$，$B=\begin{bmatrix}1&0\\0&-1\end{bmatrix}$，
则
\[
AB=\begin{bmatrix}0&-1\\1&0\end{bmatrix},\qquad
BA=\begin{bmatrix}0&1\\-1&0\end{bmatrix}.
\]
两者一个是顺时针旋转 $90^\circ$，一个是逆时针旋转 $90^\circ$。
几何上看：\emph{先翻转再旋转}与\emph{先旋转再翻转}
是两回事。第五章将把这一点讲透。
\end{remk}

由结合律可定义矩阵的幂：$A^k=AA\cdots A$（$k$ 个 $A$ 连乘），
并约定 $A^0=I$。幂是第六章讨论迭代与动力系统的基本语言。

\section{特殊矩阵与分块}

\begin{latab}{常见特殊矩阵一览（$n$ 阶方阵）}
\small\centering
\renewcommand{\arraystretch}{1.34}
\begin{tabular}{@{}lll@{}}
\toprule
\sffamily 名称 & \sffamily 定义特征 & \sffamily 关键性质\\
\midrule
单位矩阵 $I$ & $a_{ii}=1$，其余为 $0$ & 乘法的单位元，$AI=A$\\
对角矩阵 & 非对角元全为 $0$ & 乘向量即逐分量缩放\\
上三角矩阵 & $a_{ij}=0\ (i>j)$ & 行列式等于对角元之积\\
对称矩阵 & $A^{\mathsf T}=A$ & 必可正交对角化（第 6 章）\\
正交矩阵 & $Q^{\mathsf T}Q=I$ & 保长度、保角度，$\det Q=\pm 1$\\
零矩阵 $O$ & 全部元素为 $0$ & 注意 $AB=O\not\Rightarrow A=O$ 或 $B=O$\\
幂零矩阵 & 某个 $k$ 使 $A^k=O$ & 特征值必全为 $0$\\
\bottomrule
\end{tabular}
\end{latab}

\subsection{分块矩阵}

把矩阵切成小块再运算，是处理大型问题与理论推导的利器。设
\[
A=\begin{bmatrix}A_{11} & A_{12}\\ A_{21} & A_{22}\end{bmatrix},
\qquad
B=\begin{bmatrix}B_{11} & B_{12}\\ B_{21} & B_{22}\end{bmatrix},
\]
其中每个 $A_{ij}$ 与对应的 $B_{ij}$ 形状相配，则形式上有
\[
AB=\begin{bmatrix}
A_{11}B_{11}+A_{12}B_{21} & A_{11}B_{12}+A_{12}B_{22}\\
A_{21}B_{11}+A_{22}B_{21} & A_{21}B_{12}+A_{22}B_{22}
\end{bmatrix}.
\]
\textbf{规则与普通乘法完全一致，只是把“数”换成“块”。}唯一要注意的是
块的排列顺序不可交换，如 $B_{11}A_{11}\neq A_{11}B_{11}$。

\begin{exmp}{用分块看分块对角矩阵}{blkdiag}
若 $A=\diag(A_1,A_2)$（两个方阵块），则对任意 $k$ 有
$A^k=\diag(A_1^k,A_2^k)$。因此大矩阵的幂可以拆成小矩阵的幂分别计算——
这正是“解耦”思想的代数表达。
\end{exmp}

\begin{remk}{$\triangleright$ 上机：用 NumPy 验证乘法法则}{numpy-1}
{\small\ttfamily
import numpy as np\\
A = np.array([[0., 1.], [1., 0.]])\\
B = np.array([[1., 0.], [0., -1.]])\\
print(A @ B)\\
print(np.allclose(A @ B, B @ A))}
\par\smallskip
\noindent 运行后会发现前两行输出不同，最后一行打印 \texttt{False}——
这正是“矩阵乘法不满足交换律”的一个直接证据。
\end{remk}

\begin{sidebar}{人物侧记：凯莱与矩阵代数的诞生}
把矩阵当作独立的代数对象来研究，一般归功于英国数学家
\textbf{阿瑟·凯莱}（Arthur Cayley，1821—1895）。1858 年他在
\emph{A Memoir on the Theory of Matrices} 中给出了矩阵的加法、
乘法与逆的定义。有趣的是，凯莱是出于研究代数方程的需要才发明矩阵的，
他甚至没有意识到矩阵会在此后的物理学中无处不在。
今天我们习以为常的“矩阵即变换”，要等到二十世纪量子力学兴起之后才被普遍接受。
\end{sidebar}

\begin{chapkeys}
\small
\begin{multicols}{2}
\begin{itemize}[leftmargin=1.5em, itemsep=2pt]
\item 向量加法与数乘是线性代数仅有的两种基本运算。
\item 张成集合的大小取决于方向是否重复，而非向量个数。
\item 矩阵乘法要求内维相等：$(m\times n)(n\times p)=m\times p$。
\item $AB$ 的第 $j$ 列是 $A$ 各列的线性组合（系数来自 $B$ 的第 $j$ 列）。
\item 结合律成立，交换律一般不成立。
\item $(AB)^{\mathsf T}=B^{\mathsf T}A^{\mathsf T}$，顺序要颠倒。
\item 分块运算规则与普通乘法一致，但块的次序不可交换。
\end{itemize}
\end{multicols}
\end{chapkeys}

\exhead
\begin{multicols}{2}
\small
\textbf{A 组（基础计算）}
\begin{enumerate}[leftmargin=1.6em, itemsep=3pt]
\item 设 $\boldsymbol{u}=(1,-2,3)^{\mathsf T}$，$\boldsymbol{v}=(0,4,-1)^{\mathsf T}$。
      求 $2\boldsymbol{u}-3\boldsymbol{v}$，并指出 $\Span\{\boldsymbol{u},\boldsymbol{v}\}$ 的几何维数。
\item 计算下列乘积（能算则算，不能算请说明理由）：
      \[ \begin{bmatrix}1&2\\3&4\end{bmatrix}
         \begin{bmatrix}0&1\\1&0\end{bmatrix},\qquad
         \begin{bmatrix}1&2&3\end{bmatrix}
         \begin{bmatrix}4\\5\\6\end{bmatrix} \]
      \[ \begin{bmatrix}1\\2\\3\end{bmatrix}
         \begin{bmatrix}4&5&6\end{bmatrix}. \]
\item 设 $A=\begin{bmatrix}1&2\\0&1\end{bmatrix}$，求 $A^2,A^3$，并猜测 $A^n$ 的通式。
\item 求 $A=\begin{bmatrix}1&2&0\\3&-1&4\end{bmatrix}$ 的转置，并验证 $(A^{\mathsf T})^{\mathsf T}=A$。
\end{enumerate}
\textbf{B 组（推理证明）}
\begin{enumerate}[leftmargin=1.6em, itemsep=3pt, start=5]
\item 证明：若 $A^{\mathsf T}=-A$（反对称矩阵），则 $A$ 的对角元全为 $0$。
\item 证明：对任意方阵 $A$，$A+A^{\mathsf T}$ 对称，$A-A^{\mathsf T}$ 反对称，
      且 $A$ 可写成二者之和。此分解唯一吗？
\item 设 $A,B$ 为 $n$ 阶方阵且 $AB=BA$。证明 $(A+B)^2=A^2+2AB+B^2$，
      并举例说明去掉该条件后等式不再成立。
\end{enumerate}
\textbf{C 组（应用与建模）}
\begin{enumerate}[leftmargin=1.6em, itemsep=3pt, start=8]
\item 某工厂生产三种产品，需经两道工序。工序耗时矩阵 $A\in\R{2\times 3}$，
      产品计划产量向量 $\boldsymbol{q}\in\Rnum{3}$。说明 $A\boldsymbol{q}$ 的两个分量各表示什么。
\item 用 NumPy 随机生成 $100\times 100$ 矩阵，比较计算 $A^{10}$ 的两种方法
      （连乘与反复平方）的耗时，并解释差异。
\end{enumerate}
\end{multicols}

% ===================== 第 2 章 线性方程组与高斯消元 =====================
\chapter{线性方程组与高斯消元}
\label{ch:gauss}

\begin{chapoutline}
\small
\begin{multicols}{2}
\textbf{学习目标}
\begin{itemize}[leftmargin=1.5em, itemsep=1.5pt]
\item 从几何上理解方程组解的三种情形；
\item 熟练运用初等行变换把增广矩阵化为阶梯形；
\item 用主元与自由变量描述解集的完整结构；
\item 掌握可逆矩阵的判定与求法，认识 LU 分解。
\end{itemize}
\columnbreak
\textbf{关键词}
\par\smallskip
增广矩阵\quad 初等行变换\quad 主元\\
阶梯形\quad 自由变量\quad 可逆矩阵\quad LU 分解
\par\smallskip
\textbf{建议学时}
\par\smallskip
10 学时（含上机 1 学时）
\end{multicols}
\end{chapoutline}

\section{方程组的三种几何面目}

一个 $n$ 元线性方程 $a_1x_1+\cdots+a_nx_n=b$ 在 $\Rnum{2}$ 中是一条直线，
在 $\Rnum{3}$ 中是一张平面。所谓解方程组，就是求这些“平坦对象”的公共点。

\begin{exmp}{二元情形的三种可能}{2var-cases}
\[
\left\{
\begin{aligned}
x+2y&=4\\
3x- y&=1
\end{aligned}\right.
\qquad
\left\{
\begin{aligned}
x+2y&=4\\
2x+4y&=9
\end{aligned}\right.
\qquad
\left\{
\begin{aligned}
x+2y&=4\\
2x+4y&=8
\end{aligned}\right.
\]
三组方程分别对应：两直线相交（唯一解 $(6/7,\,11/7)$）、
两直线平行（无解）、两直线重合（无穷多解）。
\begin{lafig}[>=Stealth, line cap=round, scale=1.0]{二元线性方程组的三种几何情形：相交、平行与重合}
  % ============ 左：相交（x+2y=4 与 3x-y=1，交点 (6/7,11/7)）============
  \begin{scope}[xshift=0cm]
    \draw[line!80, very thin, step=0.5] (-0.9,-0.9) grid (2.9,2.9);
    \draw[qing!55, line width=0.7pt, ->] (-1.05,0) -- (3.05,0);
    \draw[qing!55, line width=0.7pt, ->] (0,-1.05) -- (0,3.05);
    \draw[zhe, line width=1.5pt] (-0.9,2.45) -- (2.9,0.55)
      node[pos=0.14, above, yshift=1pt]{\scriptsize $x{+}2y=4$};
    \draw[dai, line width=1.5pt] (0.03,-0.9) -- (1.3,2.9)
      node[pos=0.85, right, xshift=1pt]{\scriptsize $3x{-}y=1$};
    \fill[qing] (0.857,1.571) circle (2.3pt);
    \node[qing] at (1.72,1.05) {\small 唯一解};
    \node[font=\footnotesize\sffamily, color=inkl] at (1.0,-1.5) {相交};
  \end{scope}
  % ============ 中：平行（x+2y=4 与 2x+4y=9，同斜率 -1/2）============
  \begin{scope}[xshift=5.0cm]
    \draw[line!80, very thin, step=0.5] (-0.9,-0.9) grid (2.9,2.9);
    \draw[qing!55, line width=0.7pt, ->] (-1.05,0) -- (3.05,0);
    \draw[qing!55, line width=0.7pt, ->] (0,-1.05) -- (0,3.05);
    \draw[zhe, line width=1.5pt] (-0.9,2.45) -- (2.9,0.55)
      node[pos=0.14, above, yshift=1pt]{\scriptsize $x{+}2y=4$};
    \draw[dai, line width=1.5pt] (-0.9,2.7) -- (2.9,0.8)
      node[pos=0.88, below, yshift=-2pt]{\scriptsize $2x{+}4y=9$};
    \node[font=\footnotesize\sffamily, color=inkl] at (1.0,-1.5) {平行};
    \node[zhe] at (2.05,1.9) {\small 无公共点};
  \end{scope}
  % ============ 右：重合（x+2y=4 与 2x+4y=8，同一条直线 y=2-x/2）============
  \begin{scope}[xshift=10.0cm]
    \draw[line!80, very thin, step=0.5] (-0.9,-0.9) grid (2.9,2.9);
    \draw[qing!55, line width=0.7pt, ->] (-1.05,0) -- (3.05,0);
    \draw[qing!55, line width=0.7pt, ->] (0,-1.05) -- (0,3.05);
    \draw[line width=3.2pt, qing!22] (-0.9,2.45) -- (2.9,0.55);
    \draw[dai, line width=1.2pt] (-0.9,2.45) -- (2.9,0.55);
    \draw[zhe, dashed, line width=1.0pt] (-0.9,2.45) -- (2.9,0.55);
    \node[font=\footnotesize\sffamily, color=inkl] at (1.0,-1.5) {重合};
    \node[qing] at (0.95,2.55) {\small 两方程同解：同一条直线};
  \end{scope}
\end{lafig}
\end{exmp}

在更高维，情形依然只有这三种，但“为什么只有三种”要到第二章末尾
才能说清——它本质上是秩的一个结论。

\begin{lafig}[>=Stealth, line cap=round, x={(0.92cm,0.44cm)}, y={(-0.50cm,0.46cm)}, z={(0cm,1cm)}]
{三元方程组的几何：三个平面的三种位置关系——交于一点、交于一条直线、没有公共点}
  % ================= 左：唯一解 =================
  \begin{scope}
    \fill[qing!9]  (0,0,0) -- (2.2,0,0) -- (2.2,0,1.35) -- (0,0,1.35) -- cycle;
    \fill[dai!9]   (0,0,0) -- (0,2.2,0) -- (0,2.2,1.35) -- (0,0,1.35) -- cycle;
    \fill[zhe!8]   (0,0,0) -- (2.2,0,0) -- (2.2,2.2,0) -- (0,2.2,0) -- cycle;
    \draw[qing!60, line width=0.7pt] (0,0,0) -- (2.2,0,0) -- (2.2,0,1.35) -- (0,0,1.35) -- cycle;
    \draw[dai!60, line width=0.7pt]  (0,0,0) -- (0,2.2,0) -- (0,2.2,1.35) -- (0,0,1.35) -- cycle;
    \draw[zhe!60, line width=0.7pt]  (0,0,0) -- (2.2,0,0) -- (2.2,2.2,0) -- (0,2.2,0) -- cycle;
    \fill[qing] (0,0,0) circle (2.2pt);
    \node[qing, font=\footnotesize] at (0.42,0.42,0.30) {$P$};
    \node[font=\footnotesize\sffamily, color=inkl] at (1.05,1.05,-1.75) {交于一点：唯一解};
  \end{scope}
  % ================= 中：无穷多解 =================
  \begin{scope}[xshift=6.0cm]
    \fill[qing!9]  (0,0,0) -- (2.2,0,0) -- (2.2,0,1.35) -- (0,0,1.35) -- cycle;
    \fill[zhe!8]   (0,0,0) -- (2.2,0,0) -- (2.2,2.2,0) -- (0,2.2,0) -- cycle;
    \fill[dai!12]  (0,-0.35,0) -- (0,2.4,0) -- (1.45,2.4,1.45) -- (1.45,-0.35,1.45) -- cycle;
    \draw[qing!60, line width=0.7pt] (0,0,0) -- (2.2,0,0) -- (2.2,0,1.35) -- (0,0,1.35) -- cycle;
    \draw[zhe!60, line width=0.7pt]  (0,0,0) -- (2.2,0,0) -- (2.2,2.2,0) -- (0,2.2,0) -- cycle;
    \draw[dai!65, line width=0.7pt]  (0,-0.35,0) -- (0,2.4,0) -- (1.45,2.4,1.45)
      -- (1.45,-0.35,1.45) -- cycle;
    \draw[ink, line width=1.5pt] (0,-0.35,0) -- (0,2.4,0);
    \node[ink, font=\footnotesize] at (0.42,1.15,0.16) {$L$};
    \node[font=\footnotesize\sffamily, color=inkl] at (1.05,1.05,-1.75) {交于一条直线：无穷多解};
  \end{scope}
  % ================= 右：无解 =================
  \begin{scope}[xshift=12.0cm]
    \fill[qing!9]  (0,0,0) -- (2.2,0,0) -- (2.2,0,1.75) -- (0,0,1.75) -- cycle;
    \fill[zhe!8]   (0,0,0) -- (2.2,0,0) -- (2.2,2.2,0) -- (0,2.2,0) -- cycle;
    \fill[dai!12]  (0,0,0.95) -- (2.2,0,0.95) -- (2.2,2.2,0.95) -- (0,2.2,0.95) -- cycle;
    \draw[qing!60, line width=0.7pt] (0,0,0) -- (2.2,0,0) -- (2.2,0,1.75) -- (0,0,1.75) -- cycle;
    \draw[zhe!60, line width=0.7pt]  (0,0,0) -- (2.2,0,0) -- (2.2,2.2,0) -- (0,2.2,0) -- cycle;
    \draw[dai!65, line width=0.7pt]  (0,0,0.95) -- (2.2,0,0.95) -- (2.2,2.2,0.95)
      -- (0,2.2,0.95) -- cycle;
    \draw[ink!70, dashed, line width=0.8pt] (0,1.05,0) -- (0,1.05,0.95);
    \node[inkl, font=\scriptsize] at (0.5,1.15,0.5) {$d$};
    \node[font=\footnotesize\sffamily, color=inkl] at (1.05,1.05,-1.75) {相互错开：无解};
  \end{scope}
\end{lafig}

\begin{remk}{如何读这张图}{read-3planes}
左图中三个平面像房间的一个墙角，公共点 $P$ 就是解；
中图中三个平面像一本翻开的书再加一张斜插的纸，公共部分是一条直线 $L$，
通解即 $L$ 的参数表示；右图中两张水平平面相互平行，无论第三个平面怎么放，
至多与它们各交出一条平行线，三个平面没有公共点。
“点—线—空”三种交叠方式，对应后面将要定义的三种秩情形
$\rank(A)=3$、$\rank(A)=2$ 与方程组不相容。
\end{remk}

\section{增广矩阵与初等行变换}

把方程组的系数与右端项一起抄进一张数表，就得到\textbf{增广矩阵}：
\[
\left\{
\begin{aligned}
x+\ 2y+\ 2z&=6\\
2x+\ 5y+\ 7z&=15\\
-2x-4y-4z&=-12
\end{aligned}\right.
\qquad\longleftrightarrow\qquad
\left[\begin{array}{ccc|c}
1 & 2 & 2 & 6\\
2 & 5 & 7 & 15\\
-2 & -4 & -4 & -12
\end{array}\right].
\]

\begin{defn}{初等行变换}{row-op}
对增广矩阵允许以下三种操作，它们不改变方程组的解集：
\begin{enumerate}[label=(\arabic*)]
\item \textbf{交换}：交换两行的位置，记作 $R_i\leftrightarrow R_j$；
\item \textbf{倍乘}：用非零常数 $c$ 乘某一行，记作 $cR_i$；
\item \textbf{倍加}：把某一行的 $c$ 倍加到另一行，记作 $R_j\leftarrow R_j+cR_i$。
\end{enumerate}
\end{defn}

\begin{exmp}{完整的高斯消元}{gauss-full}
继续上面的例子：
\[
\left[\begin{array}{ccc|c}1&2&2&6\\2&5&7&15\\-2&-4&-4&-12\end{array}\right]
\xrightarrow[\;R_3\leftarrow R_3+2R_1\;]{\;R_2\leftarrow R_2-2R_1\;}
\left[\begin{array}{ccc|c}1&2&2&6\\0&1&3&3\\0&0&0&0\end{array}\right]
\]
消元提前结束了：第三行整行变成了零。这说明第三个方程是前两个的推论，
真正独立的方程只有两个，于是有一个自由变量。取 $z=t$，回代得
\[
y=3-3t,\qquad x=6-2(3-3t)-2t=6t.
\]
解集为
$\{\,(6t,\ 3-3t,\ t)^{\mathsf T}: t\in\R\,\}
=(0,3,0)^{\mathsf T}+t\,(6,-3,1)^{\mathsf T}$：
一条过 $(0,3,0)$、方向向量为 $(6,-3,1)^{\mathsf T}$ 的直线。
\end{exmp}

\begin{remk}{为什么行变换不改变解集}{why-rowop}
三种变换都对应“方程的等价变形”：交换是交换方程，倍乘是两边同乘非零常数，
倍加是用一个方程去消去另一个方程中的未知数。这些操作都可逆，
因此新旧方程组互为推论，解集相同。
\end{remk}

\section{阶梯形、主元与自由变量}

\begin{defn}{行阶梯形与简化行阶梯形}{rref}
称矩阵为\textbf{行阶梯形}（echelon form），若
全零行（若有）位于最下方，且每个非零行的首个非零元（\textbf{主元}）
位于上一行主元的右侧。
若进一步要求主元为 $1$ 且主元所在列的其余元素全为 $0$，
则称为\textbf{简化行阶梯形}（RREF）。
\end{defn}

\begin{keypoint}
\textbf{定理（RREF 唯一性）}：任何矩阵的简化行阶梯形是唯一的。
（行阶梯形不唯一，但主元的位置唯一。）
\end{keypoint}

\begin{exmp}{识别主元与自由变量}{pivot-read}
\[
\left[\begin{array}{ccccc|c}
\textcolor{zhe}{1} & 2 & 0 & 3 & 0 & 4\\
0 & 0 & \textcolor{zhe}{1} & 1 & 0 & 2\\
0 & 0 & 0 & 0 & \textcolor{zhe}{1} & 1\\
0 & 0 & 0 & 0 & 0 & 0
\end{array}\right]
\qquad
\text{主元列：第 }1,3,5\text{ 列；自由变量：}x_2,\ x_4.
\]
（赭色标出主元位置。）令 $x_2=s$，$x_4=t$，逐行读出
$x_1=4-2s-3t$，$x_3=2-t$，$x_5=1$，即得通解。
\end{exmp}

\begin{remk}{主元个数 = 秩}{pivots-rank}
主元的个数称为矩阵的\textbf{秩}（rank）。秩是全书最重要的数字之一：
\[
\#\{\text{主元}\}=\rank(A)=\rank(A\,|\,\boldsymbol{b})
\quad\Longleftrightarrow\quad
A\boldsymbol{x}=\boldsymbol{b}\ \text{有解}.
\]
而\textbf{自由变量的个数} $=n-\rank(A)$。
\end{remk}

\section{解的结构：特解加齐次解}

\begin{thm}{非齐次方程组解的完整结构}{sol-struct}
设 $A\in\R{m\times n}$，$\boldsymbol{b}\in\Rnum{m}$，且 $A\boldsymbol{x}=\boldsymbol{b}$ 有解。
取定一个解 $\boldsymbol{x}_p$（特解），则
\[
\{\boldsymbol{x}:A\boldsymbol{x}=\boldsymbol{b}\}
=\boldsymbol{x}_p+\Nul(A)
=\{\,\boldsymbol{x}_p+\boldsymbol{x}_h : A\boldsymbol{x}_h=\boldsymbol{0}\,\}.
\]
\end{thm}

\begin{pf}
若 $A\boldsymbol{x}=\boldsymbol{b}$，则 $A(\boldsymbol{x}-\boldsymbol{x}_p)=\boldsymbol{0}$，
故 $\boldsymbol{x}-\boldsymbol{x}_p\in\Nul(A)$。反之，若 $A\boldsymbol{x}_h=\boldsymbol{0}$，
则 $A(\boldsymbol{x}_p+\boldsymbol{x}_h)=\boldsymbol{b}$。
\end{pf}

\begin{lafig}[>=Stealth, line cap=round]{解集结构：特解 $\boldsymbol{x}_p$ 加上零空间 $\Nul(A)$ 的平移}
  \draw[line!80, very thin, step=0.5] (-2.4,-1.6) grid (2.4,1.6);
  \draw[qing!55, line width=0.7pt, ->] (-2.55,0) -- (2.6,0);
  \draw[qing!55, line width=0.7pt, ->] (0,-1.75) -- (0,1.8);
  % Nul(A)：过原点、斜率 0.5；解集：同斜率、过 x_p=(0.9,0.9)
  \draw[dai, line width=1.7pt] (-2.2,-1.1) -- (2.2,1.1)
     node[pos=0.09, below, yshift=-2pt]{\small $\Nul(A)$};
  \draw[zhe!70, dashed, line width=1.3pt] (-2.5,-0.8) -- (2.5,1.7);
  \draw[qing, line width=1.4pt] (-2.3,-0.7) -- (2.3,1.6)
     node[pos=0.82, below right, xshift=-2pt, yshift=1pt]{\small $\boldsymbol{x}_p+\Nul(A)$};
  \fill[zhe] (0.9,0.9) circle (2.2pt);
  \node[zhe, above right, xshift=-3pt] at (0.9,0.9) {\small $\boldsymbol{x}_p$};
  \draw[->, ink!60, line width=1.1pt] (0,0) -- (0.9,0.9);
  \fill[ink] (0,0) circle (1.9pt);
  \node[font=\footnotesize, color=inkl, below=3pt] at (0,-1.9)
    {解集是齐次解直线 $\Nul(A)$ 沿 $\boldsymbol{x}_p$ 平移得到的仿射直线};
\end{lafig}

\begin{remk}{一句话记住解的结构}{one-line-sol}
\textbf{非齐次的解集 = 一个特解 + 齐次的全部解。}
这条“平移”的观点在最小二乘（第 7 章）与 PCA（第 9 章）中会反复出现。
\end{remk}

\section{可逆矩阵}

\begin{defn}{逆矩阵}{inverse}
设 $A$ 为 $n$ 阶方阵。若存在 $B$ 使 $AB=BA=I$，则称 $A$ \textbf{可逆}，
$B$ 称为 $A$ 的逆，记作 $A^{-1}$。可逆矩阵必为方阵，且逆唯一。
\end{defn}

\begin{thm}{可逆的等价刻画}{invertible-thm}
对 $n$ 阶方阵 $A$，下列条件等价：
\begin{enumerate}[label=(\alph*)]
\item $A$ 可逆；
\item $A\boldsymbol{x}=\boldsymbol{0}$ 只有零解；
\item $\rank(A)=n$（每个列都是主元列）；
\item $A$ 的简化行阶梯形为 $I$；
\item $\det A\neq 0$（第三章）；
\item $A$ 的列向量张成 $\Rnum{n}$（第四章）。
\end{enumerate}
\end{thm}

\begin{prop}{逆的运算性质}{inv-law}
\[
(A^{-1})^{-1}=A,\qquad (AB)^{-1}=B^{-1}A^{-1},\qquad
(A^{\mathsf T})^{-1}=(A^{-1})^{\mathsf T},\qquad (cA)^{-1}=c^{-1}A^{-1}\ (c\neq0).
\]
\end{prop}

\begin{prop}{用行变换求逆}{inv-by-row}
$[A\,|\,I]\xrightarrow{\ \text{行变换}\ }[I\,|\,A^{-1}]$。
若左侧无法化为 $I$，则 $A$ 不可逆。
\end{prop}

\begin{exmp}{求逆并解方程组}{inv-ex}
$A=\begin{bmatrix}2&1\\1&1\end{bmatrix}$。作
$\left[\begin{array}{cc|cc}2&1&1&0\\1&1&0&1\end{array}\right]
\xrightarrow{R_1\leftrightarrow R_2}
\left[\begin{array}{cc|cc}1&1&0&1\\2&1&1&0\end{array}\right]
\xrightarrow{R_2\leftarrow R_2-2R_1}
\left[\begin{array}{cc|cc}1&1&0&1\\0&-1&1&-2\end{array}\right]
\xrightarrow[\,R_1\leftarrow R_1-R_2\,]{R_2\leftarrow -R_2}
\left[\begin{array}{cc|cc}1&0&1&-1\\0&1&-1&2\end{array}\right]$
故 $A^{-1}=\begin{bmatrix}1&-1\\-1&2\end{bmatrix}$。
于是 $A\boldsymbol{x}=\boldsymbol{b}$ 的解为 $\boldsymbol{x}=A^{-1}\boldsymbol{b}$。
\end{exmp}

\begin{remk}{数值警告：不要用逆来解方程}{no-inverse}
理论上 $\boldsymbol{x}=A^{-1}\boldsymbol{b}$ 很漂亮，数值上却应当避免：
计算 $A^{-1}$ 的代价是 $O(n^3)$ 且舍入误差更大。
工程与科学计算的标配是 LU 分解或 QR 分解。
\end{remk}

\section{LU 分解：消元的“记账法”}

高斯消元的每一步“倍加”都可以用一个初等矩阵实现。把所有消元步骤合起来，
就得到
\[
A=LU,\qquad
L=\begin{bmatrix}1&0&0\\ \ell_{21}&1&0\\ \ell_{31}&\ell_{32}&1\end{bmatrix},
\quad
U=\text{行阶梯形},
\]
其中 $L$ 是单位下三角矩阵（其元素恰为消元时用过的乘数），$U$ 是上三角矩阵。

\begin{prop}{LU 分解的用途}{lu-use}
解 $A\boldsymbol{x}=\boldsymbol{b}$ 等价于依次解两个三角方程组：
\[
L\boldsymbol{y}=\boldsymbol{b}\quad(\text{前代}),\qquad
U\boldsymbol{x}=\boldsymbol{y}\quad(\text{回代}).
\]
分解只需做一次（约 $\tfrac23 n^3$ 次乘加）；此后每换一个右端项 $\boldsymbol{b}$，
只需 $O(n^2)$ 次运算。这是科学计算中最经济的做法。
\end{prop}

\begin{remk}{$\triangleright$ 上机：SciPy 的 LU 分解}{lu-code}
{\small\ttfamily
from scipy.linalg import lu\_factor, lu\_solve\\
lu, piv = lu\_factor(A)\\
x = lu\_solve((lu, piv), b)}
\par\smallskip
\noindent 先做一次分解，之后更换任意右端项 $\boldsymbol{b}$ 都只需前代与回代两步。
\end{remk}

\begin{sidebar}{高斯与“消去法”的两千年}
消去法并非高斯首创。中国汉代的《九章算术》第八章“方程”已完整给出
用算筹在增广长方形数表上消元的算法，比高斯早约一千七百年。
高斯（1777—1855）是在用最小二乘法确定小行星谷神星轨道时
重新系统地使用了这一方法，并把它与对称正定方程组的快速解法结合起来，
因此今天以他的名字命名。1940 年代计算机出现后，
人们才意识到消元法的运算量必须精打细算，LU 分解由此成为标准记法。
\end{sidebar}

\begin{remk}{为什么要选最大的主元（部分主元法）}{pivot-stability}
若某步的主元恰好是 $0$ 或非常接近 $0$，直接消元要么无法继续、
要么产生巨大的乘数 $\ell_{ij}$，把舍入误差急剧放大。标准做法是
\textbf{部分主元法}：每一步在当前列中选绝对值最大的元素，
先做行交换再消元。这也解释了为什么实际软件中的 LU 分解总带一个
置换矩阵 $P$，形如 $PA=LU$。
\textbf{经验法则}：主元越小，解越“晃”——
第 2.5 节的条件数理论会把这句话变成可计算的判据。
\end{remk}

\begin{latab}{消元法的运算量与三种“分解记账法”（$n$ 阶一般矩阵）}
\centering\small
\begin{tabular}{@{}llll@{}}
\toprule
方法 & 分解形式 & 运算量 & 适用场景\\
\midrule
高斯消元（逐步） & 无 & $O(n^3)$ / 每个右端项 & 手算小方程组\\
LU 分解 & $PA=LU$ & 分解 $\tfrac23n^3$，回代 $O(n^2)$ & 多个右端项\\
Cholesky 分解 & $A=LL^{\mathsf T}$ & $\tfrac13n^3$（约省一半） & 对称正定矩阵\\
QR 分解 & $A=QR$ & $2n^3$ & 最小二乘（第七章）\\
\bottomrule
\end{tabular}
\end{latab}

\begin{chapkeys}
\small
\begin{multicols}{2}
\begin{itemize}[leftmargin=1.5em, itemsep=2pt]
\item 行变换不改变解集；三种变换均可逆。
\item 主元个数 = 秩；自由变量个数 $=n-\rank(A)$。
\item 简化行阶梯形唯一，行阶梯形不唯一。
\item 有解 $\Leftrightarrow$ 增广矩阵与系数矩阵同秩。
\item 解的结构：$\boldsymbol{x}_p+\Nul(A)$，即“特解 + 齐次解”。
\item 方阵可逆 $\Leftrightarrow$ 零空间只有零解 $\Leftrightarrow$ $\rank=n$。
\item $(AB)^{-1}=B^{-1}A^{-1}$，顺序颠倒。
\item LU 分解一次，多次右端项各付 $O(n^2)$。
\end{itemize}
\end{multicols}
\end{chapkeys}

\exhead
\begin{multicols}{2}
\small
\textbf{A 组（基础计算）}
\begin{enumerate}[leftmargin=1.6em, itemsep=3pt]
\item 用行变换把
      $\left[\begin{array}{ccc|c}1&2&1&3\\2&5&3&8\\1&1&-1&-1\end{array}\right]$
      化为简化行阶梯形，并写出通解。
\item 判断下列方程组是否有解、是否有唯一解：
      \[ \left\{\begin{aligned}x+2y-z&=1\\ 2x+4y-2z&=3\end{aligned}\right.,
      \qquad
      \left\{\begin{aligned}x+2y-z&=1\\ 2x+4y-2z&=2\end{aligned}\right. \]
\item 用行变换求 $A=\begin{bmatrix}1&2\\3&7\end{bmatrix}$ 的逆，
      并用它解 $A\boldsymbol{x}=(1,0)^{\mathsf T}$。
\item 求 $A=\begin{bmatrix}2&1&1\\4&3&3\\8&7&9\end{bmatrix}$ 的 LU 分解。
\end{enumerate}
\textbf{B 组（推理证明）}
\begin{enumerate}[leftmargin=1.6em, itemsep=3pt, start=5]
\item 证明：若 $A$ 可逆且 $AB=O$，则 $B=O$。举一个 $A$ 不可逆的反例。
\item 证明：可逆矩阵的逆唯一。（提示：设 $B,C$ 都是逆，考察 $B(AC)$。）
\item 设 $A$ 为 $n$ 阶方阵且 $\rank(A)=n-1$。证明 $\rank(A^2)$ 可为 $n-1$ 或 $n-2$，
      并各举一例。
\item 证明：上三角可逆矩阵的逆仍是上三角矩阵。
\end{enumerate}
\textbf{C 组（应用与建模）}
\begin{enumerate}[leftmargin=1.6em, itemsep=3pt, start=9]
\item 交通流量问题：某十字路口四条进入车道的流量满足线性方程组。
      自行设定流量并讨论“哪些流量是自由的，哪些是被约束的”。
\item 配平化学方程式 $a\,\mathrm{C_3H_8}+b\,\mathrm{O_2}\to c\,\mathrm{CO_2}+d\,\mathrm{H_2O}$：
      对 C、H、O 三种元素列线性方程组，求最小正整数解。
\end{enumerate}
\end{multicols}

% ===================== 第 3 章 行列式 =====================
\chapter{行列式}
\label{ch:det}

\begin{chapoutline}
\small
\begin{multicols}{2}
\textbf{学习目标}
\begin{itemize}[leftmargin=1.5em, itemsep=1.5pt]
\item 从有向面积出发理解二阶行列式；
\item 掌握行列式的十条基本性质并熟练运用于计算；
\item 会用余子式展开与行变换法求高阶行列式；
\item 理解行列式的几何意义：有向体积。
\end{itemize}
\columnbreak
\textbf{关键词}
\par\smallskip
有向面积\quad 置换\quad 余子式\quad 代数余子式\\
克拉默法则\quad 范德蒙德行列式
\par\smallskip
\textbf{建议学时}
\par\smallskip
8 学时
\end{multicols}
\end{chapoutline}

\section{从面积说起：二阶行列式}

\begin{defn}{二阶行列式}{det2}
\[
\det\begin{bmatrix}a & b\\ c & d\end{bmatrix}=ad-bc.
\]
\end{defn}

这个式子单看是枯燥的，但它的几何意义极其漂亮。
设 $A=\begin{bmatrix}a&b\\c&d\end{bmatrix}$ 的两列为
$\boldsymbol{u}=(a,c)^{\mathsf T}$，$\boldsymbol{v}=(b,d)^{\mathsf T}$，
则以 $\boldsymbol{u},\boldsymbol{v}$ 为邻边的平行四边形面积恰为 $|ad-bc|$；
而 $ad-bc$ 本身还带\textbf{符号}：当 $\boldsymbol{u}$ 逆时针转到 $\boldsymbol{v}$
时为正，顺时针为负。故二阶行列式是\textbf{有向面积}。
\begin{lafig}[>=Stealth, line cap=round]{二阶行列式的几何意义：以两列为边的平行四边形的有向面积}
  \begin{scope}
    \draw[line!80, very thin, step=0.5] (-0.6,-0.4) grid (2.9,2.9);
    \draw[qing!55, line width=0.7pt, ->] (-0.75,0) -- (3.05,0);
    \draw[qing!55, line width=0.7pt, ->] (0,-0.55) -- (0,3.05);
    \fill[qing!11] (0,0) -- (2.0,0.5) -- (3.0,2.5) -- (1.0,2.0) -- cycle;
    \draw[dai, line width=1.7pt, ->] (0,0) -- (2.0,0.5);
    \draw[zhe, line width=1.7pt, ->] (0,0) -- (1.0,2.0);
    \draw[dashed, ink!35] (2.0,0.5) -- (3.0,2.5) -- (1.0,2.0);
    \draw[qing, ->, line width=1.1pt] (1.15,1.35) arc[start angle=63, end angle=14, radius=1.3];
    \node[qing] at (1.55,0.92) {\small 逆时针：$+$};
    \node[font=\footnotesize, color=inkl] at (1.5,-0.95)
      {$\det A=2\cdot2-1\cdot0.5=+3.5>0$};
  \end{scope}
  \begin{scope}[xshift=6.6cm]
    \draw[line!80, very thin, step=0.5] (-0.6,-0.4) grid (2.9,2.9);
    \draw[qing!55, line width=0.7pt, ->] (-0.75,0) -- (3.05,0);
    \draw[qing!55, line width=0.7pt, ->] (0,-0.55) -- (0,3.05);
    \fill[zhe!11] (0,0) -- (1.0,2.0) -- (3.0,2.5) -- (2.0,0.5) -- cycle;
    \draw[dai, line width=1.7pt, ->] (0,0) -- (1.0,2.0);
    \draw[zhe, line width=1.7pt, ->] (0,0) -- (2.0,0.5);
    \draw[dashed, ink!35] (1.0,2.0) -- (3.0,2.5) -- (2.0,0.5);
    \draw[zhe, ->, line width=1.1pt] (0.55,1.1) arc[start angle=63, end angle=14, radius=1.3];
    \node[zhe] at (1.0,0.72) {\small 顺时针：$-$};
    \node[font=\footnotesize, color=inkl] at (1.5,-0.95)
      {交换两列，行列式变号};
  \end{scope}
\end{lafig}

\begin{remk}{行列式与可逆性}{det-invert}
$\det A\neq 0$ 恰好等价于“两列不共线”，也就是 $A$ 可逆。
因此\textbf{行列式是方阵可逆性的一个数值判据}：
面积被压成零，信息就被压没了。
\end{remk}

\section{排列定义与基本性质}

\begin{defn}{$n$ 阶行列式}{detn}
设 $A=(a_{ij})$ 为 $n$ 阶方阵。定义
\[
\det A=\sum_{\sigma\in S_n}\sign(\sigma)\,
a_{1\sigma(1)}a_{2\sigma(2)}\cdots a_{n\sigma(n)},
\]
其中 $S_n$ 是 $n$ 个数码的全部置换，$\sign(\sigma)=(-1)^{\operatorname{inv}(\sigma)}$
由逆序数 $\operatorname{inv}(\sigma)$ 决定。
\end{defn}

\begin{exmp}{三阶行列式按定义展开}{det3-def}
三阶情形有 $3!=6$ 项：
\[
\det\begin{bmatrix}a_{11}&a_{12}&a_{13}\\a_{21}&a_{22}&a_{23}\\a_{31}&a_{32}&a_{33}\end{bmatrix}
=a_{11}a_{22}a_{33}+a_{12}a_{23}a_{31}+a_{13}a_{21}a_{32}
-a_{13}a_{22}a_{31}-a_{11}a_{23}a_{32}-a_{12}a_{21}a_{33}.
\]
\end{exmp}

\begin{thm}{行列式的十条基本性质}{det-prop}
设 $A,B$ 为 $n$ 阶方阵，则
\begin{enumerate}[label=(\arabic*), itemsep=1pt]
\item $\det I=1$；
\item 交换两行（列），行列式变号；
\item 某行乘常数 $c$，行列式乘 $c$；
\item 某行加上另一行的 $c$ 倍，行列式不变；
\item 某行为两组数之和，则行列式可按该行拆成两个行列式之和；
\item 有两行（列）成比例或相同 $\Rightarrow\det A=0$；
\item $\det A^{\mathsf T}=\det A$（行与列地位平等）；
\item 上（下）三角矩阵的行列式等于对角元之积；
\item $A$ 可逆 $\Longleftrightarrow \det A\neq 0$，且 $\det A^{-1}=1/\det A$；
\item $\det(AB)=\det A\cdot\det B$（乘性！但 $\det(A+B)$ 无简单公式）。
\end{enumerate}
\end{thm}

\begin{remk}{用性质计算，而不是用定义}{det-calc}
性质 (2)(3)(4) 正是初等行变换的翻译。于是求高阶行列式的标准流程是：
\textbf{用行变换把矩阵化成三角阵，再取对角元之积}——
每做一次换行记一个负号，每做一次倍加不记号。
这比按定义展开（$n!$ 项）高效得多：$10!=3628800$。
\end{remk}

\begin{exmp}{行变换法求四阶行列式}{det4-ex}
\[
D=\begin{vmatrix}
2 & 1 & 1 & 1\\
1 & 2 & 1 & 1\\
1 & 1 & 2 & 1\\
1 & 1 & 1 & 2
\end{vmatrix}.
\]
各行元素之和都是 $5$，故把第 $2,3,4$ 列都加到第 $1$ 列，提出公因子 $5$：
\[
D=5\begin{vmatrix}1&1&1&1\\1&2&1&1\\1&1&2&1\\1&1&1&2\end{vmatrix}
=5\begin{vmatrix}1&1&1&1\\0&1&0&0\\0&0&1&0\\0&0&0&1\end{vmatrix}=5.
\]
\end{exmp}

\begin{exmp}{范德蒙德行列式}{vandermonde}
\[
V(x_1,\dots,x_n)=\det
\begin{bmatrix}
1 & 1 & \cdots & 1\\
x_1 & x_2 & \cdots & x_n\\
x_1^2 & x_2^2 & \cdots & x_n^2\\
\vdots & \vdots & & \vdots\\
x_1^{n-1} & x_2^{n-1} & \cdots & x_n^{n-1}
\end{bmatrix}
=\prod_{1\le i<j\le n}(x_j-x_i).
\]
证明思路：自下而上，用 $R_i\leftarrow R_i-x_1R_{i-1}$ 逐行消去第一列的
$x_1$ 的幂，再按第一列展开，归纳即得。
\end{exmp}

\section{余子式展开}

\begin{defn}{余子式与代数余子式}{cofactor}
设 $A$ 为 $n$ 阶方阵。划去第 $i$ 行第 $j$ 列所得的 $n-1$ 阶行列式记作 $M_{ij}$
（\textbf{余子式}），而
\[
A_{ij}=(-1)^{i+j}M_{ij}
\]
称为\textbf{代数余子式}。
\end{defn}

\begin{thm}{按行（列）展开}{laplace}
对任意固定的 $i$ 与 $j$：
\[
\det A=\sum_{k=1}^{n}a_{ik}A_{ik}\quad(\text{按第 }i\text{ 行展开})
=\sum_{k=1}^{n}a_{kj}A_{kj}\quad(\text{按第 }j\text{ 列展开}).
\]
进一步有\textbf{异乘变零}：$\sum_{k=1}^{n}a_{ik}A_{jk}=0\ (i\neq j)$。
\end{thm}

\begin{pf}
只证异乘变零。把 $A$ 的第 $j$ 行换成第 $i$ 行得到 $A'$，
则 $A'$ 有两行相同，故 $\det A'=0$。按第 $j$ 行展开 $A'$：
\[
0=\det A'=\sum_{k=1}^{n}a_{ik}A'_{jk}=\sum_{k=1}^{n}a_{ik}A_{jk},
\]
因为划去第 $j$ 行后 $A$ 与 $A'$ 完全相同。
\end{pf}

\begin{prop}{伴随矩阵}{adjugate}
$A^{*}=(A_{ji})_{n\times n}$ 称为\textbf{伴随矩阵}，满足
\[
A\,A^{*}=A^{*}A=(\det A)\,I,
\qquad
A^{-1}=\frac{1}{\det A}A^{*}\quad(\det A\neq0).
\]
\end{prop}

\begin{remk}{伴随矩阵适合理论，不适合计算}{adj-note}
用 $A^{-1}=\frac{1}{\det A}A^{*}$ 求 $n$ 阶逆需要 $n^2$ 个 $(n-1)$ 阶行列式，
运算量约 $O(n!)$。实用中请使用第 2 章的行变换法。
伴随矩阵的价值在于公式本身——它给出逆矩阵的\emph{封闭表达式}，
在理论推导（如克拉默法则）中不可替代。
\end{remk}

\section{克拉默法则}

\begin{thm}{克拉默法则}{cramer}
设 $A$ 为 $n$ 阶可逆矩阵，$\boldsymbol{b}\in\Rnum{n}$。则 $A\boldsymbol{x}=\boldsymbol{b}$
有唯一解，且
\[
x_j=\frac{\det A_j}{\det A},\qquad
A_j=[\boldsymbol{a}_1,\dots,\boldsymbol{a}_{j-1},\boldsymbol{b},
      \boldsymbol{a}_{j+1},\dots,\boldsymbol{a}_n],
\]
即把 $A$ 的第 $j$ 列换成 $\boldsymbol{b}$。
\end{thm}

\begin{pf}
在 $A\boldsymbol{x}=\boldsymbol{b}$ 两端左乘 $A^{*}$：
$(\det A)\boldsymbol{x}=A^{*}\boldsymbol{b}$。
考察第 $j$ 个分量：$\sum_i A_{ij}b_i=\sum_i A_{ij}(A\boldsymbol{x})_i
=\sum_i A_{ij}\sum_k a_{ik}x_k=\sum_k x_k\sum_i a_{ik}A_{ij}=x_j\det A$，
最后一步用了“异乘变零”。
\end{pf}

\begin{remk}{克拉默法则的地位}{cramer-note}
克拉默法则给出了解的\textbf{显式公式}，在理论分析（尤其是解对参数的依赖性）
中极有用。但它需要 $n+1$ 个行列式，运算量远高于消元法，
因此只适合 $n\le 3$ 的手算与理论研究。
\end{remk}

\section{行列式的几何意义：有向体积}

\begin{thm}{行列式 $=$ 有向体积}{det-vol}
设 $A$ 为 $n$ 阶方阵，其列向量为 $\boldsymbol{a}_1,\dots,\boldsymbol{a}_n$。
则 $\det A$ 等于这 $n$ 个向量张成的平行多面体的\textbf{有向}体积；
其绝对值等于该平行多面体的体积。
特别地，$|\det A|$ 是线性变换 $A$ 对体积的缩放因子。
\end{thm}

\begin{lafig}[>=Stealth, line cap=round]{行列式的几何意义：二维为有向面积，三维为有向体积；高维类推}
  % 左：二阶面积
  \begin{scope}
    \node[font=\small\sffamily, color=qing] at (1.6,2.75) {二维：有向面积};
    \draw[line!80, very thin, step=0.5] (-0.5,-0.3) grid (2.7,2.3);
    \fill[qing!10] (0,0) -- (1.0,0.3) -- (2.2,1.9) -- (1.2,1.6) -- cycle;
    \draw[dai, line width=1.6pt, ->] (0,0) -- (1.0,0.3);
    \draw[zhe, line width=1.6pt, ->] (0,0) -- (1.2,1.6);
    \draw[dashed, ink!35] (1.0,0.3) -- (2.2,1.9) -- (1.2,1.6);
    \node[dai] at (0.55,0.02) {\small $\boldsymbol{a}_1$};
    \node[zhe] at (0.62,1.02) {\small $\boldsymbol{a}_2$};
    \node[font=\footnotesize, color=inkl] at (1.3,-0.75)
      {$|\det A|=\,$平行四边形面积};
  \end{scope}
  % 右：三阶体积
  \begin{scope}[xshift=7.2cm, x={(0.95cm,0cm)}, y={(0.42cm,0.42cm)}, z={(0cm,1.0cm)}]
    \node[font=\small\sffamily, color=qing] at (1.3,2.85) {三维：有向体积};
    % 底面网格
    \foreach \i in {0,0.5,1.0,1.5}{
      \draw[line!65, very thin] (\i,-0.3,0) -- (\i,2.1,0);
      \draw[line!65, very thin] (-0.3,\i,0) -- (1.9,\i,0);
    }
    \coordinate (O)   at (0,0,0);
    \coordinate (A)   at (1.6,0,0);
    \coordinate (B)   at (0,1.5,0);
    \coordinate (C)   at (0,0,1.15);
    \coordinate (AB)  at (1.6,1.5,0);
    \coordinate (AC)  at (1.6,0,1.15);
    \coordinate (BC)  at (0,1.5,1.15);
    \coordinate (ABC) at (1.6,1.5,1.15);
    \fill[qing!7]  (O) -- (A) -- (AB) -- (B) -- cycle;
    \fill[qing!11] (O) -- (A) -- (AC) -- (C) -- cycle;
    \fill[qing!15] (O) -- (B) -- (BC) -- (C) -- cycle;
    \fill[qing!19] (C) -- (AC) -- (ABC) -- (BC) -- cycle;
    \fill[qing!22] (A) -- (AB) -- (ABC) -- (AC) -- cycle;
    \fill[qing!25] (B) -- (AB) -- (ABC) -- (BC) -- cycle;
    \draw[qing!75, line width=0.6pt]
      (O) -- (A) -- (AB) -- (B) -- cycle;
    \draw[qing!75, line width=0.6pt] (O) -- (C) -- (BC);
    \draw[qing!75, line width=0.6pt] (A) -- (AC) -- (ABC) -- (BC);
    \draw[qing!75, line width=0.6pt] (AB) -- (ABC);
    \draw[dai, line width=1.6pt, ->] (O) -- (A)
      node[midway, below, yshift=-2pt]{\small $\boldsymbol{a}_1$};
    \draw[lan, line width=1.6pt, ->] (O) -- (B)
      node[midway, right, xshift=1pt]{\small $\boldsymbol{a}_2$};
    \draw[zhe, line width=1.6pt, ->] (O) -- (C)
      node[midway, left, xshift=-1pt]{\small $\boldsymbol{a}_3$};
    \node[font=\footnotesize, color=inkl] at (1.3,-1.15)
      {$|\det A|=\,$平行六面体体积};
  \end{scope}
\end{lafig}

\begin{coro}{行列式的三条几何推论}{det-geo-cor}
\begin{enumerate}[label=(\arabic*), itemsep=1pt]
\item 两列共线 $\Rightarrow$ “体积”被压扁为零 $\Rightarrow\det A=0$；
\item 正交矩阵的 $|\det|=1$（旋转与反射不改变体积）；
\item $\det(AB)=\det A\,\det B$ 的几何版本：先缩放 $\det A$ 倍再缩放 $\det B$ 倍。
\end{enumerate}
\end{coro}

\begin{lafig}[>=Stealth, line cap=round]{三种初等行变换对面积（行列式）的影响：
倍加不变、对换反号、倍乘乘 $c$}
  % ============ 左：倍加（剪切，面积不变）============
  \begin{scope}
    \node[font=\small\sffamily, color=qing] at (1.35,2.55) {倍加：$r_2 \mathrel{+}= k\,r_1$};
    \fill[zhe!12] (0,0) -- (1.6,0) -- (1.6,1.1) -- (0,1.1) -- cycle;
    \fill[qing!14] (0,0) -- (1.6,0) -- (2.35,1.1) -- (0.75,1.1) -- cycle;
    \draw[dai, line width=1.5pt, ->] (0,0) -- (1.6,0)
      node[midway, below, yshift=-2pt]{\small $\boldsymbol{a}_1$};
    \draw[zhe, line width=1.3pt, ->] (0,0) -- (0.75,1.1)
      node[midway, left, xshift=-1pt]{\small $\boldsymbol{a}_2$};
    \draw[qing, line width=1.5pt, ->] (0,0) -- (2.35,1.1)
      node[midway, above right, xshift=-2pt]{\small $\boldsymbol{a}_2'\!=\boldsymbol{a}_2+k\boldsymbol{a}_1$};
    \draw[ink!55, dashed, line width=0.7pt] (0.75,1.1) -- (2.35,1.1);
    \node[font=\footnotesize\sffamily, color=inkl] at (1.35,-0.62)
      {剪切变形，面积不变：$\det$ 不变};
  \end{scope}
  % ============ 中：对换（面积不变，定向反号）============
  \begin{scope}[xshift=6.1cm]
    \node[font=\small\sffamily, color=zhe] at (1.35,2.55) {对换：$\boldsymbol{a}_1\leftrightarrow\boldsymbol{a}_2$};
    \fill[qing!14] (0,0) -- (1.6,0) -- (2.35,1.1) -- (0.75,1.1) -- cycle;
    \draw[qing!45, line width=0.7pt] (0,0) -- (1.6,0) -- (2.35,1.1) -- (0.75,1.1) -- cycle;
    \draw[dai, line width=1.5pt, ->] (0,0) -- (0.75,1.1)
      node[midway, above left, xshift=1pt]{\small $\boldsymbol{a}_1'$};
    \draw[zhe, line width=1.5pt, ->] (0,0) -- (1.6,0)
      node[midway, below, yshift=-2pt]{\small $\boldsymbol{a}_2'$};
    \draw[ink!50, dashed, line width=0.7pt] (0.75,1.1) -- (0,1.1);
    % 定向（绕行）箭头：逆时针 -> 顺时针
    \draw[ink!60, line width=0.8pt, ->] (0.42,0.47) arc[start angle=125, end angle=15, radius=0.51];
    \node[font=\scriptsize, color=inkl] at (1.0,0.14) {顺时针};
    \node[font=\footnotesize\sffamily, color=inkl] at (1.35,-0.62)
      {底与高不变，绕行方向反转：$\det$ 变号};
  \end{scope}
  % ============ 右：倍乘（面积放大 c 倍）============
  \begin{scope}[xshift=12.2cm]
    \node[font=\small\sffamily, color=dai] at (1.35,2.55) {倍乘：$r_2 \mathrel{*}= c$};
    \fill[zhe!8]  (0,0) -- (1.6,0) -- (1.6,1.1) -- (0,1.1) -- cycle;
    \fill[qing!16] (0,0) -- (1.6,0) -- (1.6,1.98) -- (0,1.98) -- cycle;
    \draw[dai, line width=1.5pt, ->] (0,0) -- (1.6,0)
      node[midway, below, yshift=-2pt]{\small $\boldsymbol{a}_1$};
    \draw[zhe, line width=1.3pt, ->] (0,0) -- (0,1.1)
      node[midway, left, xshift=-1pt]{\small $\boldsymbol{a}_2$};
    \draw[qing, line width=1.5pt, ->] (0,0) -- (0,1.98)
      node[midway, right, xshift=1pt]{\small $c\,\boldsymbol{a}_2$};
    \node[font=\footnotesize\sffamily, color=inkl] at (1.35,-0.62)
      {高变成 $c$ 倍：$\det$ 变为原来的 $c$ 倍};
  \end{scope}
\end{lafig}

\begin{keypoint}
把上图的观察写成公式，就是行变换法求行列式的全部理论依据：
\[
\text{倍加不变}\;(r_i \mathrel{+}= k r_j)\,;\qquad
\text{对换变号}\;(r_i \leftrightarrow r_j)\,;\qquad
\text{倍乘乘}\,c\;(r_i \mathrel{*}= c).
\]
消元把矩阵打成上三角形时，只有“对换变号”与“倍乘乘 $c$”两种干扰，
因此 $\det A = (\pm 1)\times(\text{主元之积})$——这正是本章算法框的出处。
\end{keypoint}

\begin{sidebar}{行列式一词的来历}
“determinant”（行列式）由德国数学家\textbf{高斯}在 1801 年的
《算术研究》中首次使用，但当时仅指二次型的判别式。
把它作为独立对象系统研究的是\textbf{柯西}（Cauchy，1812），
他引入了双下标记法与乘积公式 $\det(AB)=\det A\,\det B$。
日本数学家\textbf{关孝和}早在 1683 年就得到了消元法与行列式的雏形，
比莱布尼茨的同类工作早了十年。行列式在十九世纪曾是线性代数的中心；
二十世纪以后，“矩阵与变换”逐渐取代了它的中心地位，
但行列式作为“体积缩放因子”的几何身份从未动摇。
\end{sidebar}

\begin{chapkeys}
\small
\begin{multicols}{2}
\begin{itemize}[leftmargin=1.5em, itemsep=2pt]
\item 二阶行列式 $=$ 有向面积；$n$ 阶 $=$ 有向体积。
\item 交换两行变号；倍加一行不变；倍乘一行乘 $c$。
\item 上三角行列式 $=$ 对角元之积——这是行变换法的目标。
\item $\det A^{\mathsf T}=\det A$，行与列地位平等。
\item $\det(AB)=\det A\,\det B$；$\det(A^{-1})=1/\det A$。
\item 异乘变零：$\sum_k a_{ik}A_{jk}=0\ (i\neq j)$。
\item 克拉默法则：$x_j=\det A_j/\det A$。
\item $|\det A|$ 是线性变换的体积缩放因子。
\end{itemize}
\end{multicols}
\end{chapkeys}

\exhead
\begin{multicols}{2}
\small
\textbf{A 组（基础计算）}
\begin{enumerate}[leftmargin=1.6em, itemsep=3pt]
\item 计算下列行列式：
      \[ \begin{vmatrix}3&1\\ 5&2\end{vmatrix},\qquad
         \begin{vmatrix}1&2&3\\ 4&5&6\\ 7&8&9\end{vmatrix} \]
      \[ \begin{vmatrix}1&1&1\\ 1&2&4\\ 1&3&9\end{vmatrix}. \]
\item 用行变换求
      \[ D=\begin{vmatrix}1&2&3&4\\ 2&3&4&1\\ 3&4&1&2\\ 4&1&2&3\end{vmatrix}. \]
\item 求解方程 $\begin{vmatrix}x&1\\ 2&x-1\end{vmatrix}=0$。
\item 设 $A$ 为 $3$ 阶方阵且 $\det A=2$，求 $\det(2A)$、$\det(A^{-1})$、$\det(-A)$、$\det(A^{*})$。
\end{enumerate}
\textbf{B 组（推理证明）}
\begin{enumerate}[leftmargin=1.6em, itemsep=3pt, start=5]
\item 证明：奇数阶反对称矩阵的行列式为 $0$。
\item 证明：若 $A$ 是 $n$ 阶幂零矩阵，则 $\det A=0$。
\item 设 $\boldsymbol{a},\boldsymbol{b},\boldsymbol{c}\in\Rnum{3}$。证明
      $\det[\boldsymbol{a}+\boldsymbol{b},\boldsymbol{b}+\boldsymbol{c},\boldsymbol{c}+\boldsymbol{a}]
       =2\det[\boldsymbol{a},\boldsymbol{b},\boldsymbol{c}]$。
\item 用归纳法证明范德蒙德行列式公式。
\end{enumerate}
\textbf{C 组（应用与建模）}
\begin{enumerate}[leftmargin=1.6em, itemsep=3pt, start=9]
\item 平面上三点 $(x_1,y_1),(x_2,y_2),(x_3,y_3)$ 共线的充要条件是
      $\begin{vmatrix}x_1&y_1&1\\ x_2&y_2&1\\ x_3&y_3&1\end{vmatrix}=0$。
      用它判断 $(1,2),(3,5),(5,8)$ 是否共线，并说明该行列式的几何意义。
\item 用 NumPy 随机生成 $50\times 50$ 矩阵 $A$，比较
      $\det(A)$ 与 $\prod_i(\text{LU 分解得到的对角元})$，观察舍入误差的影响。
\end{enumerate}
\end{multicols}

% ===================== 第 4 章 向量空间 =====================
\chapter{向量空间与子空间}
\label{ch:space}

\begin{chapoutline}
\small
\begin{multicols}{2}
\textbf{学习目标}
\begin{itemize}[leftmargin=1.5em, itemsep=1.5pt]
\item 理解向量空间、子空间、基与维数的严格定义；
\item 掌握线性相关与线性无关的判定方法；
\item 求四个基本子空间并理解它们之间的正交互补关系；
\item 会做基变换与坐标转换，理解秩—零化度定理。
\end{itemize}
\columnbreak
\textbf{关键词}
\par\smallskip
子空间\quad 线性相关性\quad 基\quad 维数\\
零空间\quad 列空间\quad 秩\quad 坐标变换
\par\smallskip
\textbf{建议学时}
\par\smallskip
12 学时
\end{multicols}
\end{chapoutline}

\section{向量空间与子空间}

\begin{defn}{向量空间}{vspace}
设 $V$ 是一个非空集合，其上定义了加法与数乘。若对任意
$\boldsymbol{u},\boldsymbol{v}\in V$ 与 $c\in\R$ 都有
$\boldsymbol{u}+\boldsymbol{v}\in V$，$c\boldsymbol{u}\in V$（\textbf{封闭性}），
并且这两种运算满足八条运算律（结合律、交换律、分配律、单位元等），
则称 $V$ 为\textbf{向量空间}。
\end{defn}

\begin{exmp}{常见的向量空间}{vspace-ex}
\begin{itemize}[itemsep=1pt]
\item $\Rnum{n}$：$n$ 维列向量全体，最基本的空间；
\item $\R{m\times n}$：$m\times n$ 实矩阵全体；
\item $\mathcal{P}_k$：次数不超过 $k$ 的多项式全体；
\item $C[a,b]$：区间 $[a,b]$ 上的连续函数全体——\emph{函数也能作向量}。
\end{itemize}
\end{exmp}

\begin{defn}{子空间}{subspace}
设 $V$ 是向量空间，$W\subseteq V$ 非空。若 $W$ 对 $V$ 中的加法与数乘封闭，
则称 $W$ 为 $V$ 的\textbf{子空间}。特别地，$\{\boldsymbol{0}\}$ 与 $V$ 本身
是 $V$ 的平凡子空间。
\end{defn}

\begin{keypoint}
\textbf{判定子空间只需验证两条：}非零元相加仍在内、数乘仍在内。
（由此自动得到 $\boldsymbol{0}\in W$，因 $0\cdot\boldsymbol{v}=\boldsymbol{0}$。）
不必逐条检查八条运算律。
\end{keypoint}

\begin{exmp}{是或不是子空间}{sub-ex}
在 $\Rnum{2}$ 中：
\begin{itemize}[itemsep=1pt]
\item 直线 $y=2x$：\textbf{是}子空间（过原点且封闭）；
\item 直线 $y=2x+1$：\textbf{不是}（$\boldsymbol{0}$ 不在其上，也不封闭）；
\item 第一象限 $\{(x,y):x\ge0,\ y\ge0\}$：\textbf{不是}（$(-1)\cdot(1,1)=(-1,-1)$ 跑出去了）；
\item 集合 $\{(x,y):x^2+y^2=1\}$：\textbf{不是}（单位圆对加法不封闭）。
\end{itemize}
\end{exmp}

\begin{remk}{几何口诀：子空间必过原点}{through-origin}
$\Rnum{2}$ 的子空间只有三类：$\{\boldsymbol{0}\}$、过原点的直线、$\Rnum{2}$ 本身。
$\Rnum{3}$ 的子空间只有四类：$\{\boldsymbol{0}\}$、过原点的直线、过原点的平面、$\Rnum{3}$。
\textbf{不过原点的平面或直线绝不是子空间}——
尽管它们看起来“平坦”。
\end{remk}

\section{线性相关与线性无关}

\begin{defn}{线性相关与线性无关}{lin-dep}
向量组 $\boldsymbol{v}_1,\dots,\boldsymbol{v}_k$ 称为\textbf{线性相关}，
若存在\textbf{不全为零}的 $c_1,\dots,c_k$ 使
$c_1\boldsymbol{v}_1+\cdots+c_k\boldsymbol{v}_k=\boldsymbol{0}$；
否则称为\textbf{线性无关}。
\end{defn}

\begin{keypoint}
\textbf{把定义翻译成方程组：}线性相关 $\Leftrightarrow$
齐次方程组 $[\boldsymbol{v}_1\ \cdots\ \boldsymbol{v}_k]\,\boldsymbol{c}=\boldsymbol{0}$
有非零解 $\Leftrightarrow$ \textbf{自由变量存在} $\Leftrightarrow$
$\rank[\boldsymbol{v}_1\ \cdots\ \boldsymbol{v}_k]<k$。
\end{keypoint}

\begin{prop}{线性相关性的常用事实}{lin-facts}
\begin{enumerate}[label=(\arabic*), itemsep=1pt]
\item 单个向量 $\boldsymbol{v}$ 线性无关 $\Leftrightarrow\boldsymbol{v}\neq\boldsymbol{0}$；
\item 两个向量线性相关 $\Leftrightarrow$ 它们共线（成比例）；
\item 含零向量的向量组必线性相关；
\item \textbf{部分相关则整体相关}，\textbf{整体无关则部分无关}；
\item $n+1$ 个 $\Rnum{n}$ 中的向量必线性相关（“向量太多”）；
\item 无关向量组的\textbf{接长}仍无关，相关向量组的\textbf{截短}仍相关。
\end{enumerate}
\end{prop}

\begin{exmp}{判定三组向量}{lin-ex}
\[
\text{(a)}\ \Big\{\begin{bmatrix}1\\2\end{bmatrix},
                 \begin{bmatrix}2\\4\end{bmatrix}\Big\}
\ \text{相关（成比例）};\qquad
\text{(b)}\ \Big\{\begin{bmatrix}1\\2\end{bmatrix},
                 \begin{bmatrix}1\\3\end{bmatrix}\Big\}
\ \text{无关};\qquad
\text{(c)}\ \Big\{\begin{bmatrix}1\\0\\0\end{bmatrix},
   \begin{bmatrix}0\\1\\0\end{bmatrix},
   \begin{bmatrix}1\\1\\0\end{bmatrix}\Big\}\ \text{相关}.
\]
对 (c)：$\boldsymbol{v}_3=\boldsymbol{v}_1+\boldsymbol{v}_2$，
即 $(1,1,-1)^{\mathsf T}$ 是一组非零组合系数。
\end{exmp}

\begin{lafig}[>=Stealth, line cap=round, x={(0.92cm,0.44cm)}, y={(-0.50cm,0.46cm)}, z={(0cm,1cm)}]
{线性相关与无关的几何：左图三向量共面（相关），右图三向量不共面（无关）}
  % ================= 左：共面（相关）=================
  \begin{scope}
    \fill[qing!8] (0,0,0) -- (2.3,0.25,0) -- (2.75,2.05,0) -- (0.45,1.8,0) -- cycle;
    \draw[qing!55, line width=0.7pt] (0,0,0) -- (2.3,0.25,0) -- (2.75,2.05,0)
      -- (0.45,1.8,0) -- cycle;
    \foreach \t in {0.25,0.5,0.75}{
      \draw[qing!22, line width=0.45pt] (2.3*\t,0.25*\t,0) -- (0.45+2.3*\t,1.8+0.25*\t,0);
      \draw[qing!22, line width=0.45pt] (2.3*\t+0.45,1.8*\t,0) -- (2.3*\t,0,0);
    }
    \draw[dai, line width=1.5pt, ->] (0,0,0) -- (2.3,0.25,0)
      node[pos=0.62, below, yshift=-2pt]{\small $\boldsymbol{v}_1$};
    \draw[lan, line width=1.5pt, ->] (0,0,0) -- (0.45,1.8,0)
      node[pos=0.6, left, xshift=-2pt]{\small $\boldsymbol{v}_2$};
    \draw[zhe, line width=1.5pt, ->] (0,0,0) -- (1.5,1.0,0)
      node[pos=0.55, above right, xshift=-2pt]{\small $\boldsymbol{v}_3$};
    \draw[ink!60, dashed, line width=0.7pt, ->] (0,0,0) -- (1.5,1.0,0.95)
      node[pos=0.9, right, xshift=1pt]{\small $\boldsymbol{w}$};
    \node[font=\footnotesize\sffamily, color=inkl] at (1.15,0.95,-2.0)
      {$\boldsymbol{v}_3\in\Span\{\boldsymbol{v}_1,\boldsymbol{v}_2\}$：相关};
  \end{scope}
  % ================= 右：不共面（无关）=================
  \begin{scope}[xshift=6.4cm]
    \fill[qing!8] (0,0,0) -- (2.3,0.25,0) -- (2.75,2.05,0) -- (0.45,1.8,0) -- cycle;
    \draw[qing!55, line width=0.7pt] (0,0,0) -- (2.3,0.25,0) -- (2.75,2.05,0)
      -- (0.45,1.8,0) -- cycle;
    \draw[dai, line width=1.5pt, ->] (0,0,0) -- (2.3,0.25,0)
      node[pos=0.62, below, yshift=-2pt]{\small $\boldsymbol{v}_1$};
    \draw[lan, line width=1.5pt, ->] (0,0,0) -- (0.45,1.8,0)
      node[pos=0.6, left, xshift=-2pt]{\small $\boldsymbol{v}_2$};
    \draw[zhe, line width=1.5pt, ->] (0,0,0) -- (0.9,0.7,1.5)
      node[pos=0.72, above right, xshift=-1pt]{\small $\boldsymbol{v}_3$};
    \draw[ink!45, dashed, line width=0.6pt] (0,0,0) -- (0.9,0.7,0);
    \draw[qing!40, line width=0.5pt, dashed] (0.9,0.7,0) -- (0.9,0.7,1.5);
    \node[font=\footnotesize\sffamily, color=inkl] at (1.15,0.95,-2.0)
      {$\boldsymbol{v}_3$ 跳出平面：无关，且 $c_1\boldsymbol{v}_1+c_2\boldsymbol{v}_2
        +c_3\boldsymbol{v}_3=\boldsymbol{0}$ 只有零解};
  \end{scope}
\end{lafig}

\begin{remk}{“无关”为什么如此重要}{why-independent}
左图中 $\boldsymbol{v}_3$ 沿用 $\boldsymbol{v}_1,\boldsymbol{v}_2$ 就能造出来，
它作为“第三个坐标轴”是多余的——用它表示向量会有\emph{两套}坐标，
表达不再唯一。右图中 $\boldsymbol{v}_1,\boldsymbol{v}_2,\boldsymbol{v}_3$
互相“够不着”，任何一个向量都不能由其余向量造出，此时每个向量
都有\textbf{唯一}的分解方式。“唯一可表达”正是基与坐标概念的核心，
也是第 4.3 节的出发点。
\end{remk}

\section{基、维数与坐标}

\begin{defn}{基}{basis}
设 $V$ 是向量空间。若 $\mathcal{B}=\{\boldsymbol{v}_1,\dots,\boldsymbol{v}_k\}$
满足：(i) $\mathcal{B}$ 线性无关；(ii) $\Span\mathcal{B}=V$，
则称 $\mathcal{B}$ 为 $V$ 的一个\textbf{基}。基中向量的个数 $k$ 称为 $V$ 的
\textbf{维数}，记作 $\dim V$。
\end{defn}

\begin{thm}{基的等价刻画}{basis-thm}
设 $V$ 是 $k$ 维向量空间，则 $V$ 中任意 $k$ 个线性无关的向量必构成 $V$ 的基；
反之，$V$ 的任何基都恰含 $k$ 个向量（维数是良定义的）。
\end{thm}

\begin{defn}{坐标}{coord}
设 $\mathcal{B}=\{\boldsymbol{v}_1,\dots,\boldsymbol{v}_k\}$ 是 $V$ 的基，
$\boldsymbol{x}\in V$。存在唯一一组数 $c_1,\dots,c_k$ 使
$\boldsymbol{x}=c_1\boldsymbol{v}_1+\cdots+c_k\boldsymbol{v}_k$，
称 $[c_1,\dots,c_k]^{\mathsf T}$ 为 $\boldsymbol{x}$ 在基 $\mathcal{B}$ 下的\textbf{坐标}，
记作 $[\boldsymbol{x}]_{\mathcal{B}}$。
\end{defn}

\begin{lafig}[>=Stealth, line cap=round]{同一个向量 $\boldsymbol{x}$，在标准基下坐标是 $(2.75,2.25)$，
在新基 $\{\boldsymbol{b}_1,\boldsymbol{b}_2\}$ 下是 $(1.5,1)$}
  \draw[line!80, very thin, step=0.5] (-1.4,-1.4) grid (3.4,3.4);
  \draw[qing!55, line width=0.7pt, ->] (-1.6,0) -- (3.55,0);
  \draw[qing!55, line width=0.7pt, ->] (0,-1.6) -- (0,3.55);
  \draw[ink!55, line width=1.3pt, ->] (0,0) -- (1,0)
    node[below, yshift=-1pt]{\small $\boldsymbol{e}_1$};
  \draw[ink!55, line width=1.3pt, ->] (0,0) -- (0,1)
    node[left, xshift=-1pt]{\small $\boldsymbol{e}_2$};
  \draw[dai, line width=1.7pt, ->] (0,0) -- (2,0.5)
    node[below right, xshift=-3pt]{\small $\boldsymbol{b}_1$};
  \draw[zhe, line width=1.7pt, ->] (0,0) -- (-0.5,2)
    node[above left, xshift=3pt]{\small $\boldsymbol{b}_2$};
  \coordinate (P) at (2.75,2.25);
  \draw[dashed, ink!35] (0,0) -- (P) -- (2,0.5) -- cycle;
  \draw[dashed, ink!35] (0,0) -- (P) -- (-0.5,2) -- cycle;
  \draw[qing, line width=2.1pt, ->] (0,0) -- (P)
    node[right, xshift=1pt]{\small $\boldsymbol{x}$};
  \fill[qing] (P) circle (2.3pt);
  \node[font=\footnotesize, color=inkl] at (1.6,-1.15)
    {$\boldsymbol{x}=1.5\,\boldsymbol{b}_1+1\,\boldsymbol{b}_2
      \ \Rightarrow\ [\boldsymbol{x}]_{\mathcal{B}}=\begin{bmatrix}1.5\\1\end{bmatrix}$};
\end{lafig}

\begin{prop}{坐标变换公式}{coord-change}
设 $P=[\boldsymbol{b}_1\ \cdots\ \boldsymbol{b}_k]$（各列为新基向量），
则
\[
\boldsymbol{x}=P\,[\boldsymbol{x}]_{\mathcal{B}}
\qquad\Longleftrightarrow\qquad
[\boldsymbol{x}]_{\mathcal{B}}=P^{-1}\boldsymbol{x}.
\]
$P$ 称为\textbf{基变换矩阵}；当 $\mathcal{B}$ 也是 $\Rnum{n}$ 的基时 $P$ 可逆。
\end{prop}

\section{四个基本子空间}

设 $A\in\R{m\times n}$。与 $A$ 相关的有四个子空间，它们是整门课程的骨架。

\begin{latab}{$A\in\R{m\times n}$ 的四个基本子空间}
\small\centering
\renewcommand{\arraystretch}{1.42}
\begin{tabular}{@{}llll@{}}
\toprule
\sffamily 名称 & \sffamily 记号 & \sffamily 所在空间 & \sffamily 维数\\
\midrule
列空间   & $\Col(A)=\Span\{A\text{ 的列}\}$ & $\Rnum{m}$ & $\rank(A)=r$\\
零空间   & $\Nul(A)=\{\boldsymbol{x}:A\boldsymbol{x}=\boldsymbol{0}\}$ & $\Rnum{n}$ & $n-r$\\
行空间   & $\Row(A)=\Col(A^{\mathsf T})$ & $\Rnum{n}$ & $r$\\
左零空间 & $\Nul(A^{\mathsf T})$ & $\Rnum{m}$ & $m-r$\\
\bottomrule
\end{tabular}
\end{latab}

\begin{keypoint}
\textbf{秩—零化度定理}：$\rank(A)+\dim\Nul(A)=n$，
即 $\dim\Col(A)+\dim\Nul(A)=$ 列数。
它是“主元个数 $+$ 自由变量个数 $=$ 未知数个数”的直接翻译。
\end{keypoint}

\begin{thm}{正交互补}{orthocomp}
\[
\Row(A)\perp\Nul(A)\ \text{（同在 }\Rnum{n}\text{ 中），}
\qquad
\Col(A)\perp\Nul(A^{\mathsf T})\ \text{（同在 }\Rnum{m}\text{ 中）}.
\]
且维数恰好互补，因此
\[
\Rnum{n}=\Row(A)\oplus\Nul(A),\qquad
\Rnum{m}=\Col(A)\oplus\Nul(A^{\mathsf T}).
\]
\end{thm}

\begin{pf}
只证第一式。若 $\boldsymbol{x}\in\Nul(A)$，则 $A\boldsymbol{x}=\boldsymbol{0}$，
即 $A$ 的每一行与 $\boldsymbol{x}$ 的内积为零；而 $\Row(A)$ 是各行张成的空间，
故 $\Row(A)$ 中任一向量与 $\boldsymbol{x}$ 正交。
再由维数和为 $n$，得二者正交互补。
\end{pf}

\begin{lafig}[>=Stealth, line cap=round]{四个基本子空间示意图（Strang 图）：$A$ 把行空间一一映到列空间}
  % 左块：定义域
  \draw[qing!45, rounded corners=2mm] (-1.5,-1.5) rectangle (1.5,1.5);
  \node[font=\small\sffamily, color=qing] at (0,1.78) {定义域 $\Rnum{n}$};
  \draw[line width=1.1pt, qing!70, ->] (-1.3,0.55) -- (1.3,0.55)
     node[midway, above, yshift=-1pt]{\small $\Row(A)$ \ 维数 $r$};
  \draw[line width=1.1pt, dai!75, ->] (-1.3,-0.55) -- (1.3,-0.55)
     node[midway, below, yshift=1pt]{\small $\Nul(A)$ \ 维数 $n-r$};
  \node[font=\footnotesize, color=inkl] at (0,-1.85) {$\Row(A)\perp\Nul(A)$};
  % 映射箭头
  \draw[zhe, line width=2.0pt, ->] (1.75,0) -- (3.6,0)
     node[midway, above=2pt]{\small $A$};
  % 右块：值域
  \begin{scope}[xshift=5.4cm]
    \draw[zhe!45, rounded corners=2mm] (-1.5,-1.5) rectangle (1.5,1.5);
    \node[font=\small\sffamily, color=zhe] at (0,1.78) {值域 $\Rnum{m}$};
    \draw[line width=1.1pt, zhe!75, ->] (-1.3,0.55) -- (1.3,0.55)
       node[midway, above, yshift=-1pt]{\small $\Col(A)$ \ 维数 $r$};
    \draw[line width=1.1pt, lan!75, ->] (-1.3,-0.55) -- (1.3,-0.55)
       node[midway, below, yshift=1pt]{\small $\Nul(A^{\mathsf T})$ \ 维数 $m-r$};
    \node[font=\footnotesize, color=inkl] at (0,-1.85) {$\Col(A)\perp\Nul(A^{\mathsf T})$};
  \end{scope}
  \node[font=\footnotesize, color=inkl] at (2.7,-2.5)
    {行空间与零空间正交互补，列空间与左零空间正交互补};
\end{lafig}

\begin{exmp}{求一个矩阵的四个子空间}{four-spaces}
设
$A=\begin{bmatrix}1&2&1&3\\ 2&4&0&4\\ 1&2&1&3\end{bmatrix}$，
其简化行阶梯形为
$\begin{bmatrix}1&2&0&1\\ 0&0&1&2\\ 0&0&0&0\end{bmatrix}$，故 $r=2$。

\smallskip
\textbf{列空间}：主元在第 1、3 列，故取 $A$ 的第 1、3 列（\emph{注意取原矩阵的列}）：
$\Col(A)=\Span\{(1,2,1)^{\mathsf T},(1,0,1)^{\mathsf T}\}$。

\smallskip
\textbf{行空间}：取阶梯形的非零行：
$\Row(A)=\Span\{(1,2,0,1),(0,0,1,2)\}$。

\smallskip
\textbf{零空间}：解 $A\boldsymbol{x}=\boldsymbol{0}$。自由变量 $x_2=s$，$x_4=t$：
$x_1=-2s-t$，$x_3=-2t$，故
\[
\Nul(A)=\Span\Big\{\begin{bmatrix}-2\\1\\0\\0\end{bmatrix},
                    \begin{bmatrix}-1\\0\\-2\\1\end{bmatrix}\Big\},
\qquad \dim\Nul(A)=4-2=2.
\]

\smallskip
\textbf{左零空间}：$\dim\Nul(A^{\mathsf T})=3-2=1$，
由行变换 $R_2\leftarrow R_2-2R_1$，$R_3\leftarrow R_3-R_1$ 的组合系数得到
$(-1,-2,1)^{\mathsf T}$
（验证：$-1\cdot$行$1-2\cdot$行$2+1\cdot$行$3=\boldsymbol{0}$）。
\end{exmp}

\begin{remk}{两个易错点}{common-error}
\begin{enumerate}[label=(\alph*)]
\item 求列空间时，\textbf{必须取原矩阵的主元列}，而不是阶梯形里的列——
      行变换改变了列空间。
\item 行空间恰恰相反：\textbf{取阶梯形的非零行}即可，因为行变换不改变行空间。
\end{enumerate}
\end{remk}

\section{基变换与相似（初步）}

\begin{prop}{不同基下坐标之间的关系}{coord-rel}
设 $V$ 有两组基 $\mathcal{B},\mathcal{C}$，基变换矩阵分别为 $P,Q$
（各列为相应基向量），则
\[
[\boldsymbol{x}]_{\mathcal{C}}=Q^{-1}P\,[\boldsymbol{x}]_{\mathcal{B}}.
\]
\end{prop}

这里出现的 $Q^{-1}P$ 形式将在第五章“相似矩阵”中再次登场。
基变换并不改变向量本身，改变的只是描述它的“语言”——
这与同一几何对象在不同坐标系下有不同的方程，是同一件事。

\begin{sidebar}{格拉斯曼与维数的第一次严格化}
“维数”在今天看来理所当然，但在十九世纪中叶它还只是直观。
1844 年，德国语言学家出身的数学家\textbf{赫尔曼·格拉斯曼}
（Hermann Grassmann，1809—1877）出版《线性扩张论》，
首次把 $n$ 维线性空间作为抽象对象处理，给出了线性相关、基与维数的一般理论。
由于行文晦涩、又与当时主流的哈密顿四元数思路相左，
这部著作在生前几乎无人问津；直到皮亚诺于 1888 年给出公理化定义，
格拉斯曼的工作才被完全承认。今天线性代数教材中“向量空间”的公理体系，
正是沿着格拉斯曼—皮亚诺这条线索定型的。
\end{sidebar}

\begin{chapkeys}
\small
\begin{multicols}{2}
\begin{itemize}[leftmargin=1.5em, itemsep=2pt]
\item 子空间只需验证封闭性，且必含零向量。
\item 线性相关 $\Leftrightarrow$ 齐次方程组有非零解 $\Leftrightarrow$ 秩不足。
\item 基 $=$ 无关 $+$ 张成；基中向量个数 $=$ 维数。
\item 坐标变换：$\boldsymbol{x}=P[\boldsymbol{x}]_{\mathcal{B}}$。
\item 四个基本子空间维数：$r,\ n-r,\ r,\ m-r$。
\item 秩—零化度定理：$\rank(A)+\dim\Nul(A)=n$。
\item $\Row(A)\perp\Nul(A)$，$\Col(A)\perp\Nul(A^{\mathsf T})$。
\item 列空间取原矩阵主元列，行空间取阶梯形非零行。
\end{itemize}
\end{multicols}
\end{chapkeys}

\exhead
\begin{multicols}{2}
\small
\textbf{A 组（基础计算）}
\begin{enumerate}[leftmargin=1.6em, itemsep=3pt]
\item 判断下列集合是否为 $\Rnum{3}$ 的子空间并说明理由：
      (a) $\{(x,y,z):x+y+z=0\}$；(b) $\{(x,y,z):x+y+z=1\}$；
      (c) $\{(x,y,z):xy=0\}$；(d) $\{(x,2x,3x):x\in\R\}$。
\item 判断向量组 $(1,2,1)^{\mathsf T},(2,4,2)^{\mathsf T},(1,1,1)^{\mathsf T}$
      是否线性相关，若相关写出一组非零组合系数。
\item 求 $A=\begin{bmatrix}1&2&1\\ 2&4&2\\ 3&6&3\end{bmatrix}$
      的列空间、行空间、零空间及各维数。
\item 设 $\mathcal{B}=\{(1,1)^{\mathsf T},(1,-1)^{\mathsf T}\}$，
      求 $\boldsymbol{x}=(5,1)^{\mathsf T}$
      在 $\mathcal{B}$ 下的坐标。
\end{enumerate}
\textbf{B 组（推理证明）}
\begin{enumerate}[leftmargin=1.6em, itemsep=3pt, start=5]
\item 证明：若 $\boldsymbol{v}_1,\dots,\boldsymbol{v}_k$ 线性无关，
      且 $\boldsymbol{v}_{k+1}\notin\Span\{\boldsymbol{v}_1,\dots,\boldsymbol{v}_k\}$，
      则 $\boldsymbol{v}_1,\dots,\boldsymbol{v}_{k+1}$ 线性无关。
\item 证明：$\rank(A)=\rank(A^{\mathsf T})$。（提示：主元个数与非零行个数。）
\item 设 $A$ 为 $m\times n$ 实矩阵。证明：$A\boldsymbol{x}=\boldsymbol{0}$
      与 $A^{\mathsf T}A\boldsymbol{x}=\boldsymbol{0}$ 同解，
      进而 $\rank(A^{\mathsf T}A)=\rank(A)$。
\item 证明：$\dim(U\cap W)+\dim(U+W)=\dim U+\dim W$。
\end{enumerate}
\textbf{C 组（应用与建模）}
\begin{enumerate}[leftmargin=1.6em, itemsep=3pt, start=9]
\item 图像可视为 $\Rnum{mn}$ 中的向量。说明为什么“压缩图像”等价于
      在某个低维子空间中近似表示它。
\item 推荐系统中，用户—物品评分矩阵的“列空间”代表什么？
      试从四个基本子空间的角度解释协同过滤。
\end{enumerate}
\end{multicols}

% ===================== 第 5 章 线性变换 =====================
\chapter{线性变换}
\label{ch:trans}

\begin{chapoutline}
\small
\begin{multicols}{2}
\textbf{学习目标}
\begin{itemize}[leftmargin=1.5em, itemsep=1.5pt]
\item 理解线性映射的定义，学会“一眼看出”一个映射是否线性；
\item 掌握线性映射与矩阵的对应关系；
\item 会求核与像，并把秩—零化度定理迁移到映射语言；
\item 理解复合、逆变换与相似矩阵的几何含义。
\end{itemize}
\columnbreak
\textbf{关键词}
\par\smallskip
线性映射\quad 核\quad 像\quad 标准矩阵\\
复合\quad 可逆变换\quad 相似矩阵\quad 基表示
\par\smallskip
\textbf{建议学时}
\par\smallskip
10 学时
\end{multicols}
\end{chapoutline}

\section{什么是线性变换}

\begin{defn}{线性变换}{lin-map}
设 $V,W$ 是向量空间，$T:V\to W$ 是一个映射。若对任意
$\boldsymbol{u},\boldsymbol{v}\in V$ 与 $c\in\R$ 都有
\[
T(\boldsymbol{u}+\boldsymbol{v})=T(\boldsymbol{u})+T(\boldsymbol{v}),
\qquad
T(c\,\boldsymbol{u})=c\,T(\boldsymbol{u}),
\]
则称 $T$ 为\textbf{线性变换}（线性映射）。两条件可合并为
\[
T(c_1\boldsymbol{u}+c_2\boldsymbol{v})
=c_1T(\boldsymbol{u})+c_2T(\boldsymbol{v}).
\]
\end{defn}

\begin{keypoint}
\textbf{两条几何判据}：线性变换把直线映成直线（或压成一点），
并且\textbf{必须把原点映到原点}。因此平移、平移加旋转都不是线性变换。
\end{keypoint}

\begin{exmp}{线性与非线性}{linear-vs-not}
\begin{itemize}[itemsep=1pt]
\item $T(x,y)=(2x+y,\ x)$：\textbf{线性}；
\item $T(x,y)=(x+1,\ y)$：\textbf{不线性}（原点被搬到 $(1,0)$）；
\item $T(x,y)=(x^2,\ y)$：\textbf{不线性}（$T(2,0)=(4,0)\neq 2T(1,0)$）；
\item 求导算子 $D:f\mapsto f'$：\textbf{线性}
      （$(f+g)'=f'+g'$，$(cf)'=cf'$）；
\item 积分算子 $J:f\mapsto\int_0^x f(t)\,\mathrm{d}t$：\textbf{线性}。
\end{itemize}
\end{exmp}

\begin{prop}{线性变换的基本性质}{lin-prop}
设 $T$ 线性，则 $T(\boldsymbol{0})=\boldsymbol{0}$，
$T(-\boldsymbol{v})=-T(\boldsymbol{v})$，
且对任意线性组合有 $T\big(\sum c_i\boldsymbol{v}_i\big)=\sum c_iT(\boldsymbol{v}_i)$。
\end{prop}

\section{五种基本平面变换}

理解线性变换最好的方式，是观察它如何\textbf{扭曲一张网格}。

\begin{lafig}[>=Stealth, line cap=round]{五种基本平面线性变换：网格的形态揭示了矩阵的作用
（旋转 $\begin{psmallmatrix}\cos\theta&-\sin\theta\\ \sin\theta&\cos\theta\end{psmallmatrix}$、
剪切 $\begin{psmallmatrix}1&k\\ 0&1\end{psmallmatrix}$、
伸缩 $\begin{psmallmatrix}s_x&0\\ 0&s_y\end{psmallmatrix}$、
反射、投影）}
  \foreach \i in {-2,-1,0,1,2}{
    \draw[qing!22, very thin] (-1.7,\i*0.85) -- (1.7,\i*0.85);
    \draw[qing!22, very thin] (\i*0.85,-1.7) -- (\i*0.85,1.7);
  }
  \node[font=\footnotesize\sffamily, color=inkl] at (0,-2.05) {原网格};
  \draw[ink!60, line width=1.2pt, ->] (-1.7,0) -- (1.75,0);
  \draw[ink!60, line width=1.2pt, ->] (0,-1.7) -- (0,1.75);
  % ---------- 旋转 ----------
  \begin{scope}[xshift=4.7cm, rotate=32]
    \foreach \i in {-2,-1,0,1,2}{
      \draw[qing!30, very thin] (-1.5,\i*0.85) -- (1.5,\i*0.85);
      \draw[qing!30, very thin] (\i*0.85,-1.5) -- (\i*0.85,1.5);
    }
  \end{scope}
  \begin{scope}[xshift=4.7cm]
    \draw[dai, line width=1.6pt, ->] (0,0) -- (0.42,0.68);
    \draw[zhe, line width=1.6pt, ->] (0,0) -- (-0.68,0.42);
    \node[font=\footnotesize\sffamily, color=inkl] at (0,-2.05) {旋转};
  \end{scope}
  % ---------- 剪切 ----------
  \begin{scope}[xshift=9.4cm]
    \begin{scope}[cm={1,0.55,0,1,(0,0)}]
      \foreach \i in {-2,-1,0,1,2}{
        \draw[qing!30, very thin] (-1.5,\i*0.85) -- (1.5,\i*0.85);
        \draw[qing!30, very thin] (\i*0.85,-1.5) -- (\i*0.85,1.5);
      }
    \end{scope}
    \draw[dai, line width=1.6pt, ->] (0,0) -- (0.85,0);
    \draw[zhe, line width=1.6pt, ->] (0,0) -- (0.47,0.85);
    \node[font=\footnotesize\sffamily, color=inkl] at (0,-2.05) {剪切};
  \end{scope}
  % ---------- 伸缩 ----------
  \begin{scope}[yshift=-5.3cm]
    \begin{scope}[cm={1.45,0,0,0.62,(0,0)}]
      \foreach \i in {-2,-1,0,1,2}{
        \draw[qing!30, very thin] (-1.5,\i*0.85) -- (1.5,\i*0.85);
        \draw[qing!30, very thin] (\i*0.85,-1.5) -- (\i*0.85,1.5);
      }
    \end{scope}
    \draw[dai, line width=1.6pt, ->] (0,0) -- (1.23,0);
    \draw[zhe, line width=1.6pt, ->] (0,0) -- (0,0.53);
    \node[font=\footnotesize\sffamily, color=inkl] at (0,-2.05) {伸缩};
  \end{scope}
  % ---------- 反射 ----------
  \begin{scope}[xshift=4.7cm, yshift=-5.3cm, cm={0,-1,-1,0,(0,0)}]
    \foreach \i in {-2,-1,0,1,2}{
      \draw[qing!30, very thin] (-1.5,\i*0.85) -- (1.5,\i*0.85);
      \draw[qing!30, very thin] (\i*0.85,-1.5) -- (\i*0.85,1.5);
    }
  \end{scope}
  \begin{scope}[xshift=4.7cm, yshift=-5.3cm]
    \draw[ink!35, line width=0.8pt, dashed] (-1.6,-1.6) -- (1.6,1.6);
    \draw[dai, line width=1.6pt, ->] (0,0) -- (0,-0.8);
    \draw[zhe, line width=1.6pt, ->] (0,0) -- (-0.8,0);
    \node[font=\footnotesize\sffamily, color=inkl] at (0,-2.05) {反射};
  \end{scope}
  % ---------- 投影 ----------
  \begin{scope}[xshift=9.4cm, yshift=-5.3cm]
    \foreach \i in {-2,-1,0,1,2}{
      \draw[qing!30, very thin] (-1.5,\i*0.85) -- (1.5,\i*0.85);
      \draw[qing!30, very thin] (\i*0.85,-1.5) -- (\i*0.85,1.5);
    }
    \foreach \i in {-2,-1,0,1,2}{
      \draw[dai!85, line width=1.15pt] (-1.5,\i*0.85) -- (1.5,\i*0.85);
    }
    \draw[ink!45, line width=0.8pt, dashed, ->] (0.85,1.27) -- (0.85,0);
    \draw[ink!45, line width=0.8pt, dashed, ->] (0,1.7) -- (0,0);
    \node[font=\footnotesize\sffamily, color=inkl] at (0,-2.05) {投影};
  \end{scope}
\end{lafig}

\begin{remk}{投影是“有损”的}{proj-lossy}
前四种变换都可逆——信息没有丢失，面积最多缩放为 $|\det A|$ 倍。
投影则不可逆：一整条竖直方向的直线被压成同一个点，
面积缩放因子为 $\det=0$，因此矩阵不可逆。
“行列式为零”与“信息被压没”，说的是同一件事。
\end{remk}

\section{矩阵表示}

\begin{thm}{线性映射由基的像唯一决定}{rep-thm}
设 $T:\Rnum{n}\to\Rnum{m}$ 线性，$\{\boldsymbol{e}_1,\dots,\boldsymbol{e}_n\}$
是标准基。令 $A=[\,T(\boldsymbol{e}_1)\ \cdots\ T(\boldsymbol{e}_n)\,]\in\R{m\times n}$，
则对任意 $\boldsymbol{x}\in\Rnum{n}$，
\[
T(\boldsymbol{x})=A\boldsymbol{x}.
\]
$A$ 称为 $T$ 的\textbf{标准矩阵}。
\end{thm}

\begin{pf}
写 $\boldsymbol{x}=x_1\boldsymbol{e}_1+\cdots+x_n\boldsymbol{e}_n$，
由线性性与 $A\boldsymbol{x}$ 的列组合视角即得。
\end{pf}

\begin{keypoint}
\textbf{求标准矩阵的口诀}：把标准基向量一个个喂给 $T$，把得到的像依次排成列。
例如 $T(x,y)=(2x+y,x)$ 满足 $T(1,0)=(2,1)$，$T(0,1)=(1,0)$，
故 $A=\begin{bmatrix}2&1\\1&0\end{bmatrix}$。
\end{keypoint}

\begin{exmp}{旋转矩阵}{rot-mat}
把平面向量逆时针旋转 $\theta$ 角，则
$T(1,0)=(\cos\theta,\sin\theta)$，$T(0,1)=(-\sin\theta,\cos\theta)$，故
\[
R_\theta=\begin{bmatrix}\cos\theta&-\sin\theta\\ \sin\theta&\cos\theta\end{bmatrix},
\qquad
\det R_\theta=1,\qquad R_\theta^{\mathsf T}R_\theta=I.
\]
旋转是\textbf{正交矩阵}的典型代表：保长度、保角度、保面积。
\end{exmp}

\subsection{在一般基下的表示}

若把定义域的基换成 $\mathcal{B}$、值域的基换成 $\mathcal{C}$，则
\[
[T]_{\mathcal{C}\leftarrow\mathcal{B}}
=C^{-1}A P,
\]
其中 $P$ 的各列为 $\mathcal{B}$，$C$ 的各列为 $\mathcal{C}$。
特殊地，当 $V=W$ 且 $\mathcal{B}=\mathcal{C}$ 时，
$[T]_{\mathcal{B}}=P^{-1}AP$ ——这正是\textbf{相似矩阵}。

\section{核、像与维数公式}

\begin{defn}{核与像}{kernel-image}
设 $T:V\to W$ 线性。
\[
\ker T=\{\boldsymbol{v}\in V:T(\boldsymbol{v})=\boldsymbol{0}\}
\quad(\text{核，是 } V \text{ 的子空间}),
\]
\[
\operatorname{im}T=\{\,T(\boldsymbol{v}):\boldsymbol{v}\in V\,\}
\quad(\text{像，是 } W \text{ 的子空间}).
\]
$\dim\ker T$ 称为 $T$ 的\textbf{零度}，$\dim\operatorname{im}T$ 称为 $T$ 的\textbf{秩}。
\end{defn}

\begin{thm}{维数公式（秩—零化度定理的映射形式）}{dim-thm}
设 $T:V\to W$ 线性且 $\dim V<\infty$，则
\[
\dim\ker T+\dim\operatorname{im}T=\dim V.
\]
\end{thm}

\begin{pf}
取 $\ker T$ 的一组基 $\boldsymbol{u}_1,\dots,\boldsymbol{u}_k$，
扩充为 $V$ 的基 $\boldsymbol{u}_1,\dots,\boldsymbol{u}_k,
\boldsymbol{w}_1,\dots,\boldsymbol{w}_r$。
则 $T(\boldsymbol{w}_1),\dots,T(\boldsymbol{w}_r)$ 张成 $\operatorname{im}T$
且线性无关，故 $\dim\operatorname{im}T=r$，而 $k+r=\dim V$。
\end{pf}

\begin{remk}{维数公式的直觉}{dim-intuition}
$\dim V$ 个“自由度”被 $T$ 分成两部分：
一部分被压到零（核），一部分被完整保留（像）。
\textbf{压得越狠，像越“瘦”；核越大，像越小。}
当 $\ker T=\{\boldsymbol{0}\}$ 时 $T$ 是单射，此时像的维数最大。
\end{remk}

\begin{prop}{单射与满射的判据}{injective}
设 $T:\Rnum{n}\to\Rnum{m}$ 线性，标准矩阵为 $A$：
\[
T\ \text{单射}\Longleftrightarrow\ker T=\{\boldsymbol{0}\}
\Longleftrightarrow\rank A=n;
\qquad
T\ \text{满射}\Longleftrightarrow\operatorname{im}T=\Rnum{m}
\Longleftrightarrow\rank A=m.
\]
因此 $T$ 可逆 $\Leftrightarrow$ $m=n$ 且 $\rank A=n$。
\end{prop}

\begin{lafig}[>=Stealth, line cap=round]{秩—零化度定理的图示：定义域被 $T$ 分为“压扁的部分”
（核）与“保留的部分”，后者与像同构}
  % ================= 左：定义域 =================
  \begin{scope}
    \node[font=\small\sffamily, color=qing] at (1.5,2.45) {定义域 $V=\Rnum{3}$};
    \draw[qing!40, line width=0.8pt, rounded corners=3pt] (-1.4,-1.5) rectangle (4.0,2.0);
    % 核（一条直线）
    \fill[dai!12] (-0.9,-1.1) -- (0.6,-1.1) -- (0.6,-0.65) -- (-0.9,-0.65) -- cycle;
    \draw[dai!60, line width=0.6pt] (-0.9,-1.1) rectangle (0.6,-0.65);
    \node[dai, font=\scriptsize] at (-0.15,-0.875) {$\ker T$（压到零）};
    % 补空间
    \fill[qing!7] (-0.9,-0.45) rectangle (0.6,1.55);
    \draw[qing!45, line width=0.6pt] (-0.9,-0.45) rectangle (0.6,1.55);
    \node[qing, font=\scriptsize] at (-0.15,1.28) {$\ker T$ 的补};
    \foreach \y in {-0.2,0.15,0.5,0.85,1.2}{
      \draw[qing!35, line width=0.5pt, ->] (0.68,\y) -- (1.55,\y);
    }
    \node[font=\scriptsize, color=inkl] at (1.12,0.02) {$T$};
    % 像
    \fill[zhe!10] (1.65,-0.45) rectangle (3.55,1.55);
    \draw[zhe!55, line width=0.6pt] (1.65,-0.45) rectangle (3.55,1.55);
    \node[zhe, font=\scriptsize] at (2.6,1.28) {$\operatorname{im}T\subseteq W$};
    \node[font=\footnotesize\sffamily, color=inkl] at (1.3,-1.85)
      {$\dim V=3=\underbrace{1}_{\text{核}}+\underbrace{2}_{\text{秩}}$};
  \end{scope}
  % ================= 右：三类映射对照 =================
  \begin{scope}[xshift=9.0cm]
    \node[font=\small\sffamily, color=qing] at (2.2,2.45) {三类典型情形};
    \foreach \i/\idim/\col/\lab in {0/3/qing/{满秩：可逆}, 1/2/dai/{投影：压掉一维},
                                    2/1/zhe/{压缩：只剩一条线}}{
      \pgfmathsetmacro{\yy}{1.45-\i*1.12}
      \fill[white] (-0.75,\yy-0.46) rectangle (2.35,\yy+0.46);
      \draw[line!90, line width=0.5pt, rounded corners=2pt]
        (-0.75,\yy-0.46) rectangle (2.35,\yy+0.46);
      \node[font=\scriptsize\sffamily, text=ink, anchor=west] at (-0.58,\yy) {\lab};
      \fill[\col!14] (1.05,\yy-0.30) rectangle (1.05+0.34*\idim,\yy+0.30);
      \draw[\col!60, line width=0.5pt]
        (1.05,\yy-0.30) rectangle (1.05+0.34*\idim,\yy+0.30);
      \draw[ink!50, line width=0.5pt, ->] (0.62,\yy) -- (0.98,\yy);
    }
    \node[font=\footnotesize\sffamily, color=inkl] at (2.2,-2.15)
      {着色条宽度 $=$ 像的维数；压得越扁，像越窄};
  \end{scope}
\end{lafig}

\section{复合、逆与相似}

\begin{prop}{复合对应矩阵乘法}{compose}
设 $S:U\to V$，$T:V\to W$ 均线性，则 $T\circ S$ 线性，且
\[
[T\circ S]=[T]\,[S].
\]
\textbf{矩阵乘法之所以“那样”定义，正是为了让它表示复合}——
这是矩阵乘法最深刻的解释。
\end{prop}

\begin{remk}{顺序很重要}{order-matters}
$(T\circ S)(\boldsymbol{x})=T(S(\boldsymbol{x}))$，对应 $[T][S]$，
右边的先作用。几何上，“先剪切再旋转”的矩阵是 $R_\theta K$，
而“先旋转再剪切”是 $KR_\theta$，二者一般不等。
第一章那个 $AB\neq BA$ 的例子，如今有了完整的解释。
\end{remk}

\begin{defn}{相似矩阵}{similar}
设 $A,B$ 为 $n$ 阶方阵。若存在可逆矩阵 $P$ 使 $B=P^{-1}AP$，
则称 $A$ 与 $B$ \textbf{相似}，记作 $A\sim B$。
\end{defn}

\begin{keypoint}
相似 $=$ \textbf{同一个线性变换在不同基下的描述}。
相似矩阵共享一切“与基无关”的量：
\[
\det B=\det A,\qquad \tr B=\tr A,\qquad
\rank B=\rank A,\qquad
\text{特征多项式与特征值相同}.
\]
第六章将利用这一点，通过换基把矩阵化成最简单的形态——对角形。
\end{keypoint}

\begin{exmp}{用复合解释矩阵乘法}{compose-ex}
设 $A=\begin{bmatrix}1&1\\0&1\end{bmatrix}$（剪切），
$R=\begin{bmatrix}0&-1\\1&0\end{bmatrix}$（逆时针 $90^\circ$）。则
\[
RA=\begin{bmatrix}0&-1\\1&1\end{bmatrix}
\qquad\text{与}\qquad
AR=\begin{bmatrix}1&-1\\1&0\end{bmatrix}
\]
不同：前者是“先剪切后旋转”，后者是“先旋转后剪切”。
用第一章的网格图分别画出这两种效果，差异一目了然。
\end{exmp}

\begin{lafig}[>=Stealth, line cap=round]{乘法顺序的几何：左为 $RA$（先剪切再旋转），
右为 $AR$（先旋转再剪切），两图整体形态明显不同}
  % ============ 左：RA ============
  \begin{scope}
    \node[font=\small\sffamily, color=qing] at (1.7,2.6) {先剪切，后旋转\;($RA$)};
    \draw[line!85, very thin, step=0.55] (-2.2,-1.4) grid (2.2,2.2);
    \begin{scope}
      \clip (-2.35,-1.55) rectangle (2.35,2.25);
      \begin{scope}[cm={0,1,-1,1,(0,0)}]
        \foreach \i in {-1,0,1,2}{\draw[qing!40, line width=0.6pt] (\i,-1) -- (\i,2);}
        \foreach \j in {0,1,2}{\draw[qing!40, line width=0.6pt] (-1,\j) -- (2,\j);}
        \fill[qing!12] (0,0) rectangle (1,1);
        \draw[qing, line width=1.1pt] (0,0) rectangle (1,1);
      \end{scope}
    \end{scope}
    \draw[qing!45, line width=0.7pt, ->] (-2.4,0) -- (2.4,0);
    \draw[qing!45, line width=0.7pt, ->] (0,-1.6) -- (0,2.4);
    \node[font=\footnotesize\sffamily, color=inkl] at (1.7,-1.95)
      {单位方形 $\to$ 平行四边形 $\to$ 整体翻转立起};
  \end{scope}
  % ============ 右：AR ============
  \begin{scope}[xshift=7.0cm]
    \node[font=\small\sffamily, color=zhe] at (1.7,2.6) {先旋转，后剪切\;($AR$)};
    \draw[line!85, very thin, step=0.55] (-2.2,-1.4) grid (2.2,2.2);
    \begin{scope}
      \clip (-2.35,-1.55) rectangle (2.35,2.25);
      \begin{scope}[cm={1,1,-1,0,(0,0)}]
        \foreach \i in {-1,0,1,2}{\draw[zhe!40, line width=0.6pt] (\i,-1) -- (\i,2);}
        \foreach \j in {0,1,2}{\draw[zhe!40, line width=0.6pt] (-1,\j) -- (2,\j);}
        \fill[zhe!10] (0,0) rectangle (1,1);
        \draw[zhe, line width=1.1pt] (0,0) rectangle (1,1);
      \end{scope}
    \end{scope}
    \draw[qing!45, line width=0.7pt, ->] (-2.4,0) -- (2.4,0);
    \draw[qing!45, line width=0.7pt, ->] (0,-1.6) -- (0,2.4);
    \node[font=\footnotesize\sffamily, color=inkl] at (1.7,-1.95)
      {单位方形 $\to$ 立起 $\to$ 沿新方向剪切};
  \end{scope}
\end{lafig}

\begin{sidebar}{从“数表”到“算子”：一次观念的转身}
十九世纪的矩阵首先是一张\textbf{数表}：它来自行列式，
用于记录二次型或方程组的系数。真正把它看作\textbf{作用在空间上的算子}，
是二十世纪的事。1920 年代，量子力学发现物理态的演化由算子支配；
与此同时，诺特（E.\ Noether）与阿廷（E.\ Artin）把线性代数公理化，
从此“$A$”不再指代一堆具体的数字，而指代一个抽象的线性映射。
今天我们说“矩阵乘法不交换”，实质上是在说
“\emph{先旋转再翻转}不等于\emph{先翻转再旋转}”——
一句中学生也能听懂的话，只是它需要两千年才被写成公式。
\end{sidebar}

\begin{chapkeys}
\small
\begin{multicols}{2}
\begin{itemize}[leftmargin=1.5em, itemsep=2pt]
\item 线性变换必把原点映到原点，且把直线映成直线。
\item 标准矩阵：把标准基的像依次排成列。
\item $\dim\ker T+\dim\operatorname{im}T=\dim V$。
\item 单射 $\Leftrightarrow$ 核只有零解 $\Leftrightarrow$ $\rank A=n$。
\item 复合对应矩阵乘法，且顺序不可交换。
\item 相似矩阵共享行列式、迹、秩与特征值。
\item 投影不可逆，$\det=0$，信息被压没。
\item 旋转矩阵是正交矩阵，$\det=1$。
\end{itemize}
\end{multicols}
\end{chapkeys}

\exhead
\begin{multicols}{2}
\small
\textbf{A 组（基础计算）}
\begin{enumerate}[leftmargin=1.6em, itemsep=3pt]
\item 判断下列映射是否线性，并说明理由：
      (a) $T(x,y)=(x-y,\ 2x)$；(b) $T(x,y)=(x+y,\ 1)$；
      (c) $T(x,y,z)=(x,\ y,\ 0)$；(d) $T(x,y)=(|x|,\ y)$。
\item 求下列线性变换的标准矩阵：
      (a) $T(x,y)=(3x-2y,\ x+4y)$；(b) $T(x,y)=(y,\ x)$；
      (c) $T(x,y,z)=(x+y,\ y+z,\ z+x)$。
\item 求把平面向量关于直线 $y=x$ 作反射的标准矩阵，并验证 $A^2=I$。
\item 设 $T(x,y)=(2x+y,x)$，求 $\ker T$、$\operatorname{im}T$ 与二者的维数。
\end{enumerate}
\textbf{B 组（推理证明）}
\begin{enumerate}[leftmargin=1.6em, itemsep=3pt, start=5]
\item 证明：$T$ 线性 $\Leftrightarrow T(\boldsymbol{0})=\boldsymbol{0}$
      且 $T(\boldsymbol{u}-\boldsymbol{v})=T(\boldsymbol{u})-T(\boldsymbol{v})$。
\item 证明：线性映射 $T:V\to W$ 是单射 $\Leftrightarrow$
      它把线性无关组映为线性无关组。
\item 设 $A,B$ 相似。证明：$A^k$ 与 $B^k$ 相似，且 $A+cI$ 与 $B+cI$ 相似。
\item 证明：不存在线性变换 $T:\Rnum{2}\to\Rnum{3}$ 既是单射又是满射。
\end{enumerate}
\textbf{C 组（应用与建模）}
\begin{enumerate}[leftmargin=1.6em, itemsep=3pt, start=9]
\item 计算机图形学中，屏幕上的图形用齐次坐标 $(x,y,1)$ 表示。
      查阅资料，说明为什么引入三维齐次坐标后“平移”也能写成矩阵乘法，
      并写出平移矩阵的形式。
\item 用 NumPy 对一幅 $8\times 8$ 灰度图像分别施加剪切与投影，
      观察哪种变换无法恢复原图，并解释原因。
\end{enumerate}
\end{multicols}

% ===================== 第 6 章 特征值与对角化 =====================
\chapter{特征值、特征向量与对角化}
\label{ch:eig}

\begin{chapoutline}
\small
\begin{multicols}{2}
\textbf{学习目标}
\begin{itemize}[leftmargin=1.5em, itemsep=1.5pt]
\item 理解特征值与特征向量的几何含义；
\item 熟练求特征多项式、特征值与特征向量；
\item 掌握可对角化的充要条件与对角化步骤；
\item 会用对角化计算矩阵幂，处理动力系统与马尔可夫链问题。
\end{itemize}
\columnbreak
\textbf{关键词}
\par\smallskip
特征值\quad 特征向量\quad 特征多项式\\
代数重数\quad 几何重数\quad 正交对角化\quad 马尔可夫链
\par\smallskip
\textbf{建议学时}
\par\smallskip
12 学时（含上机 1 学时）
\end{multicols}
\end{chapoutline}

\section{变换中“不变的方向”}

一个线性变换会把大多数向量“掰弯”——改变它们的方向。
但总有少数特殊方向，变换之后\textbf{方向不变，只被拉伸}。

\begin{defn}{特征值与特征向量}{eig-def}
设 $A$ 为 $n$ 阶方阵。若存在非零向量 $\boldsymbol{v}$ 与数 $\lambda$ 使
\[
A\boldsymbol{v}=\lambda\boldsymbol{v},
\]
则称 $\lambda$ 为 $A$ 的\textbf{特征值}，$\boldsymbol{v}$ 为属于 $\lambda$ 的
\textbf{特征向量}。全体属于 $\lambda$ 的特征向量加上零向量构成
$\ker(A-\lambda I)$，称为\textbf{特征子空间}。
\end{defn}

\begin{keypoint}
\textbf{几何解读：}$A$ 作用在 $\boldsymbol{v}$ 上只是把它伸缩 $\lambda$ 倍。
$\lambda>0$ 同向拉伸，$\lambda<0$ 反向，$|\lambda|<1$ 压缩，$\lambda=0$ 压到原点。
因此\textbf{实方阵的实特征向量必然存在当且仅当该变换在某条直线上“只做伸缩”}。
\end{keypoint}

\begin{exmp}{一个直观例子}{eig-ex1}
$A=\begin{bmatrix}3&1\\0&2\end{bmatrix}$ 是上三角矩阵。
直接观察：$A(1,0)^{\mathsf T}=(3,0)^{\mathsf T}=3(1,0)^{\mathsf T}$，
故 $\lambda_1=3$，$\boldsymbol{v}_1=(1,0)^{\mathsf T}$；
$A(-1,1)^{\mathsf T}=(-2,2)^{\mathsf T}=2(-1,1)^{\mathsf T}$，
故 $\lambda_2=2$，$\boldsymbol{v}_2=(-1,1)^{\mathsf T}$。
\textbf{三角矩阵的特征值就是对角元}——这是一条应当记住的规律。
\end{exmp}

\begin{lafig}[>=Stealth, line cap=round]{特征向量与普通向量的对比：虚线为特征方向}
  \draw[line!80, very thin, step=0.5] (-1.6,-1.4) grid (3.2,2.6);
  \draw[qing!55, line width=0.7pt, ->] (-1.75,0) -- (3.35,0);
  \draw[qing!55, line width=0.7pt, ->] (0,-1.55) -- (0,2.75);
  % 特征方向（虚线）
  \draw[dai!55, dashed, line width=0.9pt] (-1.5,-0.75) -- (3.0,1.5);
  \draw[zhe!55, dashed, line width=0.9pt] (-0.75,-1.5) -- (1.5,3.0);
  % 特征向量
  \draw[dai, line width=1.9pt, ->] (0,0) -- (1.2,0.6)
     node[midway, above left, xshift=-2pt]{\small $\boldsymbol{v}_1$};
  \draw[dai!55, line width=1.9pt, ->] (0,0) -- (2.4,1.2)
     node[right, xshift=1pt]{\small $\lambda_1\boldsymbol{v}_1=2\boldsymbol{v}_1$};
  \draw[zhe, line width=1.9pt, ->] (0,0) -- (0.5,1.0)
     node[above left, xshift=1pt]{\small $\boldsymbol{v}_2$};
  \draw[zhe!55, line width=1.9pt, ->] (0,0) -- (-0.75,-1.5)
     node[below left, xshift=1pt]{\small $\lambda_2\boldsymbol{v}_2=-1.5\boldsymbol{v}_2$};
  % 普通向量被掰弯
  \draw[ink!55, line width=1.5pt, ->] (0,0) -- (1.6,1.6);
  \draw[ink!35, line width=1.5pt, ->] (0,0) -- (2.1,1.05);
  \node[font=\footnotesize, color=inkl] at (2.35,-1.0)
    {普通向量被“掰弯”；特征向量只被伸缩};
\end{lafig}

\section{特征多项式}

\begin{prop}{求特征值的标准流程}{eig-flow}
$A\boldsymbol{v}=\lambda\boldsymbol{v}$（$\boldsymbol{v}\neq\boldsymbol{0}$）
$\Longleftrightarrow(\lambda I-A)\boldsymbol{v}=\boldsymbol{0}$ 有非零解
$\Longleftrightarrow\det(\lambda I-A)=0$。称
\[
p_A(\lambda)=\det(\lambda I-A)
\]
为 $A$ 的\textbf{特征多项式}（$n$ 次，首项系数 $1$）。
\end{prop}

\begin{thm}{特征值的基本量}{eig-quant}
设 $A$ 的 $n$ 个特征值为 $\lambda_1,\dots,\lambda_n$（计重数），则
\[
\lambda_1+\lambda_2+\cdots+\lambda_n=\tr A,\qquad
\lambda_1\lambda_2\cdots\lambda_n=\det A.
\]
因此 $A$ 可逆 $\Leftrightarrow 0$ 不是 $A$ 的特征值。
\end{thm}

\begin{pf}
只证 $2$ 阶情形。$p_A(\lambda)=\lambda^2-(a_{11}+a_{22})\lambda+\det A$，
其两根之和为 $a_{11}+a_{22}=\tr A$，两根之积为 $\det A$。一般情形
由 $p_A(\lambda)=\lambda^n-\tr A\,\lambda^{n-1}+\cdots+(-1)^n\det A$ 归纳可得。
\end{pf}

\begin{defn}{代数重数与几何重数}{mult-def}
$\lambda_0$ 作为 $p_A(\lambda)$ 根的重数称为\textbf{代数重数}；
特征子空间 $\ker(\lambda_0 I-A)$ 的维数称为\textbf{几何重数}。
恒有
\[
1\le \text{几何重数}\le \text{代数重数}.
\]
\end{defn}

\begin{exmp}{几何重数小于代数重数}{mult-ex}
$A=\begin{bmatrix}2&1\\0&2\end{bmatrix}$：$p_A(\lambda)=(\lambda-2)^2$，
代数重数为 $2$。但
$A-2I=\begin{bmatrix}0&1\\0&0\end{bmatrix}$ 的秩为 $1$，
故几何重数 $=2-1=1<2$。这条“缺口”正是 $A$ 不可对角化的原因。
\end{exmp}

\begin{prop}{不同特征值的特征向量线性无关}{eig-indep}
属于\textbf{互不相同}特征值的特征向量必线性无关。
\end{prop}

\section{对角化}

\begin{defn}{可对角化}{diag-def}
若存在可逆矩阵 $P$ 使 $P^{-1}AP=D$ 为对角矩阵，则称 $A$ \textbf{可对角化}。
\end{defn}

\begin{thm}{可对角化的充要条件}{diag-thm}
$n$ 阶方阵 $A$ 可对角化
$\Longleftrightarrow$ $A$ 有 $n$ 个线性无关的特征向量
$\Longleftrightarrow$ 每个特征值的几何重数等于其代数重数。
此时 $P$ 的各列为特征向量，$D$ 的对角元为相应特征值，即
\[
AP=PD,\qquad A=PDP^{-1},\qquad D=P^{-1}AP.
\]
\end{thm}

\begin{pf}
把 $AP=PD$ 按列拆开：$A\boldsymbol{p}_i=\lambda_i\boldsymbol{p}_i$。
因此 $P$ 可逆且各列为特征向量 $\Leftrightarrow$ 存在 $n$ 个无关特征向量。
\end{pf}

\begin{keypoint}
\textbf{对角化三步法：}\quad
\textbf{第一步}：解 $\det(\lambda I-A)=0$，得全部特征值与各重数；
\textbf{第二步}：对每个 $\lambda_i$ 解 $(\lambda_i I-A)\boldsymbol{x}=\boldsymbol{0}$，
得特征子空间的基；
\textbf{第三步}：把所有特征向量按列排成 $P$，特征值按同样顺序排成 $D$。
验证 $AP=PD$（比算 $P^{-1}$ 便宜得多）。
\end{keypoint}

\begin{exmp}{完整的对角化}{diag-ex}
$A=\begin{bmatrix}1&2\\2&1\end{bmatrix}$。
\[
p_A(\lambda)=\det\begin{bmatrix}\lambda-1&-2\\-2&\lambda-1\end{bmatrix}
=(\lambda-1)^2-4=(\lambda-3)(\lambda+1).
\]
$\lambda_1=3$：$(A-3I)=\begin{bmatrix}-2&2\\2&-2\end{bmatrix}$，
得 $\boldsymbol{p}_1=(1,1)^{\mathsf T}$；
$\lambda_2=-1$：$(A+I)=\begin{bmatrix}2&2\\2&2\end{bmatrix}$，
得 $\boldsymbol{p}_2=(1,-1)^{\mathsf T}$。于是
\[
P=\begin{bmatrix}1&1\\1&-1\end{bmatrix},\quad
D=\begin{bmatrix}3&0\\0&-1\end{bmatrix},\quad
P^{-1}=\tfrac12\begin{bmatrix}1&1\\1&-1\end{bmatrix}.
\]
验证 $AP=\begin{bmatrix}3&-1\\3&-1\end{bmatrix}=PD$。\ $\checkmark$
\end{exmp}

\begin{prop}{矩阵幂的计算}{matpower}
若 $A=PDP^{-1}$，则
\[
A^k=PD^kP^{-1},\qquad
D^k=\diag(\lambda_1^k,\dots,\lambda_n^k).
\]
更一般地，$f(A)=P\,f(D)\,P^{-1}$ 对任意多项式 $f$ 成立。
\end{prop}

\begin{exmp}{斐波那契数列的封闭公式}{fib}
斐波那契递推 $F_{k+1}=F_k+F_{k-1}$ 可写成
$\begin{bmatrix}F_{k+1}\\F_k\end{bmatrix}
=A\begin{bmatrix}F_k\\F_{k-1}\end{bmatrix}$，
其中 $A=\begin{bmatrix}1&1\\1&0\end{bmatrix}$。
$A$ 的特征值为 $\dfrac{1\pm\sqrt5}{2}$（黄金比 $\varphi$ 与 $-1/\varphi$），
于是
\[
F_k=\frac{1}{\sqrt5}\left[\Big(\frac{1+\sqrt5}{2}\Big)^{k}
-\Big(\frac{1-\sqrt5}{2}\Big)^{k}\right].
\]
一个整数数列竟由两个无理数精确描述——这是对角化最令人惊讶的应用之一。
\end{exmp}

\begin{remk}{不是所有方阵都能对角化}{not-all-diag}
实\textbf{对称}矩阵一定可对角化（见下节）；
一般实方阵未必：若某个特征值的几何重数小于代数重数，
线性无关特征向量不够 $n$ 个，就“缺一不可”。
补救办法是若尔当标准形或 Schur 分解，本讲义不展开，
但应记住：\textbf{缺的从来不是特征值，而是特征向量的个数}。
\end{remk}

\begin{lafig}[>=Stealth, line cap=round]{对角化的几何：在特征基下，同一变换只是“沿两条特征轴的独立拉伸”}
  % ============ 左：标准基下的 A ============
  \begin{scope}
    \node[font=\small\sffamily, color=qing] at (1.9,2.55) {标准基下：$A=\begin{bmatrix}2&1\\ 1&2\end{bmatrix}$};
    \draw[line!85, very thin, step=0.5] (-1.8,-1.6) grid (2.2,2.0);
    \begin{scope}
      \clip (-2.0,-1.75) rectangle (2.35,2.15);
      \begin{scope}[cm={2,1,1,2,(0,0)}]
        \foreach \i in {-1,0,1}{\draw[qing!38, line width=0.6pt] (\i,-0.9) -- (\i,1.1);}
        \foreach \j in {0,1}{\draw[qing!38, line width=0.6pt] (-0.9,\j) -- (1.1,\j);}
        \fill[qing!10] (0,0) rectangle (1,1);
        \draw[qing, line width=1.1pt] (0,0) rectangle (1,1);
      \end{scope}
    \end{scope}
    \draw[qing!45, line width=0.7pt, ->] (-2.0,0) -- (2.4,0);
    \draw[qing!45, line width=0.7pt, ->] (0,-1.8) -- (0,2.2);
    \node[font=\footnotesize\sffamily, color=inkl] at (1.9,-1.95)
      {方形被拉伸的同时发生歪斜};
  \end{scope}
  % ============ 中：特征方向 ============
  \begin{scope}[xshift=6.6cm]
    \node[font=\small\sffamily, color=zhe] at (1.9,2.55) {特征方向：$\boldsymbol{u}_1,\boldsymbol{u}_2$};
    \draw[line!85, very thin, step=0.5] (-1.8,-1.6) grid (2.2,2.0);
    \draw[qing!45, line width=0.7pt, ->] (-2.0,0) -- (2.4,0);
    \draw[qing!45, line width=0.7pt, ->] (0,-1.8) -- (0,2.2);
    \draw[zhe, line width=1.6pt, ->] (-1.5,-1.5) -- (1.7,1.7)
      node[pos=0.86, above right, xshift=-2pt, inner sep=1pt]{\small $\boldsymbol{u}_1$};
    \draw[dai, line width=1.6pt, ->] (-1.7,1.7) -- (1.5,-1.5)
      node[pos=0.86, below right, xshift=-2pt, inner sep=1pt]{\small $\boldsymbol{u}_2$};
    \node[font=\footnotesize, color=zhe] at (2.0,1.35) {$\lambda_1=3$};
    \node[font=\footnotesize, color=dai] at (2.0,-1.30) {$\lambda_2=1$};
    \node[font=\footnotesize\sffamily, color=inkl] at (1.9,-1.95)
      {两条互相垂直的“不变方向”};
  \end{scope}
  % ============ 右：特征基下 ============
  \begin{scope}[xshift=13.2cm]
    \node[font=\small\sffamily, color=dai] at (1.9,2.55) {特征基下：$D=\diag(3,1)$};
    \draw[line!85, very thin, step=0.5] (-1.8,-1.6) grid (2.2,2.0);
    \begin{scope}[rotate=45]
      \foreach \i in {-1,0,1,2}{\draw[dai!40, line width=0.6pt] (0.5*\i,-1.3) -- (0.5*\i,1.3);}
      \foreach \j in {-1,0,1}{\draw[dai!40, line width=0.6pt] (-1.6,0.5*\j) -- (1.6,0.5*\j);}
      \fill[dai!14] (0,0) rectangle (1.5,0.5);
      \draw[dai, line width=1.1pt] (0,0) rectangle (1.5,0.5);
    \end{scope}
    \draw[qing!45, line width=0.7pt, ->] (-2.0,0) -- (2.4,0);
    \draw[qing!45, line width=0.7pt, ->] (0,-1.8) -- (0,2.2);
    \node[font=\footnotesize\sffamily, color=inkl] at (1.9,-1.95)
      {沿特征轴纯拉伸（$3$ 倍与 $1$ 倍）：无歪斜};
  \end{scope}
\end{lafig}

\section{实对称矩阵的正交对角化}

\begin{thm}{谱定理}{spectral}
设 $A$ 为实对称矩阵，则
\begin{enumerate}[label=(\arabic*), itemsep=1pt]
\item $A$ 的特征值全为实数；
\item 不同特征值的特征向量相互正交；
\item 存在\textbf{正交}矩阵 $Q$（$Q^{\mathsf T}Q=I$）使
      $Q^{\mathsf T}AQ=D$ 为对角矩阵。
\end{enumerate}
\end{thm}

\begin{pfsketch}
只证 (2)。设 $A\boldsymbol{u}=\lambda\boldsymbol{u}$，
$A\boldsymbol{v}=\mu\boldsymbol{v}$，$\lambda\neq\mu$。由对称性
\[
\lambda\,\boldsymbol{u}^{\mathsf T}\boldsymbol{v}
=(A\boldsymbol{u})^{\mathsf T}\boldsymbol{v}
=\boldsymbol{u}^{\mathsf T}A\boldsymbol{v}
=\mu\,\boldsymbol{u}^{\mathsf T}\boldsymbol{v},
\]
故 $(\lambda-\mu)\boldsymbol{u}^{\mathsf T}\boldsymbol{v}=0$，
即 $\boldsymbol{u}\perp\boldsymbol{v}$。(3) 的证明需要格拉姆—施密特正交化，
见第七章。
\end{pfsketch}

\begin{keypoint}{对称矩阵的对角化流程}{sym-steps}
\textbf{①} 求全部特征值；\textbf{②} 对每个特征值求特征向量；
\textbf{③} 同一特征值内的特征向量用格拉姆—施密特正交化并单位化；
\textbf{④} 把单位正交特征向量按列排成 $Q$。
\end{keypoint}

\begin{exmp}{正交对角化}{sym-ex}
$A=\begin{bmatrix}2&1\\1&2\end{bmatrix}$，
$p_A(\lambda)=(\lambda-1)(\lambda-3)$。
$\lambda=1$：$\boldsymbol{u}_1=\frac{1}{\sqrt2}(1,-1)^{\mathsf T}$；
$\lambda=3$：$\boldsymbol{u}_2=\frac{1}{\sqrt2}(1,1)^{\mathsf T}$。二者已正交，故
\[
Q=\frac{1}{\sqrt2}\begin{bmatrix}1&1\\-1&1\end{bmatrix},
\qquad
Q^{\mathsf T}AQ=\begin{bmatrix}1&0\\0&3\end{bmatrix}.
\]
几何上：先把坐标轴转到特征方向，再分别缩放 $1$ 倍与 $3$ 倍。
\end{exmp}

\section{应用一：离散动力系统}

\begin{thm}{迭代与主特征值}{dynamical}
设 $A$ 可对角化，$A=PDP^{-1}$，考察迭代 $\boldsymbol{x}_{k}=A^k\boldsymbol{x}_0$。
记 $\lambda_1$ 为模最大的特征值（\textbf{主特征值}），则
\[
\boldsymbol{x}_k=c_1\lambda_1^k\boldsymbol{p}_1
+c_2\lambda_2^k\boldsymbol{p}_2+\cdots+c_n\lambda_n^k\boldsymbol{p}_n,
\]
当 $|\lambda_1|>|\lambda_j|\ (j\ge2)$ 时，
\[
\frac{\boldsymbol{x}_k}{\lambda_1^{k}}\longrightarrow c_1\boldsymbol{p}_1
\qquad(k\to\infty).
\]
即：\textbf{无论初值如何，迭代最终都“躺”在主特征方向上}。
\end{thm}

\begin{lafig}[>=Stealth, line cap=round]{迭代序列收敛到主特征方向（幂法的原理）}
  \draw[line!80, very thin, step=0.5] (-1.3,-1.3) grid (3.0,2.2);
  \draw[qing!55, line width=0.7pt, ->] (-1.45,0) -- (3.15,0);
  \draw[qing!55, line width=0.7pt, ->] (0,-1.45) -- (0,2.35);
  \draw[dai!45, dashed, line width=0.9pt] (-1.2,-0.4) -- (2.85,0.95);
  % 迭代点列
  \fill[qing]  (0.55,0.85)  circle (2.2pt);
  \fill[qing]  (1.05,0.62)  circle (2.2pt);
  \fill[qing]  (1.45,0.62)  circle (2.2pt);
  \fill[qing]  (1.8,0.70)   circle (2.2pt);
  \fill[qing]  (2.1,0.78)   circle (2.2pt);
  \draw[qing!60, line width=0.8pt] (0.55,0.85) -- (1.05,0.62) -- (1.45,0.62)
      -- (1.8,0.70) -- (2.1,0.78);
  \node[qing] at (0.32,1.02) {\small $\boldsymbol{x}_0$};
  \node[qing] at (2.35,1.02) {\small $\boldsymbol{x}_5$};
  \draw[dai, line width=1.7pt, ->] (0,0) -- (2.55,0.85)
     node[right, xshift=2pt]{\small $\boldsymbol{p}_1$};
  \node[font=\footnotesize, color=inkl] at (0.9,-1.05)
    {$\boldsymbol{x}_{k}=A^k\boldsymbol{x}_0$ 逐步靠近主特征方向};
\end{lafig}

\begin{keypoint}{幂法：求主特征值的实用算法}{power-method}
任取非零 $\boldsymbol{x}_0$，反复计算 $\boldsymbol{x}_{k+1}=A\boldsymbol{x}_k
/\norm{A\boldsymbol{x}_k}$。归一化的目的是防止溢出。
收敛后 $\boldsymbol{x}_k\approx\pm\boldsymbol{p}_1$，
且 $\norm{A\boldsymbol{x}_k}\approx|\lambda_1|$。
这是许多大规模特征值算法（如 PageRank）的原型。
\end{keypoint}

\section{应用二：马尔可夫链}

\begin{defn}{转移矩阵}{markov}
若矩阵 $P\ge 0$ 且每列（或每行）元素之和为 $1$，则称 $P$ 为\textbf{随机矩阵}。
以 $P$ 为转移规则的系统称为\textbf{马尔可夫链}。
\end{defn}

\begin{thm}{平稳分布的存在性}{stationary}
设 $P$ 为 $n$ 阶随机矩阵且存在 $k$ 使 $P^k>0$（所有元素为正），
则 $\lambda=1$ 是 $P$ 的特征值且为单根，其余特征值满足 $|\lambda|<1$；
对应的特征向量（归一化后分量非负、总和为 $1$）称为\textbf{平稳分布} $\boldsymbol{\pi}$，
满足
\[
P\boldsymbol{\pi}=\boldsymbol{\pi},
\qquad
P^k\boldsymbol{x}_0\to\boldsymbol{\pi}\quad(k\to\infty).
\]
\end{thm}

\begin{exmp}{三城市人口迁移}{markov-ex}
某地区每年人口在城市 A、B、乡村 C 之间迁移，转移概率（列表示迁出地）：
\[
P=\begin{bmatrix}
0.85 & 0.05 & 0.10\\
0.05 & 0.90 & 0.05\\
0.10 & 0.05 & 0.85
\end{bmatrix},
\qquad \boldsymbol{x}_{k+1}=P\,\boldsymbol{x}_k.
\]
解 $P\boldsymbol{\pi}=\boldsymbol{\pi}$ 且 $\pi_1+\pi_2+\pi_3=1$，
得 $\boldsymbol{\pi}=(0.375,\ 0.375,\ 0.25)^{\mathsf T}$。
无论初始人口分布如何，长期来看 A、B 各占 $37.5\%$，C 占 $25\%$。
\textbf{这就是“平衡态”的代数含义。}
\end{exmp}

\begin{remk}{谷歌的 PageRank 是一条巨大的马尔可夫链}{pagerank}
把网页看作状态、超链接看作转移，则“网页重要性”就是随机矩阵的
平稳分布——也就是 $\lambda=1$ 的特征向量。
一个 $10^{10}$ 阶的特征值问题，催生了本世纪最有影响的线性代数应用。
\end{remk}

\begin{remk}{$\triangleright$ 上机：NumPy 求特征值}{eig-code}
{\small\ttfamily
import numpy as np\\
w, v = np.linalg.eig(A)\ \ \# w: eigenvalues, v: eigenvectors\\
np.allclose(A @ v, v * w)}
\par\smallskip
\noindent 对称矩阵请改用 \texttt{np.linalg.eigh}：更快，且保证返回实特征值与正交特征向量。
\end{remk}

\begin{sidebar}{希尔伯特、庞加莱与“谱”这个词}
“spectrum”（谱）原是物理学词汇：一束光经棱镜分出的颜色系列。
把一个算子的全体特征值称为它的谱，源于希尔伯特在积分方程方面的工作
（1904—1910）。这一命名极为贴切：对自伴算子而言，
特征值确实像光谱一样，是对象的“指纹”——
知道了谱，就知道了算子在各特征方向上的全部伸缩行为。
庞加莱在天体力学中、柯尔莫哥洛夫在概率论中，
都各自独立地把“谱”推向了应用的核心。
\end{sidebar}

\begin{chapkeys}
\small
\begin{multicols}{2}
\begin{itemize}[leftmargin=1.5em, itemsep=2pt]
\item $A\boldsymbol{v}=\lambda\boldsymbol{v}$：方向不变，只被伸缩。
\item 特征值由 $\det(\lambda I-A)=0$ 确定。
\item $\sum\lambda_i=\tr A$，$\prod\lambda_i=\det A$。
\item 几何重数 $\le$ 代数重数。
\item 可对角化 $\Leftrightarrow$ 有 $n$ 个无关特征向量 $\Leftrightarrow$ 重数处处相等。
\item $A^k=PD^kP^{-1}$，这是算矩阵幂的捷径。
\item 实对称矩阵必可正交对角化，特征向量自动正交。
\item 迭代终态由主特征值决定；马尔可夫链的平稳分布是 $\lambda=1$ 的特征向量。
\end{itemize}
\end{multicols}
\end{chapkeys}

\exhead
\begin{multicols}{2}
\small
\textbf{A 组（基础计算）}
\begin{enumerate}[leftmargin=1.6em, itemsep=3pt]
\item 求下列矩阵的特征值与全部特征向量：
      \[ \begin{bmatrix}4&1\\2&3\end{bmatrix},\qquad
         \begin{bmatrix}1&-1\\1&3\end{bmatrix},\qquad
         \begin{bmatrix}0&0&1\\0&1&0\\1&0&0\end{bmatrix}. \]
\item 判断上题中哪些矩阵可对角化，能对角化的写出 $P,D$。
\item 设 $A=\begin{bmatrix}1&2\\2&1\end{bmatrix}$，用对角化计算 $A^{10}$。
\item 已知 $A$ 为 $3$ 阶方阵，$\tr A=5$，$\det A=-6$，且有一个特征值为 $3$，
      求另外两个特征值。
\end{enumerate}
\textbf{B 组（推理证明）}
\begin{enumerate}[leftmargin=1.6em, itemsep=3pt, start=5]
\item 证明：$A$ 与 $A^{\mathsf T}$ 有相同的特征值。
\item 证明：若 $A^2=A$（幂等矩阵），则 $A$ 的特征值只能是 $0$ 或 $1$，
      且 $A$ 必可对角化。
\item 证明：属于不同特征值的特征向量线性无关（用归纳法）。
\item 设 $A$ 为实对称矩阵。证明：$A$ 正定 $\Leftrightarrow$ 全部特征值大于零。
\end{enumerate}
\textbf{C 组（应用与建模）}
\begin{enumerate}[leftmargin=1.6em, itemsep=3pt, start=9]
\item 用幂法编程求
      $A=\begin{bmatrix}4&1&0\\1&3&1\\0&1&2\end{bmatrix}$ 的主特征值与主特征向量，
      迭代 20 步，并与 \texttt{numpy.linalg.eig} 的结果比较。
\item 构造一个 $4$ 状态的马尔可夫链（如“晴天—多云—小雨—大雨”），
      求其平稳分布，并说明初值对终态的影响。
\end{enumerate}
\end{multicols}

% ===================== 第 7 章 内积、正交性与最小二乘 =====================
\chapter{内积、正交性与最小二乘}
\label{ch:ortho}

\begin{chapoutline}
\small
\begin{multicols}{2}
\textbf{学习目标}
\begin{itemize}[leftmargin=1.5em, itemsep=1.5pt]
\item 掌握内积、范数、正交的概念与几何意义；
\item 熟练使用格拉姆—施密特正交化并理解 QR 分解；
\item 掌握正交投影公式，会推导并使用最小二乘法；
\item 理解“正交 $=$ 计算最省力”这一数值观点。
\end{itemize}
\columnbreak
\textbf{关键词}
\par\smallskip
内积\quad 范数\quad 正交基\quad 施密特正交化\\
正交投影\quad 最小二乘\quad QR 分解\quad 法方程
\par\smallskip
\textbf{建议学时}
\par\smallskip
12 学时（含上机 1 学时）
\end{multicols}
\end{chapoutline}

\section{内积与范数}

\begin{defn}{内积与范数}{inner}
对 $\boldsymbol{u},\boldsymbol{v}\in\Rnum{n}$，定义
\[
\ip{\boldsymbol{u}}{\boldsymbol{v}}
=\boldsymbol{u}^{\mathsf T}\boldsymbol{v}
=u_1v_1+u_2v_2+\cdots+u_nv_n,
\qquad
\norm{\boldsymbol{v}}=\sqrt{\ip{\boldsymbol{v}}{\boldsymbol{v}}}.
\]
$\norm{\boldsymbol{v}}$ 称为 $\boldsymbol{v}$ 的\textbf{长度}（欧氏范数）。
若 $\ip{\boldsymbol{u}}{\boldsymbol{v}}=0$，则称 $\boldsymbol{u}$ 与 $\boldsymbol{v}$
\textbf{正交}，记作 $\boldsymbol{u}\perp\boldsymbol{v}$。
\end{defn}

\begin{prop}{柯西—施瓦茨不等式与三角不等式}{cauchy}
\[
\abs{\ip{\boldsymbol{u}}{\boldsymbol{v}}}\le\norm{\boldsymbol{u}}\,\norm{\boldsymbol{v}},
\qquad
\norm{\boldsymbol{u}+\boldsymbol{v}}\le\norm{\boldsymbol{u}}+\norm{\boldsymbol{v}}.
\]
由此可定义两向量的夹角
$\cos\theta=\dfrac{\ip{\boldsymbol{u}}{\boldsymbol{v}}}{\norm{\boldsymbol{u}}\norm{\boldsymbol{v}}}$。
\end{prop}

\begin{prop}{勾股定理的推广}{pythagoras}
若 $\boldsymbol{u}\perp\boldsymbol{v}$，则
$\norm{\boldsymbol{u}+\boldsymbol{v}}^2=\norm{\boldsymbol{u}}^2+\norm{\boldsymbol{v}}^2$。
更一般地，对两两正交的 $\boldsymbol{v}_1,\dots,\boldsymbol{v}_k$，
\[
\norm{c_1\boldsymbol{v}_1+\cdots+c_k\boldsymbol{v}_k}^2
=c_1^2\norm{\boldsymbol{v}_1}^2+\cdots+c_k^2\norm{\boldsymbol{v}_k}^2.
\]
\end{prop}

\section{正交基与格拉姆—施密特正交化}

\begin{defn}{正交基与标准正交基}{orthobasis}
若 $\boldsymbol{v}_1,\dots,\boldsymbol{v}_k$ 两两正交且均非零，则称其为
\textbf{正交组}；若进一步均单位化，则称\textbf{标准正交组}（orthonormal）。
若它们还张成子空间 $W$，则称为 $W$ 的\textbf{正交基}（或标准正交基）。
\end{defn}

\begin{prop}{标准正交基下坐标无须解方程}{easy-coord}
设 $\{\boldsymbol{q}_1,\dots,\boldsymbol{q}_k\}$ 是 $W$ 的标准正交基，
$\boldsymbol{w}\in W$，则
\[
\boldsymbol{w}=\ip{\boldsymbol{w}}{\boldsymbol{q}_1}\boldsymbol{q}_1
+\cdots+\ip{\boldsymbol{w}}{\boldsymbol{q}_k}\boldsymbol{q}_k,
\]
即坐标就是\textbf{逐个作内积}，无须解任何线性方程组。
这就是“正交 $=$ 计算最省力”的含义。
\end{prop}

\begin{thm}{格拉姆—施密特正交化}{gram-schmidt}
设 $\boldsymbol{v}_1,\dots,\boldsymbol{v}_k$ 线性无关。令
\[
\boldsymbol{u}_1=\boldsymbol{v}_1,\qquad
\boldsymbol{u}_j=\boldsymbol{v}_j-\sum_{i=1}^{j-1}
\frac{\ip{\boldsymbol{v}_j}{\boldsymbol{u}_i}}{\ip{\boldsymbol{u}_i}{\boldsymbol{u}_i}}
\,\boldsymbol{u}_i\quad(j=2,\dots,k),
\]
则 $\boldsymbol{u}_1,\dots,\boldsymbol{u}_k$ 两两正交且与原组张成同一子空间。
再令 $\boldsymbol{q}_j=\boldsymbol{u}_j/\norm{\boldsymbol{u}_j}$ 即得标准正交基。
\end{thm}

\begin{remk}{几何图像：逐次“减去影子”}{gs-geo}
$\boldsymbol{u}_j$ 的构造就是\textbf{把 $\boldsymbol{v}_j$ 中已经属于
$\Span\{\boldsymbol{u}_1,\dots,\boldsymbol{u}_{j-1}\}$ 的分量减掉}，
只留下与已有方向都垂直的“新方向”。
数值上应使用“改进版”（每步立即归一化或做两次正交化）以减小误差累积。
\end{remk}

\begin{exmp}{三向量正交化}{gs-ex}
$\boldsymbol{v}_1=(1,1,0)^{\mathsf T}$，
$\boldsymbol{v}_2=(1,0,1)^{\mathsf T}$，$\boldsymbol{v}_3=(-1,0,0)^{\mathsf T}$。
\[
\boldsymbol{u}_1=(1,1,0)^{\mathsf T},
\]
\[
\boldsymbol{u}_2=\boldsymbol{v}_2-\frac{\ip{\boldsymbol{v}_2}{\boldsymbol{u}_1}}
{\norm{\boldsymbol{u}_1}^{2}}\boldsymbol{u}_1
=(1,0,1)^{\mathsf T}-\frac{1}{2}(1,1,0)^{\mathsf T}
=\Big(\tfrac12,-\tfrac12,1\Big)^{\mathsf T},
\]
\[
\boldsymbol{u}_3=\boldsymbol{v}_3
-\frac{\ip{\boldsymbol{v}_3}{\boldsymbol{u}_1}}{\norm{\boldsymbol{u}_1}^{2}}\boldsymbol{u}_1
-\frac{\ip{\boldsymbol{v}_3}{\boldsymbol{u}_2}}{\norm{\boldsymbol{u}_2}^{2}}\boldsymbol{u}_2
=(-1,0,0)^{\mathsf T}+\tfrac12\boldsymbol{u}_1+\tfrac13\boldsymbol{u}_2
=\Big(-\tfrac13,\tfrac13,\tfrac13\Big)^{\mathsf T}.
\]
三个 $\boldsymbol{u}$ 两两正交（可逐一验证内积为零），
单位化后即得 $\Rnum{3}$ 的一组标准正交基。
\end{exmp}

\begin{lafig}[>=Stealth, line cap=round]{格拉姆—施密特正交化的两步：减去沿 $\boldsymbol{u}_1$
的投影得到 $\boldsymbol{u}_2$，再减去沿前两个方向的投影得到 $\boldsymbol{u}_3$}
  % ============ 左：第二步 ============
  \begin{scope}
    \node[font=\small\sffamily, color=qing] at (2.0,2.65) {第二步：从 $\boldsymbol{v}_2$ 中“扣除” $\boldsymbol{u}_1$ 方向};
    \draw[qing!45, line width=0.7pt, ->] (-2.2,0) -- (3.4,0);
    \draw[qing!45, line width=0.7pt, ->] (0,-1.35) -- (0,2.1);
    \draw[dai, line width=1.5pt, ->] (0,0) -- (2.8,0.7)
      node[pos=0.62, below, yshift=-2pt]{\small $\boldsymbol{u}_1$};
    \draw[zhe, line width=1.5pt, ->] (0,0) -- (1.1,1.75)
      node[pos=0.62, left, xshift=-2pt]{\small $\boldsymbol{v}_2$};
    % 投影
    \coordinate (P) at (0.66,0.165);
    \draw[ink!55, dashed, line width=0.7pt] (1.1,1.75) -- (P);
    \draw[ink!70, line width=0.6pt] (P) ++(0.1,0.08) -- ++(-0.12,0.096);
    \draw[qing, line width=1.5pt, ->] (P) -- (1.1,1.75)
      node[midway, left, xshift=-2pt, inner sep=1pt]{\small $\boldsymbol{u}_2$};
    \node[font=\footnotesize, color=inkl, below right, xshift=-3pt, yshift=-1pt] at (P)
      {\small $\proj\nolimits_{W}\boldsymbol{v}_2$};
    \node[font=\footnotesize\sffamily, color=inkl] at (2.0,-1.75)
      {扣除后 $\boldsymbol{u}_2\perp\boldsymbol{u}_1$，张成空间不变};
  \end{scope}
  % ============ 右：直角化过程 ============
  \begin{scope}[xshift=8.6cm]
    \node[font=\small\sffamily, color=qing] at (2.0,2.65) {逐步“扶正”：任意基 $\to$ 正交基};
    \draw[qing!45, line width=0.7pt, ->] (-2.2,0) -- (3.4,0);
    \draw[qing!45, line width=0.7pt, ->] (0,-1.35) -- (0,2.1);
    % 原始斜向量
    \draw[ink!40, line width=1.1pt, ->] (0,0) -- (1.9,1.15)
      node[pos=0.75, above, xshift=2pt, inner sep=1pt]{\small $\boldsymbol{v}_1$};
    \draw[ink!40, line width=1.1pt, ->] (0,0) -- (0.75,1.85)
      node[pos=0.75, left, xshift=-1pt, inner sep=1pt]{\small $\boldsymbol{v}_2$};
    \draw[ink!40, line width=1.1pt, ->] (0,0) -- (2.15,-0.75)
      node[pos=0.7, below, xshift=2pt, inner sep=1pt]{\small $\boldsymbol{v}_3$};
    % 正交化后
    \draw[dai, line width=1.6pt, ->] (0,0) -- (2.2,0.55)
      node[pos=0.7, below right, xshift=-2pt, inner sep=1pt]{\small $\boldsymbol{u}_1$};
    \draw[zhe, line width=1.6pt, ->] (0,0) -- (-0.5,1.9)
      node[pos=0.72, left, xshift=-1pt, inner sep=1pt]{\small $\boldsymbol{u}_2$};
    \draw[qing, line width=1.6pt, ->] (0,0) -- (-0.85,-0.55)
      node[pos=0.75, below left, xshift=2pt, inner sep=1pt]{\small $\boldsymbol{u}_3$};
    % 直角记号
    \draw[ink!55, line width=0.55pt] (0.34,0.085) -- (0.30,0.22) -- (-0.04,0.135);
    \node[font=\footnotesize\sffamily, color=inkl] at (2.0,-1.75)
      {灰色为原始向量，彩色为正交化结果};
  \end{scope}
\end{lafig}

\begin{keypoint}
正交化的本质是\textbf{“逐步扣除旧方向上的分量”}：每得到一个新方向，
就把前面所有方向上的投影一并减掉，剩下的部分自然与它们全部垂直。
扣除的只是“表示方式”，$\Span\{\boldsymbol{u}_1,\dots,\boldsymbol{u}_k\}
=\Span\{\boldsymbol{v}_1,\dots,\boldsymbol{v}_k\}$ 始终不变——
所以正交化是\textbf{换基}，不是换空间。
数值上应使用其改进版本（修正的格拉姆—施密特或 Householder），见附录 A。
\end{keypoint}

\section{正交投影}

\begin{thm}{正交分解定理}{proj-thm}
设 $W$ 是 $\Rnum{n}$ 的子空间，$\boldsymbol{b}\in\Rnum{n}$。则存在唯一的
$\widehat{\boldsymbol{b}}\in W$ 与 $\boldsymbol{z}\in W^{\perp}$ 使
$\boldsymbol{b}=\widehat{\boldsymbol{b}}+\boldsymbol{z}$。
$\widehat{\boldsymbol{b}}$ 称为 $\boldsymbol{b}$ 在 $W$ 上的\textbf{正交投影}。
\end{thm}

\begin{prop}{投影公式}{proj-formula}
若 $\{\boldsymbol{q}_1,\dots,\boldsymbol{q}_k\}$ 是 $W$ 的标准正交基，则
\[
\widehat{\boldsymbol{b}}
=\ip{\boldsymbol{b}}{\boldsymbol{q}_1}\boldsymbol{q}_1+\cdots+
  \ip{\boldsymbol{b}}{\boldsymbol{q}_k}\boldsymbol{q}_k.
\]
若 $A$ 的列构成 $W$ 的一组（不必正交的）基，则
\[
\widehat{\boldsymbol{b}}=A(A^{\mathsf T}A)^{-1}A^{\mathsf T}\boldsymbol{b}.
\]
\end{prop}

\begin{lafig}[>=Stealth, line cap=round]{正交投影：$\boldsymbol{b}$ 在子空间 $W$ 上的“影子” $\widehat{\boldsymbol{b}}$，
误差向量 $\boldsymbol{b}-\widehat{\boldsymbol{b}}$ 垂直于 $W$}
  \fill[qing!14] (-2.6,-1.64) -- (2.6,1.64) -- (2.6,1.46) -- (-2.6,-1.46) -- cycle;
  \draw[qing!55, line width=0.9pt] (-2.75,-1.64) -- (2.75,1.64);
  \node[qing, font=\small\sffamily] at (-1.85,1.35) {子空间 $W$};
  \draw[->, qing!55, line width=0.7pt] (-2.6,0) -- (2.7,0);
  \draw[->, qing!55, line width=0.7pt] (0,-1.55) -- (0,1.6);
  % 投影
  \coordinate (B) at (1.7,2.2);
  \coordinate (H) at (1.7,1.015);
  \draw[zhe, line width=1.9pt, ->] (0,0) -- (B)
     node[right, xshift=2pt]{\small $\boldsymbol{b}$};
  \draw[dai, line width=1.9pt, ->] (0,0) -- (H)
     node[midway, below right, xshift=-2pt, yshift=1pt]{\small $\widehat{\boldsymbol{b}}$};
  \draw[ink!55, dashed, line width=1.1pt] (B) -- (H);
  \node[font=\footnotesize, color=inkl] at (2.55,1.65) {$\boldsymbol{b}-\widehat{\boldsymbol{b}}$};
  \draw[dai, line width=0.9pt] (1.7,1.115) -- (1.8,1.095) -- (1.8,0.995);
  % 误差垂直记号
  \draw[->, ink!45, line width=0.9pt] (-0.9,-1.2) -- (-0.9,-0.7);
  \node[font=\footnotesize, color=inkl] at (-0.85,-1.45)
    {$\boldsymbol{b}-\widehat{\boldsymbol{b}}\ \perp\ W$};
\end{lafig}

\begin{keypoint}
\textbf{投影的最优性：}在 $W$ 的所有向量中，
$\widehat{\boldsymbol{b}}$ 是离 $\boldsymbol{b}$ 最近的那个，
且 $\norm{\boldsymbol{b}-\widehat{\boldsymbol{b}}}
\le\norm{\boldsymbol{b}-\boldsymbol{w}}$ 对一切 $\boldsymbol{w}\in W$ 成立。
这一句话就是整个最小二乘法的理论基础。
\end{keypoint}

\section{最小二乘法}

实际数据往往“没有解”。设 $A\in\R{m\times n}$（$m>n$），
方程组 $A\boldsymbol{x}=\boldsymbol{b}$ 通常无解——
$\boldsymbol{b}$ 不在列空间里。我们退而求其次：\textbf{找使残差最小的 $\boldsymbol{x}$}。

\begin{defn}{最小二乘问题}{ls-def}
求 $\widehat{\boldsymbol{x}}$ 使
\[
\norm{A\boldsymbol{x}-\boldsymbol{b}}\ \text{达到最小}.
\]
\end{defn}

\begin{thm}{法方程}{normal-eq}
$\widehat{\boldsymbol{x}}$ 是最小二乘解
$\Longleftrightarrow A\boldsymbol{x}=\boldsymbol{b}$ 的残差
$\boldsymbol{b}-A\widehat{\boldsymbol{x}}$ 垂直于 $\Col(A)$
$\Longleftrightarrow$
\[
A^{\mathsf T}A\,\widehat{\boldsymbol{x}}=A^{\mathsf T}\boldsymbol{b}.
\]
当 $A$ 各列线性无关时 $A^{\mathsf T}A$ 可逆，解唯一：
$\widehat{\boldsymbol{x}}=(A^{\mathsf T}A)^{-1}A^{\mathsf T}\boldsymbol{b}$。
\end{thm}

\begin{pf}
残差 $\boldsymbol{r}=\boldsymbol{b}-A\widehat{\boldsymbol{x}}$ 垂直于 $\Col(A)$
等价于 $\boldsymbol{r}\in\Nul(A^{\mathsf T})$，
即 $A^{\mathsf T}(\boldsymbol{b}-A\widehat{\boldsymbol{x}})=\boldsymbol{0}$，
整理即得法方程。
\end{pf}

\begin{exmp}{直线拟合}{ls-fit}
测得四组数据 $(-1,0),(0,1),(1,2),(2,4)$，拟合 $y=mx+c$。
设计矩阵与观测向量：
\[
A=\begin{bmatrix}-1&1\\0&1\\1&1\\2&1\end{bmatrix},\qquad
\boldsymbol{b}=\begin{bmatrix}0\\1\\2\\4\end{bmatrix}.
\]
计算 $A^{\mathsf T}A=\begin{bmatrix}6&2\\2&4\end{bmatrix}$，
$A^{\mathsf T}\boldsymbol{b}=\begin{bmatrix}10\\7\end{bmatrix}$。
解法方程
$\begin{bmatrix}6&2\\2&4\end{bmatrix}\widehat{\boldsymbol{x}}
=\begin{bmatrix}10\\7\end{bmatrix}$
得 $\widehat{\boldsymbol{x}}=(6/5,\,23/20)^{\mathsf T}$，即
\[
\widehat{y}=1.2\,x+1.15.
\]
残差平方和为 $\norm{A\widehat{\boldsymbol{x}}-\boldsymbol{b}}^2=0.45$。
\begin{lafig}[>=Stealth, line cap=round]{最小二乘直线拟合：不要求穿过每个点，只要求残差平方和最小}
  \begin{axis}[
      width=0.86\textwidth, height=4.9cm,
      axis lines=left, axis line style={qing!55, line width=0.7pt},
      xmin=-1.6, xmax=2.7, ymin=-0.5, ymax=4.9,
      xtick={-1,0,1,2}, ytick={0,1,2,3,4},
      tick label style={font=\footnotesize, color=inkl},
      label style={font=\small, color=ink},
      xlabel={$x$}, ylabel={$y$},
      grid=major, grid style={line!70, very thin},
    ]
    \addplot[only marks, mark=*, mark size=1.9pt, color=dai!85] coordinates
      {(-1,0) (0,1) (1,2) (2,4)};
    \addplot[zhe, domain=-1.5:2.5, samples=2, line width=1.5pt] {1.2*x+1.15};
    \node[font=\footnotesize, color=zhe, anchor=west] at (axis cs:-1.35,0.55) {$\widehat{y}=1.2x+1.15$};
  \end{axis}
\end{lafig}
\end{exmp}

\begin{remk}{为什么叫“最小二乘”}{why-ls}
因为目标函数是\textbf{残差的平方}之和。
用平方而非绝对值，好处是可导、有唯一解、且与正交投影的几何天然契合；
代价是对离群点（outlier）非常敏感。
实践中常用加权最小二乘或鲁棒回归来缓解。
\end{remk}

\section{QR 分解}

\begin{thm}{QR 分解}{qr}
设 $A\in\R{m\times n}$（$m\ge n$）各列线性无关。则存在
$Q\in\R{m\times n}$（各列为标准正交组）与上三角矩阵 $R\in\R{n\times n}$ 使
\[
A=QR.
\]
对 $A$ 的各列做格拉姆—施密特正交化即得 $Q$；
$R$ 的元素恰为正交化过程中减去的系数与各 $\boldsymbol{u}_j$ 的长度。
\end{thm}

\begin{prop}{用 QR 解最小二乘}{qr-ls}
把 $A=QR$ 代入法方程：
\[
A^{\mathsf T}A\widehat{\boldsymbol{x}}=A^{\mathsf T}\boldsymbol{b}
\ \Longrightarrow\
R\widehat{\boldsymbol{x}}=Q^{\mathsf T}\boldsymbol{b},
\]
只需一次回代，\textbf{无须构造 $A^{\mathsf T}A$，也无须求逆}。
这是数值线性代数推荐的标准做法：
$A^{\mathsf T}A$ 的条件数是 $A$ 的平方，直接解法方程会放大误差。
\end{prop}

\begin{remk}{$\triangleright$ 上机：三种解法对比}{ls-code}
{\small\ttfamily
import numpy as np\\
x1 = np.linalg.lstsq(A, b, rcond=None)[0]\\
Q, R = np.linalg.qr(A)\ \ ; x2 = np.linalg.solve(R, Q.T @ b)\\
x3 = np.linalg.solve(A.T @ A, A.T @ b)}
\par\smallskip
\noindent 三者在理想算术下相同；把 $A$ 的元素乘以 $10^{-6}$ 后再比较，
会发现 \texttt{x3} 的误差明显更大。
\end{remk}

\section{应用：傅里叶与正交多项式}

把函数看作“无穷维向量”，正交性的威力立刻放大。取区间 $[-\pi,\pi]$，
定义 $\ip{f}{g}=\int_{-\pi}^{\pi}f g\,\mathrm{d}x$，则
$\{1,\cos x,\sin x,\cos 2x,\sin 2x,\dots\}$ 是一组正交基，于是
\[
f(x)\approx\frac{a_0}{2}+\sum_{k\ge1}\bigl(a_k\cos kx+b_k\sin kx\bigr),
\qquad
a_k=\frac{1}{\pi}\int_{-\pi}^{\pi}f\cos kx\,\mathrm{d}x,
\]
系数\textbf{逐个由内积给出}——正是“标准正交基下坐标无须解方程”的无限维版本。
有限维的对应物是离散余弦变换（DCT），它是 JPEG 压缩的核心。

\begin{sidebar}{勒让德、傅里叶与“函数也是向量”}
把 $\cos kx$、$\sin kx$ 当作“坐标轴”，把函数当作“向量”——
这套观点在 1807 年\textbf{傅里叶}提出热传导理论时曾遭到严格质疑
（拉格朗日就认为任意函数不能这样展开）。
真正的澄清要等到二十世纪的希尔伯特空间理论。
有趣的是，\textbf{勒让德}多项式、切比雪夫多项式、厄米特多项式……
这一整族“正交多项式”，本质上都是在不同内积下做
格拉姆—施密特正交化的产物。第七章的方法，在无穷维依然成立。
\end{sidebar}

\begin{chapkeys}
\small
\begin{multicols}{2}
\begin{itemize}[leftmargin=1.5em, itemsep=2pt]
\item $\ip{\boldsymbol{u}}{\boldsymbol{v}}=\boldsymbol{u}^{\mathsf T}\boldsymbol{v}$，
      正交即内积为零。
\item 标准正交基下坐标 $=$ 逐个内积，无须解方程。
\item 施密特正交化：逐次减去已有方向上的“影子”。
\item 投影：$\widehat{\boldsymbol{b}}=A(A^{\mathsf T}A)^{-1}A^{\mathsf T}\boldsymbol{b}$。
\item 最小二乘 $\Leftrightarrow$ 法方程 $A^{\mathsf T}A\widehat{\boldsymbol{x}}=A^{\mathsf T}\boldsymbol{b}$。
\item QR 分解后只需 $R\widehat{\boldsymbol{x}}=Q^{\mathsf T}\boldsymbol{b}$。
\item 数值上优先用 QR，避免 $A^{\mathsf T}A$ 的条件数平方效应。
\item 正交变换保长度、保角度：$Q^{\mathsf T}Q=I$。
\end{itemize}
\end{multicols}
\end{chapkeys}

\exhead
\begin{multicols}{2}
\small
\textbf{A 组（基础计算）}
\begin{enumerate}[leftmargin=1.6em, itemsep=3pt]
\item 求 $\boldsymbol{u}=(1,2,2)^{\mathsf T}$ 与
      $\boldsymbol{v}=(-2,1,0)^{\mathsf T}$ 的内积、夹角，并验证勾股定理。
\item 用格拉姆—施密特把 $(1,1,1)^{\mathsf T},(1,0,1)^{\mathsf T}$ 正交化并单位化。
\item 求 $\boldsymbol{b}=(1,2,3)^{\mathsf T}$ 在子空间
      $W=\Span\{(1,1,0)^{\mathsf T},(0,1,1)^{\mathsf T}\}$
      上的正交投影。
\item 用最小二乘法拟合下列数据点：
      $(-1,4),(0,2),(1,1),(2,1)$，求所得直线。
\end{enumerate}
\textbf{B 组（推理证明）}
\begin{enumerate}[leftmargin=1.6em, itemsep=3pt, start=5]
\item 证明：若 $Q$ 为正交矩阵，则 $\norm{Q\boldsymbol{x}}=\norm{\boldsymbol{x}}$，
      且 $(Q\boldsymbol{x})\cdot(Q\boldsymbol{y})=\boldsymbol{x}\cdot\boldsymbol{y}$。
\item 证明：$(A^{\mathsf T}A)$ 对称，且当 $A$ 各列无关时正定。
\item 证明：最小二乘解必满足 $\norm{\boldsymbol{b}}^2
      =\norm{\widehat{\boldsymbol{b}}}^2+\norm{\boldsymbol{b}-\widehat{\boldsymbol{b}}}^2$。
\item 设 $A=QR$。证明：$\rank A=\rank R$，且 $\abs{\det(A^{\mathsf T}A)}
      =(\det R)^2$。
\end{enumerate}
\textbf{C 组（应用与建模）}
\begin{enumerate}[leftmargin=1.6em, itemsep=3pt, start=9]
\item 用二次多项式 $y=ax^2+bx+c$ 拟合上题第 4 小题的数据，
      比较残差平方和，并讨论是否“次数越高越好”。
\item 实现（或调用）Gram—Schmidt，对一个条件数较大的矩阵做 QR，
      观察 $Q^{\mathsf T}Q-I$ 的范数；再做“两次正交化”并比较。
\end{enumerate}
\end{multicols}

% ===================== 第 8 章 二次型与正定矩阵 =====================
\chapter{二次型与正定矩阵}
\label{ch:quad}

\begin{chapoutline}
\small
\begin{multicols}{2}
\textbf{学习目标}
\begin{itemize}[leftmargin=1.5em, itemsep=1.5pt]
\item 把二次齐次式写成矩阵形式 $\boldsymbol{x}^{\mathsf T}A\boldsymbol{x}$；
\item 掌握配方法与正交变换法化标准形；
\item 理解惯性定理与合同变换，熟练判定正定性；
\item 从几何上理解二次曲线、二次曲面与极值问题。
\end{itemize}
\columnbreak
\textbf{关键词}
\par\smallskip
二次型\quad 标准形\quad 惯性定理\quad 合同\\
正定\quad 顺序主子式\quad 二次曲面\quad 瑞利商
\par\smallskip
\textbf{建议学时}
\par\smallskip
10 学时
\end{multicols}
\end{chapoutline}

\section{二次型及其矩阵表示}

\begin{defn}{二次型}{quad-def}
含 $n$ 个变量的\textbf{二次齐次多项式}
\[
f(x_1,\dots,x_n)=\sum_{i=1}^{n}\sum_{j=1}^{n}a_{ij}x_ix_j
=\boldsymbol{x}^{\mathsf T}A\boldsymbol{x},
\qquad A^{\mathsf T}=A,
\]
称为 $n$ 元\textbf{二次型}，$A$ 称为它的\textbf{矩阵}。
约定 $A$ 取对称矩阵（$a_{ij}=a_{ji}$），此时表示唯一。
\end{defn}

\begin{exmp}{写出二次型的矩阵}{quad-ex}
\[
f=2x_1^2+4x_1x_2-6x_2^2
=\begin{bmatrix}x_1&x_2\end{bmatrix}
\begin{bmatrix}2&2\\2&-6\end{bmatrix}
\begin{bmatrix}x_1\\x_2\end{bmatrix}.
\]
\textbf{易错点}：交叉项 $x_1x_2$ 的系数要\emph{对半分}到 $a_{12}$ 与 $a_{21}$。
\end{exmp}

\section{标准形与配方法}

\begin{defn}{标准形}{canonical}
只含平方项的二次型
$f=d_1y_1^2+d_2y_2^2+\cdots+d_ny_n^2$
称为\textbf{标准形}，其矩阵为对角矩阵 $\diag(d_1,\dots,d_n)$。
\end{defn}

\begin{thm}{配方法（拉格朗日法）}{lagrange}
任何实二次型都可用可逆线性替换 $\boldsymbol{x}=C\boldsymbol{y}$ 化为标准形。
要点：
\begin{enumerate}[label=(\arabic*), itemsep=1pt]
\item 若含 $x_i^2$ 项，先把含 $x_i$ 的项全部配成完全平方；
\item 若不含任何平方项但有 $x_ix_j$，先作替换
      $x_i=y_i+y_j$，$x_j=y_i-y_j$ 制造平方项；
\item 重复直至只剩平方项。
\end{enumerate}
\end{thm}

\begin{exmp}{配方法全过程}{complete-square}
\[
f=x_1^2+2x_2^2+2x_1x_2+2x_1x_3
=x_1^2+2x_1(x_2+x_3)+2x_2^2
\]
\[
=(x_1+x_2+x_3)^2+x_2^2-2x_2x_3-x_3^2
=(x_1+x_2+x_3)^2+(x_2-x_3)^2-2x_3^2.
\]
令 $y_1=x_1+x_2+x_3$，$y_2=x_2-x_3$，$y_3=x_3$，
则 $f=y_1^2+y_2^2-2y_3^2$，系数为 $1,1,-2$。
\end{exmp}

\begin{thm}{正交变换法}{orth-quad}
对实对称矩阵 $A$，由谱定理存在正交矩阵 $Q$ 使 $Q^{\mathsf T}AQ=\diag(\lambda_1,\dots,\lambda_n)$。
作替换 $\boldsymbol{x}=Q\boldsymbol{y}$，则
\[
\boldsymbol{x}^{\mathsf T}A\boldsymbol{x}
=\boldsymbol{y}^{\mathsf T}\diag(\lambda_i)\boldsymbol{y}
=\sum_{i=1}^{n}\lambda_i y_i^2.
\]
\textbf{标准形的系数恰为特征值}，且正交变换不改变几何形状
（它只是转轴）。
\end{thm}

\begin{remk}{配方与转轴的区别}{two-ways}
配方法得到的标准形系数一般\textbf{不是}特征值，
对应的线性替换也非正交，因此\textbf{会改变几何形状}；
正交变换法保持形状，只是把坐标轴转到特征方向。
讨论\emph{代数}问题时用配方（快），讨论\emph{几何}问题时用正交变换（真）。
\end{remk}

\section{惯性定理与合同}

\begin{thm}{惯性定理}{sylvester}
设实二次型 $f$ 经两个可逆线性替换分别化为标准形，则
\textbf{正系数的个数}与\textbf{负系数的个数}分别相同。
正、负、零系数的个数分别记为 $p,\,q,\,n-p-q$，
称 $(p,q)$ 为 $f$ 的\textbf{符号差}（正惯性指数、负惯性指数）。
\end{thm}

\begin{defn}{合同}{congruent}
若存在可逆矩阵 $C$ 使 $B=C^{\mathsf T}AC$，则称 $A$ 与 $B$ \textbf{合同}。
合同保持对称性与惯性指数，但不保持特征值。
\end{defn}

\begin{keypoint}
\textbf{三者的关系要分清}：
相似 $B=P^{-1}AP$ 保特征值；合同 $B=C^{\mathsf T}AC$ 保惯性指数；
当 $C$ 为正交矩阵时二者重合（$C^{-1}=C^{\mathsf T}$）。
对实对称矩阵，“相似 $\Leftrightarrow$ 合同 $\Leftrightarrow$ 特征值相同”。
\end{keypoint}

\section{正定及其判定}

\begin{defn}{正定与半正定}{pd-def}
设 $A$ 对称。若对任意 $\boldsymbol{x}\neq\boldsymbol{0}$ 都有
$\boldsymbol{x}^{\mathsf T}A\boldsymbol{x}>0$，则称 $A$ \textbf{正定}；
若 $\ge 0$ 则称\textbf{半正定}；若 $-\!A$ 正定则称 $A$ \textbf{负定}；
若两种符号都出现则称\textbf{不定}。
\end{defn}

\begin{thm}{正定的等价条件}{pd-thm}
设 $A$ 为 $n$ 阶实对称矩阵，则下列条件等价：
\begin{enumerate}[label=(\alph*), itemsep=1pt]
\item $A$ 正定；
\item 全部特征值 $\lambda_i>0$；
\item 全部\textbf{顺序主子式} $\Delta_k=\det A_k>0$（$k=1,\dots,n$）；
\item 存在可逆矩阵 $C$ 使 $A=C^{\mathsf T}C$；
\item 正惯性指数 $p=n$；
\item 存在列满秩矩阵 $B$ 使 $A=B^{\mathsf T}B$。
\end{enumerate}
\end{thm}

\begin{keypoint}
\textbf{手算首选顺序主子式判别法}（Sylvester 准则），
\textbf{理论分析首选特征值}。
注意顺序主子式全为正只能判定正定，
\emph{不能}用来判定负定；负定要求 $(-1)^k\Delta_k>0$ 对所有 $k$ 成立。
\end{keypoint}

\begin{exmp}{判定正定性}{pd-ex}
判断 $A=\begin{bmatrix}2&1&1\\1&3&0\\1&0&1\end{bmatrix}$ 是否正定。
\[
\Delta_1=2>0,\qquad
\Delta_2=\begin{vmatrix}2&1\\1&3\end{vmatrix}=5>0,\qquad
\Delta_3=\det A=2(3)-1(1)+1(-3)=2>0.
\]
三个顺序主子式均为正，故 $A$ 正定。
（验证：$\tr A=6$，$\det A=2$，特征值均为正。）
\end{exmp}

\section{几何：二次曲线与二次曲面}

\begin{exmp}{二次曲线的分类}{conic}
考察 $\boldsymbol{x}^{\mathsf T}A\boldsymbol{x}=1$，其中
$A=\diag(\lambda_1,\lambda_2)$：
\begin{itemize}[itemsep=1pt]
\item $\lambda_1,\lambda_2>0$：椭圆，半轴长 $1/\sqrt{\lambda_1},\,1/\sqrt{\lambda_2}$；
\item 一正一负：双曲线；
\item $\lambda_2=0$：两条平行直线；
\item $\lambda_1,\lambda_2<0$：无实轨迹。
\end{itemize}
\begin{lafig}[>=Stealth, line cap=round]{$\boldsymbol{x}^{\mathsf T}A\boldsymbol{x}=1$ 的几何形状由特征值的符号决定}
  \begin{scope}
    \draw[qing!55, line width=0.7pt, ->] (-2.0,0) -- (2.1,0);
    \draw[qing!55, line width=0.7pt, ->] (0,-1.7) -- (0,1.85);
    \draw[qing, line width=1.7pt] (0,0) ellipse (1.6 and 1.0);
    \draw[dai, line width=1.4pt] (-1.6,0) -- (1.6,0);
    \draw[dai, line width=1.4pt] (0,-1.0) -- (0,1.0);
    \node[dai, above=1pt] at (0.8,0) {\scriptsize $1/\sqrt{\lambda_1}$};
    \node[dai, left=2pt] at (0,0.5) {\scriptsize $1/\sqrt{\lambda_2}$};
    \node[font=\footnotesize\sffamily, color=inkl] at (0,-2.1)
      {正定：椭圆（$\lambda_1,\lambda_2>0$）};
  \end{scope}
  \begin{scope}[xshift=5.2cm]
    \draw[qing!55, line width=0.7pt, ->] (-2.0,0) -- (2.1,0);
    \draw[qing!55, line width=0.7pt, ->] (0,-1.7) -- (0,1.85);
    \draw[zhe, line width=1.7pt, domain=0.46:2.05, smooth, variable=\x]
      plot ({\x},{0.62*sqrt(\x*\x-0.18)});
    \draw[zhe, line width=1.7pt, domain=0.46:2.05, smooth, variable=\x]
      plot ({\x},{-0.62*sqrt(\x*\x-0.18)});
    \draw[zhe, line width=1.7pt, domain=-2.05:-0.46, smooth, variable=\x]
      plot ({\x},{0.62*sqrt(\x*\x-0.18)});
    \draw[zhe, line width=1.7pt, domain=-2.05:-0.46, smooth, variable=\x]
      plot ({\x},{-0.62*sqrt(\x*\x-0.18)});
    \draw[dai, line width=1.1pt] (-1.95,-1.35) -- (1.95,1.35);
    \draw[dai, line width=1.1pt] (-1.95,1.35) -- (1.95,-1.35);
    \node[font=\footnotesize\sffamily, color=inkl] at (0,-2.1)
      {不定：双曲线（一正一负）};
  \end{scope}
  \begin{scope}[xshift=10.4cm]
    \draw[qing!55, line width=0.7pt, ->] (-2.0,0) -- (2.1,0);
    \draw[qing!55, line width=0.7pt, ->] (0,-1.7) -- (0,1.85);
    \draw[dai, line width=1.7pt] (0,0) ellipse (1.55 and 1.0);
    \draw[dai!55, line width=1.7pt] (0,0) ellipse (0.95 and 0.62);
    \draw[dai!30, line width=1.7pt] (0,0) ellipse (0.45 and 0.29);
    \foreach \i in {1.55,0.95,0.45}{
      \fill[dai] (\i,0) circle (1.7pt);
    }
    \node[font=\footnotesize\sffamily, color=inkl] at (0,-2.1)
      {半正定：退化为一条线（$\lambda_2=0$）};
  \end{scope}
\end{lafig}
\end{exmp}

\begin{remk}{为什么椭圆的半轴是 $1/\sqrt{\lambda_i}$}{ellipse-axis}
在特征方向 $\boldsymbol{p}_i$ 上取 $\boldsymbol{x}=t\boldsymbol{p}_i$，
则 $\boldsymbol{x}^{\mathsf T}A\boldsymbol{x}=t^2\lambda_i=1$，
故 $t=1/\sqrt{\lambda_i}$。
\textbf{特征值越大，等高线在该方向上越“窄”。}
这正是“大特征值对应大能量方向”的几何来源。
\end{remk}

\begin{lafig}[>=Stealth, line cap=round]{三维情形：二次型 $f(\boldsymbol{x})=\boldsymbol{x}^{\mathsf T}A\boldsymbol{x}$
的正、负、零特征值分别造出“碗”“鞍”与“谷”}
\begin{axis}[name=ax1, width=5.1cm, height=4.6cm, view={35}{18},
  axis lines=none, ticks=none, title style={font=\small\sffamily, text=qing},
  title={正定：椭圆抛物面},
  colormap={qmap}{rgb255=(247,248,246) rgb255=(150,185,196) rgb255=(42,88,102)},
  declare function={f1(\x,\y)=0.55*(\x*\x+\y*\y);}]
  \addplot3[surf, shader=interp, domain=-1.5:1.5, y domain=-1.5:1.5,
            samples=19, z buffer=sort] {f1(x,y)};
\end{axis}
\begin{axis}[name=ax2, width=5.1cm, height=4.6cm, view={35}{18},
  at={(ax1.east)}, anchor=west, xshift=10pt,
  axis lines=none, ticks=none, title style={font=\small\sffamily, text=zhe},
  title={不定：双曲抛物面（鞍面）},
  colormap={zmap}{rgb255=(247,248,246) rgb255=(196,178,150) rgb255=(152,100,64)},
  declare function={f2(\x,\y)=0.42*(\x*\x-\y*\y);}]
  \addplot3[surf, shader=interp, domain=-1.5:1.5, y domain=-1.5:1.5,
            samples=19, z buffer=sort] {f2(x,y)};
\end{axis}
\begin{axis}[name=ax3, width=5.1cm, height=4.6cm, view={35}{18},
  at={(ax2.east)}, anchor=west, xshift=10pt,
  axis lines=none, ticks=none, title style={font=\small\sffamily, text=dai},
  title={半正定：抛物柱面},
  colormap={dmap}{rgb255=(247,248,246) rgb255=(160,190,175) rgb255=(84,110,98)},
  declare function={f3(\x,\y)=0.55*\x*\x;}]
  \addplot3[surf, shader=interp, domain=-1.5:1.5, y domain=-1.5:1.5,
            samples=19, z buffer=sort] {f3(x,y)};
\end{axis}
\end{lafig}

\begin{remk}{从曲面回到矩阵}{surface-to-matrix}
三幅图分别是正定（$A\succ 0$，两个正特征值）、
不定（特征值一正一负，原点是\textbf{鞍点}）与
半正定（有一个零特征值，曲面沿一个方向“平”掉）的标准像。
第 8.6 节的极值判据说的正是：$H$ 正定 $\Rightarrow$ 谷底（极小值），
$H$ 不定 $\Rightarrow$ 鞍点，$H$ 半定 $\Rightarrow$ 需要进一步判别。
\end{remk}

\section{应用：极值与瑞利商}

\begin{thm}{多元函数极值的二阶判据}{hessian}
设 $f$ 在 $\boldsymbol{x}_0$ 处梯度为零，Hessian 矩阵 $H=\nabla^2 f(\boldsymbol{x}_0)$
对称。则
\[
H\ \text{正定}\Rightarrow\text{极小值；}\qquad
H\ \text{负定}\Rightarrow\text{极大值；}\qquad
H\ \text{不定}\Rightarrow\text{鞍点}.
\]
\end{thm}

\begin{defn}{瑞利商}{rayleigh}
对对称矩阵 $A$ 与 $\boldsymbol{x}\neq\boldsymbol{0}$，称
\[
R(\boldsymbol{x})=\frac{\boldsymbol{x}^{\mathsf T}A\boldsymbol{x}}
                      {\boldsymbol{x}^{\mathsf T}\boldsymbol{x}}
\]
为\textbf{瑞利商}。在单位球面上，
\[
\lambda_{\min}\le R(\boldsymbol{x})\le\lambda_{\max},
\]
且最值分别在最小、最大特征向量处取得。
\end{defn}

\begin{keypoint}
瑞利商把“特征值”重新解释为\textbf{广义能量的极值}。
主成分分析（第 9 章）中“方差最大的方向就是最大特征向量”，
本质上是瑞利商最大化的结论。
\end{keypoint}

\begin{exmp}{一个鞍点的识别}{saddle-ex}
$f(x,y)=x^2-y^2$，矩阵 $A=\diag(1,-1)$，$(0,0)$ 是临界点。
$A$ 不定，故 $(0,0)$ 是鞍点：沿 $x$ 方向是极小，沿 $y$ 方向是极大。
等高线正是双曲线 $x^2-y^2=c$，$c=0$ 时退化为一对相交直线。
\end{exmp}

\begin{sidebar}{从柯西到西尔维斯特：主子式判别法}
正定的顺序主子式判别法以\textbf{西尔维斯特}（J.\ J.\ Sylvester，1814—1897）命名，
但柯西早在 1829 年就给出了二次型惯性定理的雏形。
西尔维斯特一生跌宕：因宗教原因无法取得剑桥学位，
靠保险精算与律师职业谋生，五十岁后才在约翰斯·霍普金斯大学
获得平生第一个教职，却在晚年成为英国最负盛名的数学家之一。
他发明了大量今天仍在使用的术语——矩阵（matrix）、
约化行列式（minor）、判别式（discriminant）都出自他的笔下。
\end{sidebar}

\begin{chapkeys}
\small
\begin{multicols}{2}
\begin{itemize}[leftmargin=1.5em, itemsep=2pt]
\item 二次型 $\boldsymbol{x}^{\mathsf T}A\boldsymbol{x}$，$A$ 取对称矩阵。
\item 交叉项系数要对半分到 $a_{ij}$ 与 $a_{ji}$。
\item 正交变换化标准形，系数恰为特征值。
\item 惯性定理：正、负系数个数在可逆替换下不变。
\item 合同保惯性指数，相似保特征值；正交时二者一致。
\item 正定 $\Leftrightarrow$ 特征值全正 $\Leftrightarrow$ 顺序主子式全正。
\item 椭圆半轴长 $=1/\sqrt{\lambda_i}$。
\item 瑞利商在特征方向上取到极值。
\end{itemize}
\end{multicols}
\end{chapkeys}

\exhead
\begin{multicols}{2}
\small
\textbf{A 组（基础计算）}
\begin{enumerate}[leftmargin=1.6em, itemsep=3pt]
\item 写出下列二次型的矩阵：
      $f=3x_1^2+2x_2^2-x_3^2+4x_1x_2-2x_2x_3$。
\item 用配方法把 $f=x_1^2+2x_2^2+4x_3^2+2x_1x_2+2x_1x_3$ 化为标准形。
\item 用正交变换把 $f=5x_1^2+5x_2^2+x_3^2-2x_1x_2$ 化为标准形。
\item 判断下列矩阵是否正定：
      $\begin{bmatrix}2&1\\1&2\end{bmatrix}$，
      $\begin{bmatrix}1&2\\2&1\end{bmatrix}$，
      $\begin{bmatrix}1&0\\0&0\end{bmatrix}$。
\end{enumerate}
\textbf{B 组（推理证明）}
\begin{enumerate}[leftmargin=1.6em, itemsep=3pt, start=5]
\item 证明：若 $A$ 正定，则 $A$ 可逆且 $A^{-1}$ 也正定。
\item 证明：若 $A,B$ 均为 $n$ 阶正定矩阵，则 $A+B$ 正定。
      问 $AB$ 何时正定？
\item 证明：实对称矩阵 $A$ 正定 $\Leftrightarrow$ 存在列满秩矩阵 $B$ 使 $A=B^{\mathsf T}B$。
\item 证明：二次型 $f=\boldsymbol{x}^{\mathsf T}A\boldsymbol{x}$ 在
      $\norm{\boldsymbol{x}}=1$ 上的最大值等于 $A$ 的最大特征值。
\end{enumerate}
\textbf{C 组（应用与建模）}
\begin{enumerate}[leftmargin=1.6em, itemsep=3pt, start=9]
\item 用 pgfplots 或 Python 绘制
      $f(x,y)=x^2+2y^2+2xy$ 的等高线，
      并说明该二次型是否正定、形状如何随系数变化。
\item 投资组合理论中，协方差矩阵 $\Sigma$ 的正定性有什么含义？
      若 $\Sigma$ 半正定而非正定，说明组合中存在什么关系？
\end{enumerate}
\end{multicols}

% ===================== 第 9 章 奇异值分解与应用 =====================
\chapter{奇异值分解与应用}
\label{ch:svd}

\begin{chapoutline}
\small
\begin{multicols}{2}
\textbf{学习目标}
\begin{itemize}[leftmargin=1.5em, itemsep=1.5pt]
\item 理解任意矩阵都有奇异值分解，并掌握其构造；
\item 用“单位圆变成椭圆”的几何图像理解 SVD；
\item 由 SVD 读出四个基本子空间的正交基；
\item 掌握最优低秩逼近与主成分分析（PCA）。
\end{itemize}
\columnbreak
\textbf{关键词}
\par\smallskip
奇异值\quad 奇异向量\quad 伪逆\quad 最优低秩逼近\\
Eckart--Young 定理\quad PCA\quad 数据压缩
\par\smallskip
\textbf{建议学时}
\par\smallskip
10 学时（含上机 2 学时）
\end{multicols}
\end{chapoutline}

\section{为什么要 SVD}

特征分解 $A=PDP^{-1}$ 有两个限制：只对方阵有效，
且即便有效也要求可对角化。而\textbf{奇异值分解对任意矩阵都存在}——
长方形的、奇异的、甚至全零的。它是线性代数的“万能分解”，
也是当今数据科学中使用最频繁的矩阵工具。

\begin{thm}{奇异值分解}{svd-thm}
设 $A\in\R{m\times n}$。则存在正交矩阵
$U\in\R{m\times m}$、$V\in\R{n\times n}$ 与非负数
$\sigma_1\ge\sigma_2\ge\cdots\ge\sigma_r>0$（$r=\rank A$），使
\[
A=U\,\Sigma\,V^{\mathsf T},
\qquad
\Sigma=\begin{bmatrix}\diag(\sigma_1,\dots,\sigma_r)&O\\O&O\end{bmatrix}.
\]
$\sigma_i$ 称为\textbf{奇异值}，$U,V$ 的列分别称为左、右\textbf{奇异向量}。
\end{thm}

\begin{prop}{奇异值与特征值的关系}{svd-eig}
$\sigma_i^2$ 是 $A^{\mathsf T}A$ 与 $AA^{\mathsf T}$ 的非零特征值；
$V$ 的列是 $A^{\mathsf T}A$ 的特征向量，$U$ 的列是 $AA^{\mathsf T}$ 的特征向量，
且 $A\boldsymbol{v}_i=\sigma_i\boldsymbol{u}_i$。
因此
\[
\sigma_i=\sqrt{\lambda_i(A^{\mathsf T}A)},\qquad
\sigma_1\sigma_2\cdots\sigma_r=\text{$A$ 各行（列）体积的缩放因子}.
\]
\end{prop}

\section{几何图像：圆如何变成椭圆}

SVD 最漂亮的地方在于它可以\textbf{一步一步画出来}。
把 $A$ 分解成“旋转—拉伸—旋转”：
\[
A=U\,\Sigma\,V^{\mathsf T},
\]
其中 $U$ 与 $V^{\mathsf T}$ 是旋转（或反射），$\Sigma$ 是沿坐标轴的拉伸。

\begin{lafig}[>=Stealth, line cap=round]{SVD 的四步几何：单位圆经 $V^{\mathsf T}$ 转正、被 $\Sigma$ 沿坐标轴拉伸、
再经 $U$ 转到最终方向，得到椭圆 $\sigma_1\boldsymbol{u}_1,\ \sigma_2\boldsymbol{u}_2$}
  % 1) 单位圆
  \begin{scope}
    \draw[line!80, very thin, step=0.5] (-1.7,-1.4) grid (1.7,1.7);
    \draw[qing!55, line width=0.7pt, ->] (-1.85,0) -- (1.85,0);
    \draw[qing!55, line width=0.7pt, ->] (0,-1.55) -- (0,1.85);
    \draw[qing, line width=1.7pt] (0,0) circle (1.25);
    \draw[dai, line width=1.5pt, ->] (0,0) -- (0.884,0.884);
    \draw[zhe, line width=1.5pt, ->] (0,0) -- (-0.884,0.884);
    \node[font=\footnotesize\sffamily, color=inkl] at (0,-1.85) {① 单位圆};
  \end{scope}
  % 2) V^T 旋转
  \begin{scope}[xshift=4.55cm, rotate=-28]
    \draw[line!80, very thin, step=0.5] (-1.7,-1.4) grid (1.7,1.7);
    \draw[qing, line width=1.7pt] (0,0) circle (1.25);
  \end{scope}
  \begin{scope}[xshift=4.55cm]
    \draw[qing!55, line width=0.7pt, ->] (-1.85,0) -- (1.85,0);
    \draw[qing!55, line width=0.7pt, ->] (0,-1.55) -- (0,1.85);
    \draw[dai, line width=1.5pt, ->] (0,0) -- (1.104,-0.587);
    \draw[zhe, line width=1.5pt, ->] (0,0) -- (0.587,1.104);
    \node[font=\footnotesize\sffamily, color=inkl] at (0,-1.85) {② $V^{\mathsf T}$：转正};
  \end{scope}
  % 3) Sigma 拉伸
  \begin{scope}[xshift=9.1cm]
    \draw[line!80, very thin, step=0.5] (-1.9,-1.4) grid (2.3,1.7);
    \draw[qing!55, line width=0.7pt, ->] (-2.05,0) -- (2.45,0);
    \draw[qing!55, line width=0.7pt, ->] (0,-1.55) -- (0,1.85);
    \draw[zhe, line width=1.7pt] (0,0) ellipse (1.95 and 0.72);
    \draw[dai, line width=1.5pt, ->] (0,0) -- (1.95,0);
    \draw[zhe, line width=1.5pt, ->] (0,0) -- (0,0.72);
    \node[font=\footnotesize\sffamily, color=inkl] at (0.2,-1.85) {③ $\Sigma$：拉伸};
  \end{scope}
  % 4) U 旋转
  \begin{scope}[xshift=13.65cm, rotate=17]
    \draw[line!80, very thin, step=0.5] (-1.9,-1.4) grid (2.3,1.7);
    \draw[qing!55, line width=0.7pt, ->] (-2.05,0) -- (2.45,0);
    \draw[qing!55, line width=0.7pt, ->] (0,-1.55) -- (0,1.85);
    \draw[zhe, line width=1.7pt] (0,0) ellipse (1.95 and 0.72);
  \end{scope}
  \begin{scope}[xshift=13.65cm]
    \draw[dai, line width=1.5pt, ->] (0,0) -- (1.866,0.571);
    \draw[zhe, line width=1.5pt, ->] (0,0) -- (-0.211,0.690);
    \node[font=\footnotesize\sffamily, color=inkl] at (0.2,-1.85) {④ $U$：再转};
  \end{scope}
\end{lafig}

\begin{keypoint}
\textbf{奇异值就是椭圆的半轴长。}
最大的 $\sigma_1$ 告诉我们：这个变换在\textbf{哪个方向}上拉伸得最厉害、
拉伸了多少倍。因此
\[
\norm{A}_2=\sigma_1,\qquad
\kappa(A)=\sigma_1/\sigma_n
\]
分别是矩阵的谱范数与条件数。$\sigma_n=0$ 当且仅当 $A$ 不可逆。
\end{keypoint}

\section{由 SVD 读出四个基本子空间}

把 $A=U\Sigma V^{\mathsf T}$ 按列展开：
\[
A=\sum_{i=1}^{r}\sigma_i\boldsymbol{u}_i\boldsymbol{v}_i^{\mathsf T}.
\]
于是\textbf{四个基本子空间的正交基全部现成}：

\begin{latab}{由 SVD 读出四个基本子空间的正交基}
\small\centering
\renewcommand{\arraystretch}{1.45}
\begin{tabular}{@{}lll@{}}
\toprule
\sffamily 子空间 & \sffamily 正交基 & \sffamily 来源\\
\midrule
列空间 $\Col(A)$ & $\boldsymbol{u}_1,\dots,\boldsymbol{u}_r$ & $U$ 的前 $r$ 列\\
左零空间 $\Nul(A^{\mathsf T})$ & $\boldsymbol{u}_{r+1},\dots,\boldsymbol{u}_{m}$ & $U$ 的后 $m-r$ 列\\
行空间 $\Row(A)$ & $\boldsymbol{v}_1,\dots,\boldsymbol{v}_r$ & $V$ 的前 $r$ 列\\
零空间 $\Nul(A)$ & $\boldsymbol{v}_{r+1},\dots,\boldsymbol{v}_n$ & $V$ 的后 $n-r$ 列\\
\bottomrule
\end{tabular}
\end{latab}

这正是第四章 Strang 图的“SVD 版本”：$V$ 与 $U$ 分别把
$\Rnum{n}$ 与 $\Rnum{m}$ 分解成两两正交的两块，
而 $A$ 把 $V$ 的前 $r$ 个方向一一映到 $U$ 的前 $r$ 个方向，并放大 $\sigma_i$ 倍。

\section{最优低秩逼近}

\begin{thm}{Eckart--Young 定理}{eckart}
在所有秩不超过 $k$ 的矩阵中，取 $A$ 的前 $k$ 项截断
\[
A_k=\sum_{i=1}^{k}\sigma_i\boldsymbol{u}_i\boldsymbol{v}_i^{\mathsf T}
\]
使误差的谱范数与 Frobenius 范数同时最小，且
\[
\norm{A-A_k}_2=\sigma_{k+1},\qquad
\norm{A-A_k}_F^2=\sigma_{k+1}^2+\sigma_{k+2}^2+\cdots+\sigma_r^2.
\]
\end{thm}

\begin{keypoint}
\textbf{用 $k$ 个奇异值就能“最好地”概括整个矩阵。}
误差由被丢弃的奇异值平方和刻画。
若奇异值衰减得快（多数真实数据如此），
那么保留前几项就能还原几乎全部信息——这就是\textbf{低秩压缩}。
\end{keypoint}

\begin{exmp}{一个手算的小例子}{svd-ex}
$A=\begin{bmatrix}4&0\\0&2\end{bmatrix}$ 已经是对角的，故
$\Sigma=\diag(4,2)$，$U=V=I$，$\sigma_1=4$，$\sigma_2=2$。
取 $k=1$ 的逼近 $A_1=\begin{bmatrix}4&0\\0&0\end{bmatrix}$，
误差 $\norm{A-A_1}_2=2=\sigma_2$。
\end{exmp}

\begin{remk}{$\triangleright$ 上机：奇异值衰减}{svd-code}
{\small\ttfamily
import numpy as np\\
U, s, Vt = np.linalg.svd(A, full\_matrices=False)\\
k = 10\\
Ak = (U[:, :k] * s[:k]) @ Vt[:k]}
\par\smallskip
\noindent 先画 \texttt{plt.semilogy(s)}：若奇异值在 $10$ 附近骤降，
说明数据本质上是 $10$ 维的。
\end{remk}

\begin{lafig}[>=Stealth, line cap=round]{奇异值衰减与逼近误差：左为典型数据的奇异值谱，
右为保留 $k$ 项后的相对误差 $\bigl(\sum_{i>k}\sigma_i^2\bigr)^{1/2}/\norm{A}_F$}
\begin{axis}[name=sx1, width=7.0cm, height=5.1cm,
  xlabel={序号 $i$}, ylabel={$\sigma_i$},
  label style={font=\small\sffamily, text=ink},
  tick label style={font=\footnotesize},
  axis lines=left, xtick={1,3,5,7,9,11,13,15},
  ymin=0, ymax=105, xmin=0.4, xmax=15.6,
  every axis plot/.append style={line width=1.1pt}]
  \addplot[qing, mark=*, mark size=1.25pt] coordinates {
    (1,100)(2,63)(3,38)(4,21)(5,12.5)(6,8.2)(7,5.6)(8,3.9)
    (9,2.7)(10,1.9)(11,1.3)(12,0.9)(13,0.6)(14,0.4)(15,0.25)};
  \draw[zhe!70, dashed, line width=0.7pt] (axis cs:5,0) -- (axis cs:5,100)
    node[pos=0.90, right=2pt, font=\scriptsize\sffamily, text=zhe] {取 $k=5$};
\end{axis}
\begin{axis}[name=sx2, width=7.0cm, height=5.1cm,
  at={(sx1.east)}, anchor=west, xshift=14pt,
  xlabel={保留个数 $k$}, ylabel={相对误差},
  label style={font=\small\sffamily, text=ink},
  tick label style={font=\footnotesize},
  axis lines=left, xtick={1,5,10,15}, ytick={0,25,50,75,100},
  ymin=0, ymax=105, xmin=0.4, xmax=15.6,
  every axis plot/.append style={line width=1.1pt}]
  \addplot[dai, thick, mark=square*, mark size=1.15pt] coordinates {
    (1,93.7)(2,77.5)(3,63.6)(4,50.4)(5,38.9)(6,29.7)(7,22.3)
    (8,16.2)(9,11.4)(10,7.8)(11,5.1)(12,3.2)(13,1.8)(14,0.9)(15,0)};
  \draw[qing!60, dashed, line width=0.7pt] (axis cs:0.4,10) -- (axis cs:15.6,10)
    node[pos=0.06, above=1pt, font=\scriptsize\sffamily, text=inkl] {$10\%$ 容差};
  \draw[ink!60, dotted, line width=0.7pt] (axis cs:5,0) -- (axis cs:5,38.9);
\end{axis}
\end{lafig}

\begin{remk}{如何读这两幅图}{read-decay}
左图是“谱”：奇异值从 $\sigma_1=100$ 一路衰减到 $0.25$，
说明该数据虽然名义上有 15 个坐标，能量却集中在最前面几个方向。
右图是“账单”：保留 $k=5$ 时相对误差约 $39\%$，保留 $k=10$ 时已不足 $8\%$；
若业务容差为 $10\%$，则只需存储前 9 个奇异分量。
\textbf{衰减越陡，压缩越狠}——这就是第 9.6 节图像压缩的全部原理。
\end{remk}

\section{主成分分析（PCA）}

设有一组中心化后的数据点
$\boldsymbol{x}_1,\dots,\boldsymbol{x}_m\in\Rnum{n}$
（即每个坐标都已减去均值），把它们按行排成矩阵 $X\in\R{m\times n}$。
样本协方差矩阵为
\[
S=\frac{1}{m-1}X^{\mathsf T}X,
\]
它是对称半正定矩阵。设 $S$ 的特征分解为
$S=Q\diag(\lambda_1\ge\lambda_2\ge\cdots\ge\lambda_n)Q^{\mathsf T}$。

\begin{thm}{主成分的方向}{pca-dir}
数据在单位向量 $\boldsymbol{w}$ 上的投影方差为
$\boldsymbol{w}^{\mathsf T}S\boldsymbol{w}$。
由瑞利商（第八章），方差最大的方向是\textbf{最大特征值对应的特征向量}，
最大方差就是 $\lambda_1$。
依次类推：第 $k$ 个主成分是第 $k$ 大特征向量，方差为 $\lambda_k$。
\end{thm}

\begin{lafig}[>=Stealth, line cap=round]{主成分分析：$\boldsymbol{v}_1$ 是方差最大的方向，
$\boldsymbol{v}_2$ 与之正交并承接剩余方差}
  \begin{axis}[
      width=0.80\textwidth, height=6.4cm,
      axis lines=left, axis line style={qing!55, line width=0.7pt},
      xmin=-3.6, xmax=4.2, ymin=-2.6, ymax=2.9,
      xtick={-3,-2,-1,0,1,2,3,4}, ytick={-2,-1,0,1,2},
      tick label style={font=\footnotesize, color=inkl},
      label style={font=\small, color=ink},
      xlabel={$x_1$}, ylabel={$x_2$},
      grid=major, grid style={line!70, very thin},
    ]
    \addplot[only marks, mark=*, mark size=1.35pt, color=qing!75,
             mark options={fill=qing!30}] coordinates {
      (-3.1,-1.05) (-2.5,-0.55) (-2.0,-0.95) (-1.5,-0.25) (-1.1,-0.75)
      (-0.6,-0.10) (-0.2,-0.45) (0.3,0.20) (0.7,-0.15) (1.2,0.45)
      (1.6,0.05) (2.0,0.75) (2.5,0.45) (2.9,1.05) (3.4,0.85) (3.7,1.45)};
    \addplot[dai, domain=-3.4:4.0, samples=2, line width=1.7pt] {0.335*x+0.10};
    \addplot[zhe, domain=-1.55:1.95, samples=2, line width=1.5pt]
      {-2.97*x+0.95};
    \node[dai, font=\footnotesize] at (axis cs:3.35,1.62) {$\boldsymbol{v}_1$};
    \node[zhe, font=\footnotesize] at (axis cs:1.62,2.35) {$\boldsymbol{v}_2$};
    \node[font=\footnotesize, color=inkl, anchor=south west] at (axis cs:-3.5,-2.45)
      {$\sigma_1$ 方向：数据最“长”的方向；$\sigma_2$ 方向：剩余方差的方向};
  \end{axis}
\end{lafig}

\begin{prop}{PCA 就是 SVD}{pca-svd}
对中心化数据矩阵 $X$ 作 SVD：$X=U\Sigma V^{\mathsf T}$。
则
\[
V\ \text{的各列}=\text{主成分方向，}\qquad
\sigma_i^2/(m-1)=\text{第 } i \text{ 主成分的方差}.
\]
\textbf{因此 PCA 无须显式构造 $X^{\mathsf T}X$，直接对 $X$ 做 SVD 即可}——
这既更快也更稳定。
\end{prop}

\begin{remk}{PCA 在做什么}{pca-meaning}
PCA 寻找一个\textbf{新的正交坐标系}，使数据在这个坐标系下各方向互不相关，
并且方差按从大到小排列。保留前 $k$ 个坐标，
就是\textbf{在所有 $k$ 维线性子空间中，投影误差最小的一个}——
这正是一次 Eckart--Young 逼近。
图像压缩、人脸识别（特征脸）、降噪、可视化，用的都是这同一个动作。
\end{remk}

\section{应用：图像压缩}

一张灰度图像就是一个非负矩阵 $A\in\R{m\times n}$。
取前 $k$ 个奇异值，只需存储
$k(m+n+1)$ 个数（$U_k$、$\Sigma_k$、$V_k^{\mathsf T}$），
而原图要 $mn$ 个。压缩比约为
\[
\rho=\frac{k(m+n+1)}{mn}.
\]

\begin{latab}{一幅 $512\times 512$ 图像在不同秩下的存储量与误差}
\small\centering
\renewcommand{\arraystretch}{1.4}
\begin{tabular}{@{}ccccc@{}}
\toprule
\sffamily 秩 $k$ & \sffamily 存储数 & \sffamily 压缩比 $\rho$ &
\sffamily 相对误差 $\norm{A-A_k}_F/\norm{A}_F$ & \sffamily 视觉效果\\
\midrule
5    & 5\,165   & 2.0\%  & 0.31 & 轮廓可辨，细节丢失\\
20   & 20\,660  & 7.9\%  & 0.17 & 主要结构清晰\\
50   & 51\,650  & 19.7\% & 0.10 & 接近原图\\
100  & 103\,300 & 39.4\% & 0.06 & 肉眼难辨差异\\
512  & 262\,144 & 100\%  & 0    & 原图\\
\bottomrule
\end{tabular}
\end{latab}

（表中误差为示意值，实际取决于图像的奇异值衰减速度。
关键观察是：\textbf{误差由被丢弃的奇异值决定，而非由像素数量决定}。）

\begin{sidebar}{贝尔特拉米与若尔当：SVD 的两条发现路径}
奇异值分解通常归于贝尔特拉米（E.\ Beltrami，1873）与若尔当（C.\ Jordan，1874），
二人几乎同时独立给出双线性形式的分解。
但真正让 SVD 成为计算利器的，是 1936 年 Eckart--Young 定理
给出了“截断即最优”的保证，
以及 1965 年 Golub 提出的稳定算法（基于对双边 Jacobi 或 Householder 变换）。
今天，从 Google 的 PageRank 到 Netflix 的推荐系统，
从人脸识别到引力波信号的去噪，背后都是同一条分解式
$A=U\Sigma V^{\mathsf T}$。
\end{sidebar}

\begin{chapkeys}
\small
\begin{multicols}{2}
\begin{itemize}[leftmargin=1.5em, itemsep=2pt]
\item 任意矩阵都有 SVD：$A=U\Sigma V^{\mathsf T}$。
\item $\sigma_i=\sqrt{\lambda_i(A^{\mathsf T}A)}$，即椭圆半轴长。
\item $A\boldsymbol{v}_i=\sigma_i\boldsymbol{u}_i$。
\item $U,V$ 的列给出四个基本子空间的正交基。
\item $A=\sum_i\sigma_i\boldsymbol{u}_i\boldsymbol{v}_i^{\mathsf T}$。
\item Eckart--Young：截断 $A_k$ 是最优低秩逼近，误差 $=\sigma_{k+1}$。
\item PCA：主成分方向 $=$ 协方差矩阵的特征向量 $=$ $V$ 的列。
\item 条件数 $\kappa(A)=\sigma_1/\sigma_n$。
\end{itemize}
\end{multicols}
\end{chapkeys}

\exhead
\begin{multicols}{2}
\small
\textbf{A 组（基础计算）}
\begin{enumerate}[leftmargin=1.6em, itemsep=3pt]
\item 求 $A=\begin{bmatrix}3&0\\0&2\end{bmatrix}$ 与
      $B=\begin{bmatrix}2&0\\0&3\end{bmatrix}$ 的奇异值，并说明
      它们的特征值与奇异值有何不同。
\item 求 $A=\begin{bmatrix}1&1\\0&0\end{bmatrix}$ 的 SVD。
\item 设 $\sigma_1\ge\sigma_2\ge\sigma_3$ 为某矩阵的奇异值，
      求秩 $1$ 最佳逼近的谱范数误差与 Frobenius 误差。
\item 已知 $A$ 为 $4\times 3$ 矩阵，$\rank A=2$。写出 $\Sigma$ 的形状，
      并指出 $U$、$V$ 的哪些列张成 $\Nul(A^{\mathsf T})$。
\end{enumerate}
\textbf{B 组（推理证明）}
\begin{enumerate}[leftmargin=1.6em, itemsep=3pt, start=5]
\item 证明：$\sigma_i^2$ 是 $A^{\mathsf T}A$ 的特征值，
      且 $V$ 的列是相应的特征向量。
\item 证明：$\norm{A}_F^2=\sigma_1^2+\cdots+\sigma_r^2$。
\item 证明：$\rank(A)=r$（非零奇异值的个数）。
\item 设 $A$ 对称正定。证明 $A$ 的奇异值就是它的特征值。
\end{enumerate}
\textbf{C 组（应用与建模）}
\begin{enumerate}[leftmargin=1.6em, itemsep=3pt, start=9]
\item 用 NumPy 对一幅灰度图做 SVD，分别取 $k=5,20,50$ 重构并保存图像，
      计算压缩比与相对误差，画出奇异值的半对数图。
\item 对一组二维人工数据（沿某方向拉长的椭圆分布）做 PCA，
      画出主成分方向，并验证第一主成分的方差等于最大特征值。
\end{enumerate}
\end{multicols}


\appendix
% ===================== 附录 =====================
\chapter{常用公式速查}
\label{app:formula}

\begin{center}
{\small\sffamily\color{inkl}本附录将全书主要公式按章汇总，供考前与做题时快速查阅。}
\end{center}

\begin{multicols}{2}
\small\sloppy
\setstretch{1.2}

\minihead{向量与矩阵运算}
\begin{itemize}[leftmargin=1.4em, itemsep=2pt]
\item $(AB)_{ij}=\sum_k a_{ik}b_{kj}$，要求 $(m\times n)(n\times p)$
\item $(AB)^{\mathsf T}=B^{\mathsf T}A^{\mathsf T}$，
      $(AB)^{-1}=B^{-1}A^{-1}$
\item $A^{\mathsf T}=A$ 对称；$A^{\mathsf T}=-A$ 反对称
\item $A^k=PD^kP^{-1}$（可对角化时）
\end{itemize}

\minihead{方程组与秩}
\begin{itemize}[leftmargin=1.4em, itemsep=2pt]
\item 有解 $\Leftrightarrow \rank(A\,|\,\boldsymbol{b})=\rank(A)$
\item 解集 $=\boldsymbol{x}_p+\Nul(A)$
\item 主元个数 $=\rank$；自由变量 $=n-\rank$
\item $\rank(A)+\dim\Nul(A)=n$
\item $\rank(A)=\rank(A^{\mathsf T})\le\min(m,n)$
\end{itemize}

\minihead{行列式}
\begin{itemize}[leftmargin=1.4em, itemsep=2pt]
\item $\det A=\sum_{\sigma}\sign(\sigma)\prod_i a_{i\sigma(i)}$
\item 上三角：$\det A=\prod_i a_{ii}$
\item $\det(AB)=\det A\,\det B$，$\det(A^{\mathsf T})=\det A$
\item $\det(A^{-1})=1/\det A$，$\det(cA)=c^n\det A$
\item 按行展开：$\det A=\sum_k a_{ik}A_{ik}$
\item 异乘变零：$\sum_k a_{ik}A_{jk}=0\ (i\neq j)$
\item 克拉默：$x_j=\det A_j/\det A$
\end{itemize}

\minihead{四个基本子空间}
\begin{itemize}[leftmargin=1.4em, itemsep=2pt]
\item $\Col(A)\subset\Rnum{m}$，维数 $r$
\item $\Nul(A)\subset\Rnum{n}$，维数 $n-r$
\item $\Row(A)=\Col(A^{\mathsf T})\subset\Rnum{n}$，维数 $r$
\item $\Nul(A^{\mathsf T})\subset\Rnum{m}$，维数 $m-r$
\item $\Row(A)\perp\Nul(A)$，$\Col(A)\perp\Nul(A^{\mathsf T})$
\end{itemize}

\minihead{线性变换}
\begin{itemize}[leftmargin=1.4em, itemsep=2pt]
\item 标准矩阵 $A=[\,T\boldsymbol{e}_1\ \cdots\ T\boldsymbol{e}_n\,]$
\item $\dim\ker T+\dim\operatorname{im}T=\dim V$
\item $[T\circ S]=[T][S]$
\item 相似 $B=P^{-1}AP$：保 $\det$、$\tr$、秩、特征值
\end{itemize}

\minihead{特征值}
\begin{itemize}[leftmargin=1.4em, itemsep=2pt]
\item $p_A(\lambda)=\det(\lambda I-A)$
\item $\sum\lambda_i=\tr A$，$\prod\lambda_i=\det A$
\item 几何重数 $\le$ 代数重数
\item 可对角化 $\Leftrightarrow$ $n$ 个无关特征向量
\item 对称矩阵：$Q^{\mathsf T}AQ=D$，$Q^{\mathsf T}Q=I$
\item 瑞利商：$\lambda_{\min}\le\dfrac{\boldsymbol{x}^{\mathsf T}A\boldsymbol{x}}
      {\boldsymbol{x}^{\mathsf T}\boldsymbol{x}}\le\lambda_{\max}$
\end{itemize}

\minihead{内积与正交}
\begin{itemize}[leftmargin=1.4em, itemsep=2pt]
\item $\norm{\boldsymbol{v}}=\sqrt{\boldsymbol{v}^{\mathsf T}\boldsymbol{v}}$
\item 柯西—施瓦茨：$\abs{\boldsymbol{u}\cdot\boldsymbol{v}}
      \le\norm{\boldsymbol{u}}\norm{\boldsymbol{v}}$
\item 投影：$\widehat{\boldsymbol{b}}=A(A^{\mathsf T}A)^{-1}A^{\mathsf T}\boldsymbol{b}$
\item 法方程：$A^{\mathsf T}A\widehat{\boldsymbol{x}}=A^{\mathsf T}\boldsymbol{b}$
\item $A=QR$：$R\widehat{\boldsymbol{x}}=Q^{\mathsf T}\boldsymbol{b}$
\end{itemize}

\minihead{二次型与正定}
\begin{itemize}[leftmargin=1.4em, itemsep=2pt]
\item 正定 $\Leftrightarrow$ $\lambda_i>0$ $\Leftrightarrow$ $\Delta_k>0$
      $\Leftrightarrow$ $A=C^{\mathsf T}C$
\item 负定：$(-1)^k\Delta_k>0$ 对一切 $k$
\item 惯性定理：符号差 $(p,q)$ 不变
\item 椭圆半轴长 $1/\sqrt{\lambda_i}$
\end{itemize}

\minihead{奇异值分解}
\begin{itemize}[leftmargin=1.4em, itemsep=2pt]
\item $A=U\Sigma V^{\mathsf T}$，$\sigma_i=\sqrt{\lambda_i(A^{\mathsf T}A)}$
\item $A\boldsymbol{v}_i=\sigma_i\boldsymbol{u}_i$
\item $A=\sum_{i=1}^{r}\sigma_i\boldsymbol{u}_i\boldsymbol{v}_i^{\mathsf T}$
\item $\norm{A}_2=\sigma_1$，$\norm{A}_F^2=\sum\sigma_i^2$
\item $\norm{A-A_k}_2=\sigma_{k+1}$
\item $\kappa(A)=\sigma_1/\sigma_n$
\end{itemize}

\end{multicols}

% ===================== 附录 B =====================
\chapter{Python 与 MATLAB 命令对照}
\label{app:code}

\begin{center}
{\small\sffamily\color{inkl}
约定：\texttt{A}, \texttt{b} 为已定义的矩阵与向量；下表为最常用的核心操作。}
\end{center}

\begin{multicols}{2}
\small\sloppy
\setstretch{1.22}

\minihead{基本操作}
\begin{itemize}[leftmargin=1.4em, itemsep=2pt]
\item 建矩阵：
      \texttt{np.array([[1,2],[3,4]])} \ / \ \texttt{[1 2; 3 4]}
\item 转置：\texttt{A.T} \ / \ \texttt{A'}
\item 乘法：\texttt{A @ B} \ / \ \texttt{A*B}
\item 逐元素：\texttt{A * B} \ / \ \texttt{A.*B}
\item 单位阵：\texttt{np.eye(n)} \ / \ \texttt{eye(n)}
\item 零矩阵：\texttt{np.zeros((m,n))} \ / \ \texttt{zeros(m,n)}
\end{itemize}

\minihead{解方程组}
\begin{itemize}[leftmargin=1.4em, itemsep=2pt]
\item \texttt{np.linalg.solve(A, b)} \ / \ \texttt{A\textbackslash b}
\item LU：
      \texttt{scipy.linalg.lu\_factor}，\texttt{lu\_solve}
      \ / \ \texttt{[L,U,P] = lu(A)}
\item 逆：\texttt{np.linalg.inv(A)} \ / \ \texttt{inv(A)}
\item 伪逆：\texttt{np.linalg.pinv(A)} \ / \ \texttt{pinv(A)}
\end{itemize}

\minihead{秩与阶梯形}
\begin{itemize}[leftmargin=1.4em, itemsep=2pt]
\item 秩：\texttt{np.linalg.matrix\_rank(A)} \ / \ \texttt{rank(A)}
\item RREF（Python 需 \texttt{sympy}）：
      \texttt{sympy.Matrix(A).rref()}
\item 行列式：\texttt{np.linalg.det(A)} \ / \ \texttt{det(A)}
\end{itemize}

\minihead{特征值}
\begin{itemize}[leftmargin=1.4em, itemsep=2pt]
\item 一般矩阵：
      \texttt{np.linalg.eig(A)} \ / \ \texttt{eig(A)}
\item 对称矩阵：
      \texttt{np.linalg.eigh(A)} \ / \ \texttt{eig(A)}
\item 特征多项式：
      \texttt{np.poly(A)} \ / \ \texttt{poly(A)}
\end{itemize}

\minihead{正交化与最小二乘}
\begin{itemize}[leftmargin=1.4em, itemsep=2pt]
\item QR：\texttt{np.linalg.qr(A)} \ / \ \texttt{[Q,R] = qr(A)}
\item 最小二乘：
      \texttt{np.linalg.lstsq(A, b)} \ / \ \texttt{A\textbackslash b}
      （NumPy 须加 \texttt{rcond=None}）
\item 投影：\texttt{A @ np.linalg.solve(A.T@A, A.T @ b)}
\end{itemize}

\minihead{奇异值分解}
\begin{itemize}[leftmargin=1.4em, itemsep=2pt]
\item \texttt{np.linalg.svd(A)} \ / \ \texttt{[U,S,V]=svd(A)}
\item 截断：\texttt{(U[:,:k]*s[:k]) @ Vt[:k]}
\item 条件数：\texttt{np.linalg.cond(A)} \ / \ \texttt{cond(A)}
\item 范数：\texttt{np.linalg.norm(A)} \ / \ \texttt{norm(A)}
\end{itemize}

\minihead{作图（Python）}
\begin{itemize}[leftmargin=1.4em, itemsep=2pt]
\item 散点：\texttt{plt.scatter(x, y)}
\item 拟合线：\texttt{plt.plot(x, m*x+c)}
\item 奇异值衰减：\texttt{plt.semilogy(s)}
\item 等高线：\texttt{plt.contour(X, Y, F, levels=12)}
\end{itemize}

\end{multicols}

\begin{keypoint}
\textbf{三条数值经验}：
(1) 解方程组用 \texttt{solve}，不要用 \texttt{inv(A)@b}；
(2) 对称矩阵用 \texttt{eigh}，不要用 \texttt{eig}；
(3) 判断相等一律用 \texttt{np.allclose}，不要用 \texttt{==}。
浮点运算的舍入误差，是理论数学与计算数学之间最常见的一条鸿沟。
\end{keypoint}

% ===================== 附录 C =====================
\chapter{中英术语对照}
\label{app:terms}

\begin{multicols}{2}
\small
\setstretch{1.18}
\begin{itemize}[leftmargin=1.2em, itemsep=1.5pt, label={}]
\item 向量空间 \hfill vector space
\item 子空间 \hfill subspace
\item 线性组合 \hfill linear combination
\item 张成 \hfill span
\item 线性相关 / 无关 \hfill linearly dependent / independent
\item 基 \hfill basis
\item 维数 \hfill dimension
\item 坐标 \hfill coordinates
\item 行阶梯形 \hfill echelon form
\item 简化行阶梯形 \hfill reduced echelon form (RREF)
\item 主元 \hfill pivot
\item 自由变量 \hfill free variable
\item 秩 \hfill rank
\item 增广矩阵 \hfill augmented matrix
\item 初等行变换 \hfill elementary row operations
\item 可逆矩阵 \hfill invertible matrix
\item 行列式 \hfill determinant
\item 余子式 / 代数余子式 \hfill minor / cofactor
\item 伴随矩阵 \hfill adjugate matrix
\item 克拉默法则 \hfill Cramer's rule
\item 列空间 \hfill column space
\item 零空间 \hfill null space
\item 行空间 \hfill row space
\item 左零空间 \hfill left null space
\item 秩—零化度定理 \hfill rank--nullity theorem
\item 线性变换 \hfill linear transformation
\item 核 \hfill kernel
\item 像 \hfill image / range
\item 复合 \hfill composition
\item 相似矩阵 \hfill similar matrices
\item 特征值 / 特征向量 \hfill eigenvalue / eigenvector
\item 特征多项式 \hfill characteristic polynomial
\item 代数重数 / 几何重数 \hfill algebraic / geometric multiplicity
\item 对角化 \hfill diagonalization
\item 谱定理 \hfill spectral theorem
\item 正交矩阵 \hfill orthogonal matrix
\item 内积 \hfill inner product
\item 范数 \hfill norm
\item 标准正交基 \hfill orthonormal basis
\item 格拉姆—施密特 \hfill Gram--Schmidt
\item 正交投影 \hfill orthogonal projection
\item 最小二乘 \hfill least squares
\item 法方程 \hfill normal equations
\item 二次型 \hfill quadratic form
\item 惯性定理 \hfill Sylvester's law of inertia
\item 合同 \hfill congruent
\item 正定 \hfill positive definite
\item 顺序主子式 \hfill leading principal minor
\item 瑞利商 \hfill Rayleigh quotient
\item 奇异值分解 \hfill singular value decomposition
\item 奇异值 / 奇异向量 \hfill singular value / vector
\item 条件数 \hfill condition number
\item 低秩逼近 \hfill low-rank approximation
\item 主成分分析 \hfill principal component analysis
\end{itemize}
\end{multicols}

% ===================== 参考文献 =====================

% ---- YangTex 补注（唯一新增文字）：原稿列了参考文献但正文未引用，
%      补一句以便 PDF 上的引用编号可以悬停反查；原稿文字一字未删改。 ----
\par\noindent 延伸阅读可参见~\cite{tongji,lay,strang}。\par
\begin{thebibliography}{9}
\small
\setstretch{1.15}

\bibitem{tongji}
同济大学数学系. 线性代数（第六版）. 高等教育出版社, 2014.

\bibitem{lay}
David C.\ Lay, Steven R.\ Lay, Judi J.\ McDonald.
\emph{Linear Algebra and Its Applications} (5th ed.). Pearson, 2015.
（中译本：线性代数及其应用，机械工业出版社.）

\bibitem{strang}
Gilbert Strang. \emph{Introduction to Linear Algebra} (6th ed.).
Wellesley-Cambridge Press, 2023.

\bibitem{axler}
Sheldon Axler. \emph{Linear Algebra Done Right} (4th ed.). Springer, 2024.

\bibitem{horn}
Roger A.\ Horn, Charles R.\ Johnson.
\emph{Matrix Analysis} (2nd ed.). Cambridge University Press, 2013.

\bibitem{trefethen}
Lloyd N.\ Trefethen, David Bau III.
\emph{Numerical Linear Algebra}. SIAM, 1997.

\bibitem{goodfellow}
Ian Goodfellow, Yoshua Bengio, Aaron Courville.
\emph{Deep Learning}. MIT Press, 2016.（第 2 章为线性代数速览.）

\end{thebibliography}

\vspace{0.6cm}
\begin{flushright}
\cjkkai
线性代数的全部内容，可以压缩成一句话：\\
\emph{看清一个变换在哪些方向上只做伸缩。}\\[4pt]
{\sffamily\footnotesize\color{inkl}—— 全书完 ——}
\end{flushright}


\end{document}
